\chapter{模型压缩与加速}
\label{ch:11}

\paragraph*{本章阅读锚点}
本章的锚点是效率。原始 3DGS 以大量高斯换实时渲染；压缩与加速方法要减少高斯数、参数量、带宽、排序或训练开销，同时守住画质。


3D~Gaussian Splatting~以显式、离散的高斯基元表示场景，带来了实时渲染的巨大优势，但也付出了沉重的存储代价。一个典型的~Mip-NeRF 360~场景训练后包含数百万个高斯，每个高斯携带~59~个浮点参数（位置~3、旋转~4、缩放~3、不透明度~1、球谐系数~48），总计超过~700~MB。这一体量远超手机、VR（虚拟现实，Virtual Reality）~头显等边缘设备的承载能力，也使网络传输变得不切实际。模型压缩与加速因此成为~3DGS~走向实用化的关键技术瓶颈。

本章围绕"如何将百万级高斯的庞大模型压缩到可部署规模"这一核心问题，系统梳理了~161~篇相关工作，按~6~个主题展开，形成两条紧密关联的技术主线：

\textbf{减少基元数量}。第~\ref{sec:11_1}~节讨论高斯剪枝与稀疏化——从训练后模型中识别并移除冗余高斯，保留~10\%--50\% 即可维持可接受的渲染质量。第~\ref{sec:11_5}~节讨论结构化与锚点表示——用锚点驱动的层级结构替代逐高斯独立表示，从根本上减少冗余基元的产生。

\textbf{压缩基元参数}。第~\ref{sec:11_2}~节讨论编码与压缩——用量化、向量量化（VQ）、熵编码和上下文建模等技术压缩每个高斯的~59~维参数向量。第~\ref{sec:11_3}~节讨论训练加速与高效表示——在训练阶段即控制高斯增长，从源头减少需要压缩的基元数量。第~\ref{sec:11_4}~节讨论视频与流式渲染——将压缩技术与流式传输结合，实现渐进式场景加载。第~\ref{sec:11_6}~节覆盖其他压缩方法与硬件加速。

\begin{figure}[htbp]
\centering
\resizebox{\textwidth}{!}{\begin{tikzpicture}[
  route/.style={rectangle, rounded corners=5pt, draw=#1, fill=#1!10,
    text width=3cm, minimum height=0.9cm, align=center, font=\small},
  method/.style={rectangle, rounded corners=3pt, draw=#1, fill=#1!6,
    font=\scriptsize, text width=2.8cm, align=center, minimum height=0.6cm},
  arr/.style={-{Stealth[length=4pt]}, techgray!60},
  >=Stealth
]

  \node[route=techgray, text width=4cm, font=\small\bfseries] (root) at (31.2, 5.6) {模型压缩与加速方法分类};
  \node[route=techblue] (b0) at (7.3, 3.1) {高斯剪枝};
  \node[method=techblue] (b0_l0) at (1.5, 0.6) {评分剪枝\\{\tiny LightGaussian, PUP 3D-GS, GaussianPOP}};
  \node[method=techblue] (b0_l1) at (5.1, 0.6) {掩码剪枝\\{\tiny MaskGaussian, SVR-GS, LP-3DGS}};
  \node[method=techblue] (b0_l2) at (9.5, 0.6) {正则化/最优传输\\{\tiny SafeguardGS, Clean-GS, Gaussian Herding}};
  \node[route=techgreen] (b1) at (19.2, 3.1) {编码压缩};
  \node[method=techgreen] (b1_l0) at (14.0, 0.6) {向量量化\\{\tiny Compact 3DGS, ContraGS, SALVQ}};
  \node[method=techgreen] (b1_l1) at (17.4, 0.6) {熵编码\\{\tiny HAC++, ContextGS, CAT-3DGS, HEMGS}};
  \node[method=techgreen] (b1_l2) at (21.0, 0.6) {前馈/后训练\\{\tiny FCGS, MesonGS++, POTR}};
  \node[route=techorange] (b2) at (29.2, 3.1) {训练加速};
  \node[method=techorange] (b2_l0) at (25.0, 0.6) {频率调制\\{\tiny Opti3DGS, AutoOpti3DGS}};
  \node[method=techorange] (b2_l1) at (29.2, 0.6) {内存/流程优化\\{\tiny DashGaussian, SkipGS, Taming 3DGS, Gaussians on a Diet}};
  \node[route=techpurple] (b3) at (33.4, 3.1) {视频流式};
  \node[method=techpurple] (b3_l0) at (33.4, 0.6) {渐进传输\\{\tiny PRoGS, GoDe, Matryoshka GS}};
  \node[method=techpurple] (b3_l1) at (37.2, 0.6) {时序一致性\\{\tiny L3GS, LapisGS}};
  \node[route=techblue] (b4) at (41.1, 3.1) {锚点表示};
  \node[method=techblue] (b4_l0) at (41.1, 0.6) {锚点驱动\\{\tiny Scaffold-GS, SOGS}};
  \node[method=techblue] (b4_l1) at (44.8, 0.6) {层级/混合\\{\tiny Octree-GS, HyRF, NeuralGS}};
  \node[route=techgray] (b5) at (48.4, 3.1) {硬件加速};
  \node[method=techgray] (b5_l0) at (48.4, 0.6) {边缘部署\\{\tiny ECO-GS}};
  \node[method=techgray] (b5_l1) at (52.0, 0.6) {极端压缩\\{\tiny NiFi, ExGS, SymGS}};
  \draw[arr] (root) -- (b0);
  \draw[arr] (root) -- (b1);
  \draw[arr] (root) -- (b2);
  \draw[arr] (root) -- (b3);
  \draw[arr] (root) -- (b4);
  \draw[arr] (root) -- (b5);
  \draw[arr] (b0) -- (b0_l0);
  \draw[arr] (b0) -- (b0_l1);
  \draw[arr] (b0) -- (b0_l2);
  \draw[arr] (b1) -- (b1_l0);
  \draw[arr] (b1) -- (b1_l1);
  \draw[arr] (b1) -- (b1_l2);
  \draw[arr] (b2) -- (b2_l0);
  \draw[arr] (b2) -- (b2_l1);
  \draw[arr] (b3) -- (b3_l0);
  \draw[arr] (b3) -- (b3_l1);
  \draw[arr] (b4) -- (b4_l0);
  \draw[arr] (b4) -- (b4_l1);
  \draw[arr] (b5) -- (b5_l0);
  \draw[arr] (b5) -- (b5_l1);
\end{tikzpicture}}
\caption{模型压缩与加速方法分类}
\label{fig:ch11_compression_taxonomy}
\end{figure}

\begin{figure}[htbp]
\centering
\resizebox{\textwidth}{!}{\begin{tikzpicture}[
  box/.style={rectangle, rounded corners=4pt, draw=#1, fill=#1!10,
    text width=3.2cm, align=center, font=\small, minimum height=1cm},
  sbox/.style={rectangle, rounded corners=3pt, draw=#1!60, fill=#1!5,
    font=\scriptsize, text width=2.4cm, align=center, minimum height=0.5cm},
  arr/.style={-{Stealth[length=6pt]}, thick, techgray},
  darr/.style={-{Stealth[length=6pt]}, thick, techgray, dashed},
  >=Stealth
]

  \node[box=techblue] (s1) at (3.7, 0.6) {高斯剪枝(评分/掩码/正则化)};
  \node[box=techgreen] (s2) at (12.3, 0.6) {编码压缩(向量量化/熵编码)};
  \node[box=techorange] (s3) at (20.7, 0.6) {渲染加速(LOD/流式/Web)};
  \node[font=\scriptsize, text=techgray] at (3.7, 1.8) {高斯数$\downarrow$90\%};
  \node[font=\scriptsize, text=techgray] at (12.3, 1.8) {存储$\downarrow$50$\times$};
  \node[font=\scriptsize, text=techgray] at (20.7, 1.8) {渲染$\rightarrow$实时};
  \draw[arr] (s1) -- (s2);
  \draw[arr] (s2) -- (s3);
\end{tikzpicture}}
\caption{模型压缩三阶段 pipeline：剪枝→编码→渲染加速}
\label{fig:ch11_compression_pipeline}
\end{figure}

%=====================================================================
\section{高斯剪枝与稀疏化}
\label{sec:11_1}

本节讨论从训练后或训练中的~3DGS~模型中识别并移除冗余高斯基元的方法，包括四条技术路线：基于重要性分数的剪枝、基于可学习掩码的剪枝、基于几何一致性的剪枝，以及基于最优传输的全局简化。

\subsection{问题背景}

以下逐一展开四条技术路线。

3DGS~的训练过程通过自适应密度控制（ADC）不断添加新高斯来拟合场景细节，但缺乏对高斯总量的有效约束。训练完成后的高斯集合中，大量基元对渲染质量的贡献极低。它们要么被其他高斯完全遮挡（冗余），要么覆盖的区域极小（微不足道）。剪枝（pruning）的目标是识别并移除这些低贡献高斯，在不显著降低渲染质量的前提下将模型压缩~2--10~倍。

剪枝面临的核心挑战是\textbf{重要性度量}的设计：如何准确评估每个高斯对最终渲染图像的贡献？直觉地说，一个"重要"的高斯应当满足两个条件：（1）它对训练视角的渲染像素有显著的颜色贡献；（2）移除它后渲染质量会明显下降。然而，这两个条件之间存在微妙的张力。某些高斯在当前训练视角下贡献不大，但对未见视角可能至关重要。重要性度量的设计直接影响剪枝后的质量保留程度。

\subsubsection{基于重要性分数的剪枝}

本节介绍通过为每个高斯计算全局重要性分数来指导剪枝的方法，从一阶累积贡献（LightGaussian）、二阶灵敏度分析（PUP 3D-GS）到解析渲染误差（GaussianPOP）逐步演进。

\paragraph{全局视角评估高斯显著性}

本方法修改了~3DGS~的\textbf{高斯管理环节}：用\textbf{全局显著性分数引导的渐进剪枝与蒸馏}替代\textbf{ADC~的"只增不减"密度控制}，\textbf{渲染管线（$\alpha$-混合与可微光栅化）保持不变}。

\textbf{问题与动机。}3DGS~的~ADC（自适应密度控制）策略在训练过程中不断克隆和分裂高斯以增强场景覆盖，但缺乏对高斯总量的有效约束。它只有"增加"机制（clone/split），没有"删除"机制（prune），导致训练后高斯数量从数千膨胀到数百万。一个直接的想法是"按不透明度剪枝"：移除不透明度最低的高斯。然而，不透明度只是高斯的局部属性，一个不透明度较低但覆盖面积很大的高斯可能对大量像素有不可忽视的颜色贡献。LightGaussian（2023.11）~\cite{2311_17245}~受深度学习中网络剪枝（network pruning，通过移除不重要的神经元或滤波器来压缩网络的技术）的启发，提出需要一个\textbf{全局视角}来评估每个高斯的真实重要性。这要求不仅看它"自身有多透明"，还要看它"在所有训练视角的渲染中实际贡献了多少颜色"。其核心思想是：一个高斯的重要性不取决于它在单张图像中的表现，而取决于它对所有训练视角渲染结果的累积影响。

\textbf{全局显著性分数。}具体而言，LightGaussian~的定义从~3DGS~的~$\alpha$-混合渲染方程出发。直觉上，3DGS~的渲染过程是一个逐层叠加的~$\alpha$-混合：

\[
C(\mathbf{p}) = \sum_{i=1}^{N} \alpha_i T_i \mathbf{c}_i
\]

其中 $\alpha_i$ 为第~$i$~个高斯在像素 $\mathbf{p}$ 处的不透明度，$T_i = \prod_{j=1}^{i-1}(1-\alpha_j)$ 为透射率（即该高斯前方所有高斯"漏过来"的光线比例），$\mathbf{c}_i$ 为颜色。LightGaussian~的关键洞察是：第~$i$~个高斯对像素 $\mathbf{p}$ 的有效颜色贡献为 $\alpha_i T_i \mathbf{c}_i$，贡献权重 $\alpha_i T_i$ 同时取决于高斯自身的不透明度和它前方所有高斯的遮挡。因此，LightGaussian~在所有训练视角的所有像素上累积每个高斯的贡献权重，得到全局显著性分数：

\begin{equation}
I(g_i) = \sum_{v \in \mathcal{V}} \sum_{\mathbf{p} \in \mathcal{P}_v} \mathbb{1}(g_i, r_{\mathbf{p}}) \cdot \alpha_i(\mathbf{p}) \cdot T_i(\mathbf{p}) \cdot V_{\mathrm{norm}}(g_i)
\end{equation}

其中 $\mathcal{V}$ 为训练视角集合，$\mathcal{P}_v$ 为视角~$v$~中的像素集合，$\mathbb{1}(g_i, r_{\mathbf{p}})$ 是指示函数（当高斯~$g_i$~与射线~$r_{\mathbf{p}}$~相交时为~1，否则为~0），$V_{\mathrm{norm}}(g_i)$~为归一化体积项。用白话解释这个公式：对每个高斯，数一数"它被多少条渲染射线命中"、每次命中时"它的不透明度有多大贡献"、以及"它的物理体积有多大"。该分数越高，说明高斯~$g_i$~对渲染结果的贡献越大，越不应被剪枝。

\textbf{体积归一化的设计细节。}公式中的体积项~$V(g_i) = \frac{4}{3}\pi a_i b_i c_i$（其中~$a_i, b_i, c_i$~为高斯的三个缩放维度）直接引入会过度偏向大面积的背景高斯，导致细小但关键的几何细节被误判为"不重要"。LightGaussian~因此引入了自适应归一化：先将每个高斯的体积除以全局第~90~百分位体积（$V_{\max 90}$，即所有高斯体积从大到小排序后第~90~百分位的值）进行归一化并截断到~$[0, 1]$，再通过可调指数~$\gamma$~控制体积项的影响力，即~$V_{\mathrm{norm}}(g_i) = (V_{\mathrm{raw}})^\gamma$，其中~$V_{\mathrm{raw}} = \min(V(g_i)/V_{\max 90},\, 1)$。这一归一化策略确保体积大的背景高斯不会仅凭尺寸就主导显著性排名。

\textbf{渐进剪枝三步流程。}基于全局显著性分数，LightGaussian~采用渐进式三步流程：（1）\textbf{按分数排序并剪枝}：将所有高斯按全局显著性分数排序，移除分数最低的一批高斯；（2）\textbf{恢复微调}：对剩余高斯进行~5000~次迭代的微调（仅使用光度损失，不执行额外的~clone/split~操作），让幸存高斯自适应地补偿被移除高斯的信息损失；（3）\textbf{球谐蒸馏}：通过知识蒸馏将高阶球谐系数压缩到低阶表示（下详）。

\textbf{球谐知识蒸馏。}在未压缩的~3DGS~表示中，球谐（SH，Spherical Harmonics，一种用正交基函数表示球面上函数的数学工具）系数占据了~81.3\% 的特征维度——每个高斯需要~$(45+3)$~个浮点值来存储~3~阶球谐。直接降低~SH~阶数可以节省存储，但会丢失镜面反射和视角相关的外观细节。LightGaussian~设计了教师-学生蒸馏方案：教师模型保留完整的~SH~系数，学生模型使用截断的低阶~SH。蒸馏损失基于两者在渲染图像上的像素级差异：

\begin{equation}
\mathcal{L}_{\mathrm{distill}} = \frac{1}{HW} \sum_{i=1}^{HW} \| I_{\mathrm{teacher}}(r_i; [R|\mathbf{t}]) - I_{\mathrm{student}}(r_i; [R|\mathbf{t}]) \|_2^2
\end{equation}

其中 $[R|\mathbf{t}]$ 为渲染相机的旋转和平移，$H$、$W$~为图像高度和宽度。用白话解释：让低阶学生模型尽量复现高阶教师模型在各个视角下的渲染结果，从而在降低~SH~阶数的同时最大程度保留视角相关的外观效果。值得注意的是，蒸馏时的相机位置并非仅使用训练视角，而是在训练相机位置附近按高斯分布采样伪视角~$\mathbf{t}_{\mathrm{pseudo}} = \mathbf{t}_{\mathrm{train}} + \mathcal{N}(0, \sigma^2)$，以增强学生对未见视角的泛化能力。

LightGaussian~在~Mip-NeRF 360~数据集上实现了约~15$\times$~压缩率和~200+~FPS~渲染速度，同时保持了可接受的渲染质量。其贡献在于确立了"全局显著性评估~+~渐进剪枝~+~恢复微调~+~知识蒸馏"的基本范式，后续大量工作沿用了这一思路。

\paragraph{基于二阶灵敏度的原则性剪枝}

与~LightGaussian~的一阶累积贡献不同，PUP 3D-GS（2024.06）~\cite{2406_10219}~改用\textbf{二阶灵敏度分析}来评估每个高斯的剪枝影响。LightGaussian~的全局显著性分数回答的是"每个高斯对渲染贡献了多少颜色"，但这个问题隐含一个盲点：贡献大不代表移除后误差就大。如果周围有其他高斯能补偿它的贡献，那移除它的实际损失可能很小。PUP 3D-GS~转而回答一个更精确的问题："如果扰动某个高斯的参数，重建误差会增加多少？"

\textbf{从~Fisher~信息矩阵获取灵敏度。}考虑~$L_2$~重建损失关于高斯参数的~Hessian~矩阵（即二阶导数矩阵）$\nabla_G^2 L_2$。对~Hessian~做展开后会发现两项：一项是渲染图像梯度 $\nabla_G I_G$ 的外积，另一项含重建残差 $(I_G - I_{\mathrm{gt}})$ 的二阶导数。关键洞察是：在已收敛的~3DGS~模型上，残差趋近于零，因此第二项消失，Hessian~退化为：

\begin{equation}
\nabla_G^2 L_2 \approx \sum_{\mathbf{p}} \nabla_G I_G(\mathbf{p})\, \nabla_G I_G(\mathbf{p})^\top
\end{equation}

这正是统计学中的\textbf{Fisher~信息矩阵}——它只依赖于输入相机位姿（不依赖输入图像），因此计算成本远低于完整的~Hessian。注意，这一近似的前提是模型已经充分收敛。如果残差仍然较大，则~Fisher~近似会低估真实的灵敏度。

\textbf{分块灵敏度与空间参数选择。}完整的~Fisher~矩阵对所有~$G$~个高斯的参数求二阶导，规模为二次方，计算代价极高。PUP 3D-GS~发现\textbf{只需保留每个高斯自身的参数块}（即块对角近似~$\mathbf{H}_i = \nabla_{G_i} I_G\, \nabla_{G_i} I_G^\top$，忽略高斯间的交叉项）就足以衡量单个高斯的灵敏度。这相当于假设各高斯之间的影响是可分离的。进一步，论文发现\textbf{只需使用空间均值~$\mathbf{x}_i$~和缩放~$\mathbf{s}_i$~参数}即可获得有效的灵敏度分数，不需要旋转和不透明度参数。作者推测，这是因为场景到输入视角的投影几何不变性主要由位置和尺度决定，旋转在不变量存在时不会改变~3D~几何。最终的灵敏度分数定义为该参数块~Fisher~矩阵的\textbf{对数行列式}：

\begin{equation}
U_i = \log|\det(\nabla_{\mathbf{x}_i, \mathbf{s}_i} I_G\, \nabla_{\mathbf{x}_i, \mathbf{s}_i} I_G^\top)|
\end{equation}

直觉上，对数行列式衡量的是该高斯参数周围"误差盆地"（error basin）的体积——值越小，说明该高斯的参数可以在很大范围内变化而不影响重建误差，因此它很可能是冗余的。

\textbf{Patch-wise~计算与多轮管线。}为降低~Fisher~矩阵的计算开销，PUP 3D-GS~将渲染图像降采样为~$4\times4$~像素的~patch（即低分辨率渲染），在每个~patch~上独立计算~Fisher~近似后再对全部视角求和。实验表明，patch-wise~计算的灵敏度分数与全分辨率像素级计算高度相关。剪枝管线方面，PUP 3D-GS~沿用了~LightGaussian~的"剪枝~$\to$~微调恢复"基本流程，但将单轮扩展为\textbf{多轮~prune-refine}：第一轮剪除~80\%~低灵敏度高斯并微调~5000~步，第二轮在剩余高斯中再剪~50\%（累计~90\%~剪枝率）并再次微调。论文发现多轮逐步剪枝的效果一致优于同等总剪枝率的单轮剪枝。原因在于每轮微调后，剩余高斯会重新分配误差盆地，使得下一轮的灵敏度估计更准确。

\textbf{实验对比。}在~Mip-NeRF 360~数据集上~90\%~剪枝率（仅保留~10\%~高斯）的设定下，PUP 3D-GS~在几乎全部场景的全部指标上优于~LightGaussian。唯一的例外是~Tanks and Temples~数据集的~PSNR，PUP 3D-GS~低~0.36~dB。论文还指出，LightGaussian~的体积归一化项倾向于保留更大的高斯（因为大高斯在更多像素上可见），而~PUP 3D-GS~的灵敏度分析不存在这种偏差。实验中~LightGaussian~剪枝后场景的高斯协方差分布明显偏向大高斯，而~PUP 3D-GS~的分布更均匀。结合~LightGaussian~框架中的~Vectree（向量树量化，一种基于树结构的向量量化方法）量化，PUP 3D-GS~的最终模型在三个数据集上分别实现了~51.70$\times$、52.56$\times$、51.06$\times$~的压缩比。

\paragraph{从贡献权重到直接渲染误差}

前述方法（LightGaussian、PUP 3D-GS）各有其间接性。LightGaussian~用全局显著性分数（$\alpha_i T_i$~的累积）度量高斯的"颜色贡献量"，但贡献大不意味着移除后视觉误差就大。如果该高斯的颜色与它后方高斯的颜色相近，移除它几乎不影响画面。PUP 3D-GS~用~Fisher~信息矩阵度量参数的"灵敏度"，但它衡量的是参数扰动对重建误差的敏感程度，而非直接回答"移除这个高斯后画面会变差多少"。GaussianPOP（2026.02）~\cite{2602_06830}~的出发点正是这个缺口：能否从~3DGS~渲染方程出发，\textbf{解析地推导}出移除每个高斯后的\textbf{逐像素颜色误差}？

\textbf{核心推导：从渲染方程分解到解析误差。}GaussianPOP~的关键思想是将渲染方程在任意高斯~$k$~处拆分为三段：前景累积色~$P_{k-1}$（高斯~$k$~之前的所有颜色叠加）、高斯~$k$~自身的贡献~$T_k \alpha_k \mathbf{c}_k$、以及被高斯~$k$~衰减的背景色~$T_k(1-\alpha_k)\mathbf{b}_{k+1}$。这里~$\mathbf{b}_{k+1}$~是高斯~$k$~正后方所有后续高斯叠加出的"背景颜色"，定义为~$\mathbf{b}_{k+1} = \sum_{j=k+1}^{N} \alpha_j T_{k+1,j} \mathbf{c}_j / T_{k+1}$（即归一化后的后续累积色）。当移除高斯~$k$~时，它的贡献消失，但更关键的是背景色不再被~$\alpha_k$~衰减——后续高斯的透射率从~$T_{k+1} = T_k(1-\alpha_k)$~跳变为~$T_k$，因此背景色的权重从~$T_k(1-\alpha_k)$~增大为~$T_k$。将"移除前"与"移除后"的渲染颜色做差，前景项~$P_{k-1}$~完全抵消，最终得到\textbf{精确的逐像素平方误差}（不是上界近似）：

\begin{equation}
\mathrm{SE}_k(\mathbf{p}) = \| T_k \alpha_k (\mathbf{c}_k - \mathbf{b}_{k+1}) \|_2^2
\end{equation}

这个公式揭示了一个重要直觉：移除高斯~$k$~的视觉误差取决于两个因素的乘积。一是它的有效贡献权重~$T_k \alpha_k$（与~LightGaussian~的显著性分数一致），二是它的颜色~$\mathbf{c}_k$~与后方背景~$\mathbf{b}_{k+1}$~的差异。如果一个高斯贡献权重很大，但它的颜色与背景几乎相同（例如一朵白色云的内部高斯，被后方同为白色的其他高斯包围），那么移除它几乎不产生误差。这正是~LightGaussian~的显著性分数无法区分的情况。

\textbf{背景项的高效计算。}上述公式看似简洁，但隐含一个计算挑战：背景色~$\mathbf{b}_{k+1}$~依赖于高斯~$k$~之后所有高斯的累积，朴素方法需要对每个高斯分别做一次完整渲染，计算量为~$O(MV)$（$M$~个高斯~$\times$~$V$~个视角）。GaussianPOP~的关键算法贡献是\textbf{render-once, compute-locally}~策略：（1）执行一次标准前向渲染，生成每个像素的贡献高斯排序列表；（2）对排序列表做一次并行前缀求和（parallel prefix sum），同时缓存每个高斯的累积色~$P_k$~和透射率~$T_k$；（3）利用渲染最终色~$C_{\mathrm{render}}$~与缓存值的代数关系，反推出~$\mathbf{b}_{k+1} = (C_{\mathrm{render}} - P_k) / T_{k+1}$。这样在每个高斯的局部寄存器中独立完成误差计算，无需重新读取全局数据。这一策略将~$M$~次完整渲染替换为一次渲染加上并行局部操作，使误差量化阶段高度高效。

\textbf{简化框架。}GaussianPOP~的三阶段流程为：（1）在全部训练视角上运行上述误差量化，为每个高斯累积全局误差分数~$\Sigma\mathrm{SE}$；（2）按该分数全局排序，移除分数最低的~$P\%$~高斯；（3）微调恢复。值得注意的是，GaussianPOP~的误差准则\textbf{不需要真实图像}。它只度量"移除前后渲染结果自身的差异"，因此可以直接应用于任何预训练~3DGS~模型。GaussianPOP~同时支持训练中剪枝（按固定步长触发剪枝调度）和训练后简化（迭代式重新量化，每轮剪枝后重新计算误差分数再剪下一轮）。实验表明，在同等压缩率下，GaussianPOP~在几乎所有场景和指标上优于~LightGaussian~和~PUP 3D-GS，尤其在训练后简化的高压缩率设定下优势明显。这一结果验证了"直接度量渲染误差"比"间接度量颜色贡献或参数灵敏度"更准确地反映了剪枝的真实影响。

\subsubsection{基于可学习掩码的剪枝}

本路线介绍通过为每个高斯学习可微分掩码来在训练过程中自动决定保留或移除的剪枝方法。

基于重要性分数的剪枝在训练后一次性决定哪些高斯保留，但"重要性"是一个动态概念——训练早期不重要的基元可能在后期变得关键。基于可学习掩码的剪枝将保留/移除决策建模为优化问题，让训练过程自动发现最优的高斯子集。

\paragraph{概率性掩码剪枝}

前述基于重要性分数的方法（如~LightGaussian、GaussianPOP）都遵循"先训练、再评估、后剪枝"的范式。训练结束后根据某个快照下的统计量（不透明度、梯度贡献等）一次性决定哪些高斯保留。但"重要性"本质上是训练过程中的动态概念：训练早期不重要的基元可能在后期变得关键，训练后期冗余的高斯可能曾在早期承担过关键的粗结构建模角色。MaskGaussian（2024.12）~\cite{2412_20522}~将保留/移除决策从训练后的一次性判断变为训练中的持续优化。它不永久移除任何高斯，而是为每个高斯学习一个存在概率分布。

\textbf{掩码采样与存在概率。}具体地，MaskGaussian 为每个高斯维护~2~个可学习的掩码分数，通过~Gumbel-Softmax（一种可微的离散采样技巧）从中采样一个可微分的二元掩码~$M_i \in \{0, 1\}$，表示该高斯在当前迭代中是"存在"还是"缺席"。直觉上，这类似于对每个高斯进行持续投票——每个训练迭代都重新采样，让每个高斯有机会证明自己的价值，而非一次性被宣判"不重要"。

\textbf{掩码光栅化（masked rasterization）}是方法的核心机制。与将掩码直接乘到高斯属性（如缩放或不透明度）上的做法不同，MaskGaussian 保留了完整的高斯属性用于投影和~$\alpha$-过滤，只在光栅化阶段将掩码嵌入颜色累积和透射率演化中：

\begin{equation}
C(\mathbf{x}) = \sum_{i=1}^{N} M_i \, c_i \, \alpha_i \, T_i, \quad T_{i+1} = M_i (1-\alpha_i) T_i + (1-M_i) T_i
\end{equation}

当~$M_i = 1$~时，高斯正常贡献颜色并消耗透射率；当~$M_i = 0$~时，高斯的颜色贡献被屏蔽，透射率也不被消耗——相当于该高斯在渲染中"透明"。这一设计确保了前向渲染结果在缺席高斯处的正确性。

\textbf{梯度的关键作用}在于让被屏蔽的高斯仍然知道自己的潜在贡献。定义~$b_{i+1}$~为第~$i$~个高斯之后所有高斯渲染的颜色，则掩码的梯度为：

\begin{equation}
\frac{\partial L}{\partial M_i} = \alpha_i T_i \cdot \frac{\partial L}{\partial C(\mathbf{x})} \cdot (c_i - b_{i+1})
\end{equation}

用白话解释这个公式：梯度的符号取决于"使用这个高斯的颜色"比"使用它后面的颜色"是否更有利于降低损失。如果被屏蔽的高斯确实对渲染有帮助，它会收到正梯度，提高其存在概率；反之，如果其他高斯已经很好地补偿了它的缺席，梯度趋于零或负值，存在概率逐渐降低。值得注意的是，梯度公式中已经包含了~$\alpha_i T_i$——即基于重要性分数的剪枝方法所使用的指标——同时还额外捕获了被屏蔽高斯的颜色与场景需求之间的关系，这是分数方法无法衡量的。

\textbf{训练策略与剪枝。}为了控制最终的高斯数量，MaskGaussian 引入均方稀疏正则化项：

\begin{equation}
L_m = \left(\frac{1}{N} \sum_{i=1}^{N} M_i\right)^2, \quad L = L_{\text{render}} + \lambda_m L_m
\end{equation}

其中~$\lambda_m$~是平衡超参数，实验表明在不同训练阶段使用不同的~$\lambda_m$~效果最佳。剪枝时，对每个高斯采样~10~次，移除从未被采样的高斯。这一剪枝在每个致密化步骤时执行，致密化结束后每~1000~次迭代执行一次。实验显示，MaskGaussian~在移除超过~60\%~高斯的同时仅造成~0.02~dB~的~PSNR~下降。

SVR-GS（2025.09）~\cite{2509_11116}~进一步指出 MaskGaussian 的全局均值正则化~$L_m$~对所有高斯施加相同的稀疏压力，而像素级重建损失关注的是局部区域，二者之间存在错位。SVR-GS 提出空间变化正则化器，沿每条光线渲染像素级空间掩码，将稀疏压力施加在真正不重要的位置。相比 MaskGaussian，SVR-GS 进一步将高斯数量减少~1.79$\times$，PSNR~仅额外下降~0.50~dB。

\paragraph{学习自动剪枝}

前述方法（包括~MaskGaussian）都需要人工指定剪枝强度。MaskGaussian 通过超参数~$\lambda_m$~间接控制稀疏程度，基于重要性分数的方法则需要设定剪枝阈值或比例。然而，每个场景的最优剪枝比例差异显著：一个简单室内场景可能在~50\%~剪枝率下仍保持~PSNR~不变，而一个复杂的室外大场景在~30\%~剪枝率时就会出现明显的质量下降。这意味着，为每个场景寻找最优剪枝比例往往需要多轮试错训练，极大地增加了实际部署成本。LP-3DGS（2024.05）~\cite{2405_18784}~将剪枝比例本身变为可学习的量，让模型自动确定每个场景的最优剪枝程度。

\textbf{Gumbel-Sigmoid 可微掩码。}LP-3DGS~为每个高斯初始化一个实值掩码参数~$m_i$，然后通过~Gumbel-Sigmoid~函数将其二值化为~0/1~决策。与~Compact3D~等先前工作使用的直通估计器（STE，Straight-Through Estimator）不同，STE~在反向传播时直接跳过阈值函数的梯度，导致实值掩码和二值掩码之间存在不一致。训练后的~Sigmoid~输出散布在~$(0,1)$~区间，难以准确转换为~0/1~二值决策。Gumbel-Sigmoid~通过引入两个~Gumbel~噪声样本~$g_0, g_1$~实现可微的近二值化采样：

\begin{equation}
\text{gs}(m) = \frac{1}{1 + \exp\bigl(-(\sigma(m) + g_0 - g_1) / \tau\bigr)}
\end{equation}

其中~$\sigma(m) = \ln(m / (1-m))$~为~logit~变换，$\tau$~为温度参数。Gumbel~噪声的引入使得输出天然趋向~0~或~1，而非停留在中间值，因此可以更准确地近似二值掩码决策，同时保持函数完全可微。

\textbf{掩码与重要性分数的结合。}LP-3DGS~将掩码直接作用在重要性分数之上，而非像~MaskGaussian~那样嵌入光栅化过程。具体地，剪枝后的有效不透明度为：

\begin{equation}
o_i^{\text{mask}} = o_i \cdot \text{gs}(m_i \cdot S_i)
\end{equation}

其中~$S_i$~是预定义的重要性分数（如~RadSplat~的最大光线贡献或~Mini-Splatting~的累积权重）。这一设计使得~LP-3DGS~与多种已有的重要性度量兼容。通过~$L_1$~正则化~$R_{\text{mask}} = \frac{1}{N}\sum_{i=1}^{N}|m_i|$~鼓励掩码稀疏化，剪枝比例不再是人工指定的超参数，而是由~$L_1$~正则化强度~$\lambda_m$~间接决定的自动结果。实践中，掩码学习仅需~500~次迭代即可收敛，之后执行一次性剪枝并微调至原始训练总步数。

\subsubsection{基于几何一致性的剪枝}

本路线介绍从几何角度出发，利用空间一致性和多视角可见性来评估高斯重要性的剪枝方法。前述方法主要从渲染贡献的角度评估高斯重要性。另一条路线从几何角度出发：一个"好的"高斯应当具有空间一致性——它应该被其他高斯环绕，且在多个视角下都可见。

\paragraph{像素级剪枝的安全保障}

SafeguardGS（2024.05）~\cite{2405_17793}~系统性地将~3DGS~剪枝分为场景级剪枝和像素级剪枝两类，并揭示了它们在极端压缩下的截然不同的行为。两种剪枝的核心区别在于评估范围。场景级剪枝（如~LightGaussian、Mini-Splatting）将每个高斯的贡献在所有像素上聚合或取最大值，得到一个全局分数，然后按阈值或比例统一裁剪。像素级剪枝（如~EfficientGS）则独立地在每条光线~$r$~的范围内评估高斯贡献，逐像素标记保留哪些基元，最终通过逻辑"或"合并所有像素的掩码。

\textbf{场景级剪枝的局限}可以用一个直觉来理解：想象一个高斯均匀地对~100~条光线贡献极小（移除后几乎无影响），而另一个高斯仅对~1~条光线贡献巨大（移除后该像素将出现空白）。使用求和式的场景级分数（如~Mini-Splatting 的~$Score_{MS}(G_i, R) = \sum_{r \in R} w_i$），前者可能因覆盖面广而获得更高分数，反而在剪枝中被保留；后者则因仅服务一条光线而面临被移除的风险。更严重的是，场景级剪枝在极端压缩下缺乏安全下界。如果参数选择不当，可能某些光线完全没有被任何高斯覆盖，导致渲染出现空洞。论文实验证实了这一点：在~3DGS~训练~30k~迭代后于~20k~迭代处执行剪枝，场景级方法在保留约~10\% 高斯时~PSNR~急剧下降。像素级方法则始终保持至少约~35~万高斯的下界，避免灾难性质量坍塌。

\textbf{像素级剪枝的天然安全保障}在于其评估粒度。像素级方法保证每条光线至少保留一个高斯，提供了自然的剪枝下界——即使采用最激进的~top-1~设置（每条光线仅保留贡献最大的一个高斯），场景也不会出现空白。这是场景级剪枝无法提供的保障。

\textbf{SafeguardGS 的核心贡献}是将最优的像素级剪枝策略与精心设计的分数函数相结合。论文从~3DGS~渲染方程出发，系统探索了由颜色相似度（~$CS_1, CS_2, CS_3$~）、像素接近度（~$PC$~）、混合权重（~$w_i$~）等因子组合而成的多种分数函数。其中颜色相似度和像素接近度的定义如下：

\begin{align}
PC(X_r, X_i) &= \exp(-\|X_r - X_i\|_2) \\
CS_1(c_r, c_i) &= 1 - \text{mean}(\|c_r - c_i\|_1) \\
CS_2(c_r, c_i) &= \exp(-\text{mean}(\|c_r - c_i\|_1))
\end{align}

其中~$X_r$~和~$c_r$~分别是光线~$r$~对应的像素坐标和颜色，$X_i$~是高斯投影~$G_i$~的中心位置。这些函数将各因子的值域压缩到~$[0,1]$，避免某一因子在分数计算中占据主导。通过在~Synthetic NeRF~和~Mip-NeRF 360~等数据集上的大量~top-k~参数扫描实验，论文发现~$Score_{SG\_36}$（结合颜色相似度~$CS_2$~和混合权重~$w_i$）在每基元表达能力（PSNR-per-Primitive）上始终领先其他函数——它能完全消除与真实像素颜色差异大的低质量高斯，同时保留对高频区域（如地板纹理）至关重要的基元。

最终，SafeguardGS~使用~$Score_{SG\_36}$~配合~top-1~准则逐像素剪枝，将场景压缩约~90\%。在~Mip-NeRF 360~真实场景上，SafeguardGS~仅造成~0.034~dB~的~PSNR~下降，同时达到最高的每基元表达能力（0.8E-4~PSNR/Prim.），显著优于同等压缩率下的场景级方法。

\paragraph{语义引导的浮游体移除}

前述~SafeguardGS~通过像素级渲染贡献来评估高斯重要性——一个高斯是否值得保留，取决于它在渲染过程中对像素颜色的贡献大小。然而，3DGS~训练中还存在一类特殊的退化高斯——浮游体（floaters）：它们远离实际物体表面，悬浮在场景的空白空间中，呈现为半透明碎片。浮游体的独特性在于，它们对训练视角的渲染影响微乎其微（因为被前景遮挡或落在背景区域），因此基于渲染贡献的剪枝方法难以有效识别它们；但在新视角下，这些浮游体暴露出来，显著降低渲染质量。Clean-GS（2026.01）~\cite{2601_00913}~从语义角度切入这一问题：不再依赖渲染贡献，而是利用少量语义分割掩码判断一个高斯是否属于目标物体。

\textbf{语义掩码引导的三阶段过滤。}Clean-GS~的输入是一个已训练的~3DGS~模型和~$M$~张带二值语义掩码的视角（掩码标注目标物体区域）。核心流程分为三个逐层收紧的阶段。第一阶段\textbf{白名单过滤}：将每个高斯投影到所有掩码视角上，若它在任何视角中投影到目标物体区域（白色像素），则加入白名单；否则直接移除。这一阶段利用空间投影证据快速剔除背景和环境高斯。第二阶段\textbf{深度缓冲颜色验证}：白名单中仍可能包含投影位置正确但颜色不匹配的浮游体。对于每个掩码视角，仅渲染白名单高斯并通过深度缓冲取前层，然后比较每个前层高斯的~DC~颜色与对应像素的实际颜色——若欧氏距离超过阈值~$\theta = 0.40$~（RGB~空间），则将其移除。该阶段保守地仅移除有明确颜色失配证据的高斯，对于在所有掩码视角中都被遮挡的高斯则予以保留。第三阶段\textbf{邻域异常值移除}：经过前两阶段后，残余伪影通常在~3D~空间中孤立分布。Clean-GS~计算每个高斯到其~$k = 10$~个最近邻的平均距离，移除超过第~95~百分位数的离群点，从而清除空间上与主体物体脱节的孤立浮游体。

\textbf{与~SafeguardGS~的关键差异。}SafeguardGS~在渲染管线内通过分数函数和~top-k~准则逐像素剪枝，属于纯几何贡献评估；Clean-GS~则引入外部语义信息（分割掩码），从"这个高斯是否属于目标物体"的角度做判断，而非"这个高斯对渲染贡献多大"。在标注成本上，Clean-GS~仅需~3~张语义分割掩码（由~SAM~自动生成，约占~302~张训练视图的~1\%），即可提供足够的语义信号完成物体分离。在~Tanks and Temples~数据集的~Temple~场景上，Clean-GS~将~525{,}717~个高斯压缩至约~198{,}000~个（约~62\%~压缩率），文件大小从~125~MB~降至~47~MB，同时渲染质量基本保持。整个处理流程在~96~核~CPU~上仅需约~37~秒（其中邻域移除阶段额外消耗~16~秒），无需~GPU~参与。

\subsubsection{基于最优传输的全局简化}

本路线介绍将高斯简化建模为全局分布匹配问题的最优传输方法。前述方法以逐高斯的方式评估重要性并进行剪枝，缺乏全局视角的优化保证。一类新兴方法将高斯简化建模为全局分布匹配问题。

\paragraph{全局高斯混合简化}

前述方法（LightGaussian、SafeguardGS~等）均以逐高斯的方式评估重要性并移除低分基元。这类策略存在根本局限：每次移除决策只考虑单个高斯的局部贡献，无法保证简化后的高斯集合在全局分布意义下最优——正如逐个丢弃包裹无法保证剩余包裹的总重量最接近原始总重。Gaussian Herding（2025.06）~\cite{2506_09534}~从根本上改变了简化的范式：不再"移除"高斯，而是将~$N$~个高斯的信息重新分配（redistribute）到~$K \ll N$~个高斯中。这类似于快递合并——不是扔掉包裹，而是将多个小包裹的内容整合到少数大包裹中，使得最终送达的货物总量不变。

\paragraph{概率建模与最优传输框架}

Gaussian Herding 的出发点是将~3DGS~的几何建模为非归一化高斯混合（unnormalized Gaussian mixture）：每个高斯基元的~PDF~为~$\phi_i(x) = |\Sigma_i|^{-1} \exp(-\frac{1}{2}(x - \mu_i)^\top \Sigma_i^{-1}(x - \mu_i))$，$N$~个高斯共同构成场景的隐式表面函数~$f(x) \approx \sum_{i=1}^N \phi_i(x)$。在此基础上，简化被建模为高斯混合简化（Gaussian Mixture Reduction, GMR）问题：寻找一个仅含~$m$~个分量的高斯混合~$\tilde{f}_m(x) = \sum_{j=1}^m \tilde{\phi}_j(x)$，使其在最优传输意义下最接近原始~$f(x)$。

\paragraph{复合传输散度与几何简化}

衡量两个高斯混合之间距离的核心工具是复合传输散度（Composite Transportation Divergence, CTD）。直觉地说，CTD~将最优传输从"搬运质量点"推广到"搬运整个高斯分布"：每个高斯分量被视为空间中的一个离散单元，CTD~计算将原始混合中的所有高斯"搬运"到目标混合所需的最小代价：

\begin{equation}
\mathcal{T}_c(f_N, \tilde{f}_m) = \min_{\pi_{ij}} \left\{ \sum_{i=1}^N \sum_{j=1}^m \pi_{ij} \, c\bigl((\cdot;\mu_i,\Sigma_i),\, (\cdot;\tilde{\mu}_j,\tilde{\Sigma}_j)\bigr) \right\}
\label{eq:ch11_ctd}
\end{equation}

其中~$\pi_{ij}$~为传输计划（transport plan），表示从原始第~$i$~个高斯搬运到目标第~$j$~个高斯的质量，需满足边际约束~$\sum_j \pi_{ij} = w_i$、$\sum_i \pi_{ij} = \tilde{w}_j$；$c(\cdot,\cdot)$~是衡量两个高斯分布之间距离的代价函数。该公式可理解为：在所有可行的"搬运方案"中，选择总代价最小的那个。

直接在整个场景上优化~CTD~的计算代价为~$O(N^2)$，对~$N = 10^5$~量级的~3DGS~场景不可接受。Gaussian Herding 引入两项关键优化。第一，基于~KD-Tree~的分块策略（Blockwise GMR）：仅用高斯中心~$\{\mu_i\}_{i=1}^N$~构建~KD-Tree（因为观察到协方差矩阵的特征值通常很小，高斯中心足以刻画空间分布），每次分裂取中位数坐标值，在深度~$d = \lfloor \log_2(N/s) \rfloor$~处产生~$2^d$~个叶块，每块至多含~$s$~个高斯（$s \ll N$），然后在每个叶块内独立执行~GMR，将~$s$~个高斯减少到~$m = \lceil s \cdot r \rceil$~个（$r$~为保留率）。分块后总复杂度从~$O(N^2)$~降至~$O(s^2 \cdot N)$——以~$N = 10^5$、$s = 10^3$、$r = 0.05$~为例，每块迭代开销从~$10^8$~降至约~$10^5$~次运算，且各块可并行。

第二，高效代价函数。Zhang~等（2025.06）~\cite{2506_09534}~指出，直接用~KL~散度作为~$c(\cdot,\cdot)$~虽在理论上最优，但需要~$O(s^3 n)$~次协方差矩阵求逆且存在数值不稳定问题。他们设计了如下高效替代：

\begin{equation}
c\bigl((\cdot;\mu,\Sigma),\, (\cdot;\mu',\Sigma')\bigr) = \|\mu - \mu'\|_2^2 + \|\Sigma - \Sigma'\|_F^2
\label{eq:ch11_efficient_cost}
\end{equation}

即代价等于均值向量的~$L^2$~距离加上协方差矩阵的~Frobenius~范数距离。该代价函数无需矩阵求逆，且更新步骤退化为加权平均：$\tilde{\mu}_j^{(t)} = \frac{\sum_i \pi_{ij} w_i \mu_i}{\sum_i \pi_{ij} w_i}$，$\tilde{\Sigma}_j^{(t)} = \frac{\sum_i \pi_{ij} w_i \Sigma_i}{\sum_i \pi_{ij} w_i}$。整个~GMR~求解过程本质上是高斯空间中的~$k$-means~聚类变体：分配步骤将每个高斯分到最近的聚类中心，更新步骤将中心替换为聚类内高斯的加权~barycenter。

\paragraph{几何与外观解耦优化}

Gaussian Herding 将简化分为两个解耦阶段。第一阶段（几何简化）：通过上述~Blockwise GMR~得到~$m$~个紧凑高斯的几何参数~$(\tilde{\mu}_j, \tilde{\Sigma}_j)$，此阶段仅修改位置和协方差，不触及外观属性。第二阶段（外观优化）：为每个简化高斯从原始集合中找最近邻（$i^* = \arg\min_i \|\mu_i - \tilde{\mu}_j\|_2$）初始化颜色~$c_j$~和不透明度~$\alpha_j$，然后在标准~3DGS~渲染管线中冻结几何参数，仅通过反向传播微调外观属性~$T'$~步。解耦的设计动机在于：GMR~输出的混合权重~$\tilde{w}_j$~不一定满足不透明度~$\alpha_j \in [0,1]$~的约束，直接使用会导致渲染伪影；通过额外的外观微调，可以修正这一偏差并恢复高频细节。

\paragraph{实验表现}

实验表明，Gaussian Herding 仅保留~10\%~的高斯即可达到与原始~3DGS~几乎相同的渲染质量（在~PSNR、SSIM、LPIPS（ Learned Perceptual Image Patch Similarity）等多项指标上差异均可忽略不计）。其关键优势在于"信息重分配"而非"移除"：被合并的高斯信息通过~CTD~优化融入了保留下来的高斯的位置和形状参数之中，而非简单地丢弃。这一机制很好地解释了为何在同等压缩率下，基于全局最优传输的方法往往优于逐个剪枝的方案。

MMGS（2026.05）~\cite{2605_19304}~进一步将多视角一致性引入全局简化，通过多视角~3D~高斯贡献排序机制过滤基元，再用最优传输聚合算法合并冗余基元。该方法仅用~10\%~基元即可达到~SOTA~渲染质量，同时训练速度提升~10$\times$。

\paragraph{免训练的高斯合并简化}

前述~Gaussian Herding 通过全局最优传输将~$N$~个高斯的信息重分配到~$K$~个高斯中，虽然保证了全局分布最优，但需要~GPU~上的渲染监督来微调外观。NanoGS（2026.03）~\cite{2603_16103}~走了一条完全不同的路线——免训练的局部成对合并。该方法不依赖任何训练图像或~GPU~计算，直接在~CPU~上对已有的~.ply~文件执行简化，其核心动机在于：边缘设备（移动端、嵌入式平台）通常不具备~GPU~和训练数据，但仍有实时渲染简化模型的需求，因此简化方案必须做到"开箱即用"——给定一个已训练好的~.ply~文件，无需任何额外信息即可输出简化版本。

\paragraph{稀疏合并图与成对合并范式}

NanoGS 将简化建模为图上的渐进式边坍缩（progressive edge collapse）问题。首先，对无结构的~$N$~个高斯基元构建~$k$-近邻（KNN，K-Nearest Neighbors）稀疏合并图~$G = (V, E)$：每个节点对应一个高斯基元，每条边连接空间上最近邻的~$k$~个候选对。这使候选合并对的数量从全部配对的~$O(N^2)$~降至~$O(kN)$，为大规模场景上的可扩展简化奠定基础。简化的执行流程类似于网格简化中的边坍缩：每轮从图中选择一组不相交（disjoint）的低代价边，将各边的两个端点高斯合并为一个新高斯，然后局部刷新邻域关系以发现新的候选对，如此迭代直至达到目标压缩率。

\paragraph{合并代价：几何失真与外观差异}

为评估两个高斯~$g_i$~和~$g_j$~是否值得合并，NanoGS~设计了包含几何与外观两个分量的合并代价函数。几何分量衡量"用单个高斯近似双高斯混合"引入的失真。具体而言，将每个高斯基元建模为非归一化高斯质量密度~$p_i(x) = w_i \, \mathcal{N}(x; \boldsymbol{\mu}_i, \boldsymbol{\Sigma}_i)$（其中质量权重~$w_i$~同时取决于基元的空间尺度与不透明度），则合并代价定义为混合密度~$p_i + p_j$~与合并后单高斯~$q_m$~之间的~I-散度（I-divergence）：$D_{\text{merge}}(i,j) = I(p_i + p_j \| q_m)$。由于双高斯混合对单高斯的~I-散度没有闭式解，NanoGS 采用蒙特卡罗采样（每对少量固定样本）进行数值估计。外观分量则定义为两个高斯外观特征向量（如球谐系数）之间的~$L^2$~距离：$D_{\text{app}}(i,j) = \| \mathbf{f}_i - \mathbf{f}_j \|_2^2$。总代价~$C(i,j) = D_{\text{merge}}(i,j) + D_{\text{app}}(i,j)$~仅在几何接近且外观相似时取低值，从而自然过滤掉不应合并的基元对。

\paragraph{质量保持矩匹配（MPMM）}

确定合并对象后，核心问题是如何将两个高斯融合为一个。NanoGS~提出了质量保持矩匹配（Mass-Preserved Moment Matching, MPMM）算子。与经典的矩匹配不同，MPMM 以质量权重（mass weight）而非简单的不透明度作为加权依据——质量权重~$w_i = (2\pi)^{3/2} \sqrt{|\boldsymbol{\Sigma}_i|} \cdot \alpha_i$~同时编码了基元的空间体积（协方差行列式）和不透明度，使空间范围更大或更不透明的高斯基元在合并中拥有更大的话语权。合并后的参数由质量加权矩确定：

\begin{align}
W &= w_i + w_j, \quad \boldsymbol{\mu}_m = \frac{w_i \boldsymbol{\mu}_i + w_j \boldsymbol{\mu}_j}{W} \label{eq:ch11_nanogs_mean} \\
\boldsymbol{\Sigma}_m &= \frac{1}{W}\!\left[\sum_{k \in \{i,j\}} w_k \boldsymbol{\Sigma}_k\right] + \frac{1}{W}\!\left[\sum_{k \in \{i,j\}} w_k (\boldsymbol{\mu}_k - \boldsymbol{\mu}_m)(\boldsymbol{\mu}_k - \boldsymbol{\mu}_m)^\top\right] \label{eq:ch11_nanogs_cov}
\end{align}

其中第一项为质量加权协方差均值（保留平均形状），第二项为均值离散项（膨胀协方差以覆盖两个高斯之间的间隙）。这一设计的关键直觉是：合并后的高斯既需保留原始高斯的典型局部形状，又需扩大到覆盖两者的联合空间范围，避免在剧烈压缩时出现空洞。不透明度则采用~Porter-Duff~"source-over"合成规则~$\alpha_m = 1 - (1-\alpha_i)(1-\alpha_j)$，即两个独立覆盖事件的概率并，防止不透明度无界增长。

\paragraph{与~Gaussian Herding 的关键差异}

NanoGS 与~Gaussian Herding 虽然都追求"信息保持"而非"信息丢弃"，但技术路线截然不同。Gaussian Herding 以~CTD~全局最优传输为框架，在~KD-Tree~分块内执行~$k$-means~风格的高斯聚类，需要~GPU~上的外观微调来修正简化引入的偏差。NanoGS 则完全免除训练环节：（1）~MPMM~的质量加权矩匹配在合并操作中同时完成几何参数和外观特征的聚合，无需后续微调；（2）所有计算（KNN~图构建、代价评估、边坍缩）均在~CPU~上执行，直接输入~.ply~文件即可输出简化结果；（3）渐进式批量边坍缩天然支持任意目标压缩率，用户可在运行时动态选择停止点。这使得~NanoGS 特别适合在资源受限的部署场景中使用——例如，在移动设备上接收一个完整的~3DGS~场景文件后，按设备算力选择合适的压缩率进行本地简化，再执行实时渲染。

\begin{takeaway}
高斯剪枝与稀疏化通过减少高斯基元数量降低存储和渲染开销。技术路线从简单的基于重要性分数的剪枝发展到信息保持的合并策略（NanoGS 的质量加权矩匹配），核心权衡是压缩率与渲染质量的平衡。NanoGS 免除训练环节、纯 CPU 执行的特点使其特别适合资源受限的部署场景。
\end{takeaway}


%=====================================================================

\subsection{近期方法脉络：后训练压缩、准确紧凑高效表示、剪枝与重要性评估等}

这一组近期工作延续 3DGS、3DGS pruning、dynamic Gaussian compression，主要把改动集中在压缩、优化、加速、增量改进等环节。下面只保留相对既有路线的变化，便于和前文方法直接对照。

\paragraph{Smaller and Faster 3DGS via Post-Training Dictionary Learning}
Smaller and Faster 3DGS via Post-Training Dictionary Learning（2026.05）~\cite{2605_30396} 延续 3DGS 在“后训练压缩”上的基本设定，主要改变是：基于字典学习的3DGS后训练压缩框架，用共享字典表示高斯属性并避免重新训练原模型。3DGS 的显式高斯表示和可微渲染能力保持为基础。
其设计动机是：提出首个面向3DGS的字典学习式后训练压缩管线，可直接作用于已有3DGS模型，无需重训或改动架构，并在多个基准场景上兼顾压缩率、画质和渲染速度。
具体而言，方法把训练好的高斯属性映射到共享字典和稀疏/低维编码，通过后处理方式替代逐高斯独立参数存储；由于不改变原始3DGS训练流程，它可以作为模型发布前的压缩步骤接入多种3DGS变体。
适用边界是：压缩效果依赖字典容量与场景属性分布；极高压缩率下仍可能牺牲细节；主要验证静态场景。

\paragraph{ACE-GS}
ACE-GS（2026.06）~\cite{2606_21244} 延续 3DGS 在“准确紧凑高效表示”上的基本设定，主要改变是：ACE-GS：渐进式优化和精细primitive management框架。3DGS 的显式高斯表示和可微渲染能力保持为基础。
其设计动机是：把高斯管理视为突破质量-速度-压缩折衷的核心，避免过度激进剪枝导致高频细节丢失。
具体而言，方法采用渐进优化策略，在训练不同阶段调整高斯生长、剪枝和属性优化，使表示逐步紧凑且保持细节。
适用边界是：具体收益依赖场景复杂度和管理策略超参数；与专用压缩方法的组合仍需研究。

\paragraph{REFINE}
REFINE（2026.06）~\cite{2606_09074} 延续 3DGS pruning 在“剪枝与重要性评估”上的基本设定，主要改变是：REFINE：rendering-free primitive importance，用解析近似的rendering-aware Hessian field估计删除高斯的感知误差。3DGS 的显式高斯表示和可微渲染能力保持为基础。
其设计动机是：提出渲染无关但渲染感知的primitive importance metric，建模可见性、投影几何和内容自适应权重，绕开昂贵的逐候选forward rendering。
具体而言，方法解析近似每个高斯移除对感知误差的影响，形成各向异性感知权重场，并据此排序/剪枝primitive。
适用边界是：Hessian近似可能在复杂遮挡/视角依赖外观下失准；剪枝阈值仍需控制。

\paragraph{VEDAL}
VEDAL（2026.06）~\cite{2606_02346} 延续 3DGS pruning 在“剪枝与稀疏化”上的基本设定，主要改变是：VEDAL：Variational Error-Driven Asynchronous Learning，将高斯剪枝表述为变分自由能最小化，使用误差门控和不确定性头。3DGS 的显式高斯表示和可微渲染能力保持为基础。
其设计动机是：把3DGS pruning从启发式重要性评分推进到信息论式变分目标，引入逐高斯预测误差门控和可学习先验的不确定性剪枝决策。
具体而言，方法为高斯是否保留建模潜变量，通过free energy目标同时惩罚重建误差和复杂度；异步机制只在高误差/不确定区域触发剪枝更新。
适用边界是：引入不确定性头和变分目标增加训练复杂度；超参数对压缩-质量平衡敏感。

\paragraph{Renderable Partial Representations for Dynamic Gaussian Splatting under Incomplete Delivery}
Renderable Partial Representations for Dynamic Gaussian Splatting under Incomplete Delivery（2026.06）~\cite{2606_17212} 可作为“流式传输与部分表示”中的新锚点候选。相较于 dynamic Gaussian compression，主要改变是：Renderable Partial Representations：把动态高斯组织为可独立寻址的时空cluster，每个cluster有base和refinement层。3DGS 的显式高斯表示和可微渲染能力保持为基础。
其设计动机是：把动态高斯压缩从完整文件/前缀优化转向真实交互传输中的partial state优化，使缺失局部时空数据时仍能渲染合理画面。
具体而言，高斯primitive按时空cluster组织，训练时采样部分依赖图并渲染多种反事实交付状态，用图像空间损失优化不完整表示质量。
这一锚点判断来自当前分类中的机制关系：候选锚点：扩展，且相关性 0.85；后续方法可围绕 从完整前缀优化扩展到可渲染部分状态 继续比较和扩展。
适用边界是：系统和网络条件复杂；cluster设计影响质量和传输效率。
\section{编码与压缩}
\label{sec:11_2}

\subsection{问题背景}

本节讨论五条技术路线：（1）量化与向量量化；（2）熵编码与上下文建模；（3）前馈压缩与通用编解码器；（4）后训练压缩管线；（5）基于图像/视频编码器的压缩。

剪枝通过减少高斯数量来降低存储，但每个保留高斯的~59~维参数向量仍然以~32~位浮点数存储，存在大量冗余。编码与压缩的目标是进一步减少每个高斯的参数存储开销。直觉地说，如果将所有高斯的参数排列成一个~$N \times 59$~的矩阵，压缩可以从三个维度展开：（1）\textbf{量化}，即降低每个参数的位宽（从~32~位降至~8~位甚至更低）；（2）\textbf{预测}，即利用高斯之间的空间相关性，将绝对值编码转为残差编码；（3）\textbf{熵编码}，即利用参数值的统计分布特性进行无损压缩。

3DGS~压缩与传统的图像/视频压缩存在本质区别。图像和视频的像素排列在规则的~2D~网格上，自然适合变换编码和上下文建模。而~3DGS~的高斯分布是高度不规则的~3D~点云，缺乏天然的~2D~邻域结构。如何在这种不规则数据上建立有效的上下文模型，是~3DGS~压缩的核心挑战。

\subsubsection{量化与向量量化}

本节介绍通过降低参数位宽和码本编码来压缩高斯属性的方法，从奠基性的~Compact 3DGS~框架出发，沿自适应位宽、码本直接训练和格向量量化等方向展开。

\paragraph{开创性的量化与向量量化框架}

Compact 3D Gaussian Representation（2023.11）~\cite{2311_13681}~是~3DGS~压缩领域的奠基性工作，建立了"可学习掩码+网格神经场+码本向量量化（VQ，Vector Quantization）"的三阶段压缩范式，后续大量工作在此基础上演进。

\textbf{问题分析：参数冗余的三个来源。}
原始~3DGS~每个高斯存储~59~维、32~位浮点参数：位置~3~维、不透明度~1~维、旋转~4~维（四元数）、缩放~3~维、球谐系数~48~维（3~阶、16~个基函数$\times$3~个颜色通道）。一个典型大场景（如~Mip-NeRF 360）常常产生数百万个高斯，存储开销超过~1~GB。仔细分析可以发现冗余来自三个层面：（1）许多高斯体积微小，对渲染贡献可忽略——\textbf{数量冗余}；（2）相邻高斯的颜色高度相似——\textbf{空间冗余}；（3）大部分高斯的几何形状（缩放和旋转）变化有限——\textbf{属性冗余}。Compact 3DGS~的三个阶段分别瞄准这三类冗余。

\textbf{第一阶段：基于体积的可学习掩码（消除数量冗余）。}
原始~3DGS~在训练过程中不断克隆和分裂高斯以细化细节，但产生大量冗余基元。Compact 3DGS~为每个高斯引入一个连续掩码参数~$m_n \in \mathbb{R}$，经~sigmoid~和阈值化生成二值掩码：

\begin{equation}
    M_n = \mathrm{sg}\bigl(\mathbb{1}[\sigma(m_n) > \tau] - \sigma(m_n)\bigr) + \sigma(m_n),
\end{equation}

其中~$\sigma$~为~sigmoid~函数，$\tau$~为掩码阈值，$\mathrm{sg}(\cdot)$~为~stop-gradient~算子。这里用到了\textbf{直通估计器}（Straight-Through Estimator, STE）技巧——前向传播正常取整为~0/1，反向传播时梯度绕过硬阈值，直接流过~sigmoid，从而让离散的掩码决策变得可微分。掩码同时作用于缩放~$\hat{s}_n = M_n \cdot s_n$~和不透明度~$\hat{o}_n = M_n \cdot o_n$，兼顾了高斯的体积和透明度两个维度的贡献。训练中配合掩码正则化损失~$\mathcal{L}_m = \frac{1}{N}\sum_{n=1}^{N} \sigma(m_n)$~平衡渲染精度与高斯数量，每次致密化迭代后直接移除~$M_n = 0$~的高斯。掩码完成后无需存储掩码参数本身，因为被掩码的高斯已经从场景中删除。在~Playroom~场景上，仅此一项即可带来~140\%~的存储效率提升和~65\%~的渲染加速。

\textbf{第二阶段：网格神经场替代球谐系数（消除空间冗余）。}
在~59~维参数中，球谐系数占~48~维（约~80\%），是存储的大头。然而相邻高斯的颜色往往高度相似，逐高斯存储独立球谐系数存在大量空间冗余。Compact 3DGS~的思路是：不再为每个高斯显式存储颜色参数，而是构建一个所有高斯共享的轻量级~3D~特征网格，按需查询。具体而言，采用~Instant-NGP~的\textbf{哈希网格}（Hash Grid）编码高斯位置，取出特征向量后与视角方向~$\mathbf{d}$~拼接，送入浅层~MLP~解码为颜色：

\begin{equation}
    c_n(\mathbf{d}) = f\bigl(\mathrm{contract}(\mathbf{p}_n),\, \mathbf{d};\, \Theta\bigr),
\end{equation}

其中~$\mathrm{contract}(\cdot)$~借鉴~Mip-NeRF 360~将无界空间压缩到有界范围内（$|\mathbf{p}| \leq 1$~保持不变，$|\mathbf{p}| > 1$~映射到~$[1, 2)$~球壳），$\Theta$~为~MLP~参数。直觉上，这就像一张隐式的~3D"颜色地图"：不需要为每个点单独记录颜色，只需知道它在空间中的位置，就能查到对应特征并解码出视角相关的颜色。消融实验表明，即使不减少高斯数量，仅此替换即可获得超过~3~倍的存储效率提升。

\textbf{第三阶段：残差向量量化压缩几何属性（消除属性冗余）。}
旋转和缩放参数虽然只有~7~维（四元数~4~维+缩放~3~维），但观察到同一场景中大部分高斯的几何形状极为相似——场景由大量微小高斯拼合而成，每个单独高斯并不需要太大的多样性。

\textbf{向量量化}（Vector Quantization, VQ，一种将连续向量映射到有限码本中最近码字的离散化技术）的直觉类似"字典编码"：与其存储每个高斯的完整参数向量，不如学习一个包含~$K$~个典型形状的"码本"，每个高斯只需记录一个整数索引指向最相似的码字。

Compact 3DGS~进一步采用\textbf{残差向量量化}（Residual VQ, R-VQ）——将~$L$~级~VQ~串联：第一级码本捕捉粗粒度形状，后续各级码本逐级逼近残差。以旋转为例，第~$l$~级的量化结果为：

\begin{equation}
    \hat{\mathbf{r}}_n^l = \sum_{k=1}^{l} Z^k[i_n^k], \quad i_n^l = \arg\min_k \|Z^l[k] - (\mathbf{r}_n - \hat{\mathbf{r}}_n^{l-1})\|_2^2,
\end{equation}

其中~$Z^l \in \mathbb{R}^{C \times 4}$~是第~$l$~级的码本（大小为~$C$），$i_n^l$~为第~$n$~个高斯在第~$l$~级选择的码字索引，$\mathbf{r}_n$~为原始旋转向量。直觉上，第一级码本存储"典型旋转的大类"，第二级修正残差使其更精确，以此类推。码本的训练目标使量化误差最小化：

\begin{equation}
    \mathcal{L}_r = \frac{1}{NC} \sum_{k=1}^{L} \sum_{n=1}^{N} \|\mathrm{sg}[\mathbf{r}_n - \hat{\mathbf{r}}_n^{k-1}] - Z^k[i_n^k]\|_2^2,
\end{equation}

其中~$\mathrm{sg}[\cdot]$~为~stop-gradient~算子，确保码本能从残差中学习而非被梯度绕过。缩放参数~$s$~使用相同的~R-VQ~流程。整个框架端到端联合训练，总损失为：

\begin{equation}
    \mathcal{L} = \mathcal{L}_{\mathrm{ren}} + \lambda_m \mathcal{L}_m + \mathcal{L}_r + \mathcal{L}_s,
\end{equation}

其中~$\mathcal{L}_{\mathrm{ren}}$~为~L1+SSIM~渲染损失，$\lambda_m$~控制掩码强度。R-VQ~仅在最后~1000~次迭代中激活（以~K-means~初始化码本），避免训练早期的计算开销。消融实验表明，几何码本单独可带来约~30\%~的存储压缩，且对渲染质量几乎无影响。

\textbf{压缩率分解与后处理。}
综合三个阶段，Compact 3DGS~在各数据集上平均实现约~15$\times$~的端到端压缩比。此外，对不透明度和哈希网格参数进行~8~位最小最大量化，对~R-VQ~索引和量化参数进行~Huffman~编码，可额外压缩~40\%以上，最终总压缩比超过~25$\times$。以~Deep Blending~数据集为例，Compact 3DGS~在~PSNR~上甚至略优于原始~3DGS，同时存储效率提升~28$\times$、渲染速度提升近~40\%。

\paragraph{Content-Aware Radiance Fields：自适应位宽量化}

针对现有辐射场量化方法依赖人工试错选择位宽、无法在降低位宽同时保证渲染质量可控的局限，Content-Aware Radiance Fields~提出让位宽分配与场景复杂度对齐的自适应方案。Compact 3DGS~对所有属性通道使用固定位宽（如~8~位最小最大量化），未考虑不同场景区域对精度的差异化需求——纹理丰富的墙面需要高位宽保持细节，而平坦天空即使~4~位也冗余。Content-Aware Radiance Fields（2024.10）~\cite{2410_19483}~将位宽本身变为可训练的连续变量：引入浮点``软位宽''$b$（训练时连续，推理时取整为~$B = \lfloor b \rfloor$），并将量化范围参数~$r_v$、偏移~$v_\text{max}$~和软位宽~$b$~全部设为可学习参数，通过梯度下降联合优化。为解决``降低位宽必然损害渲染质量''这一对抗性矛盾，作者设计了对抗式内容感知量化（A-CAQ）：位宽学习损失~$\mathcal{L}^b = \| \mathcal{L}^\text{NeRF} - \mathcal{L}^\text{metric} \| + \sum_i \omega_i B_i$~以渲染损失与目标精度阈值的差距为位宽搜索空间提供指引，加权惩罚项~$\omega_i$~按组件重要性分配——例如对可学习码本施加更高惩罚以优先压缩其位宽。相比~Compact 3DGS~的固定~8~位方案，Content-Aware Radiance Fields~在~Instant-NGP~基线上将位宽从~8~位压缩至~4--6~位范围（不同层不同），同时保持~PSNR~下降不超过~0.2\,dB。

\paragraph{ContraGS：码本直接训练} 针对``先训练后压缩''范式带来的训练中期资源浪费问题——Compact 3DGS 的 R-VQ 码本仅在完整模型训练完毕后的最后 1000 次迭代中才以 K-means 初始化引入，导致训练中期所有高斯参数均以完整连续形式存在，需消耗大量 GPU 显存与计算资源，而这些本可在压缩表示上提前节省。Compact 3DGS~的~R-VQ~码本是在完整模型训练完毕后的最后~1000~次迭代中才引入的（以~K-means~初始化），这意味着先训练一个完整模型再压缩——浪费了大量存储和计算。ContraGS（2025.09）~\cite{2509_03775}~的核心改动是将参数估计建模为贝叶斯推断问题，在码本压缩表示上直接训练，跳过``先训练完整模型再压缩''的两步流程。具体而言，每个高斯的属性不再是连续向量，而是指向码本~$SH$（球谐系数）和~$SR$（缩放+旋转）的整数指针~$g2sh$~和~$g2sr$；训练过程通过~Metropolis-Hastings (MH)~采样在状态空间~$S = \langle G, C \rangle$（高斯集合~+ 码本）上进行：参数更新~$q_\text{update}$~用~SGLD（Stochastic Gradient Langevin Dynamics，随机梯度朗之万动力学）~梯度步（98\%概率），分裂和合并操作~$q_\text{split}$/$q_\text{merge}$~动态调整码本大小（各~1\%概率），接受概率~$A(S \to S_\text{split}) = \min(1, e^{-\lambda_\text{SH}} q_\text{sm}(c'-c''))$~确保只有降低总体能量（重建损失~+ 码本大小惩罚）的转移才被接受。初始训练时每个高斯对应独立的码本向量（一对一映射），随着~MH~采样的进行，相似的向量被合并为共享码本条目，码本自然收缩。

\paragraph{SALVQ：场景自适应格向量量化} 针对均匀标量量化（USQ）在空间填充效率与场景统计适应性上的双重局限——其轴对齐的超立方体量化单元难以贴合实际数据分布，且对各场景属性的统计特性缺乏自适应能力——Compact 3DGS~及其后续方法使用均匀标量量化（USQ），其量化单元是轴对齐的超立方体，空间填充效率低。SALVQ（2025.09）~\cite{2509_13482}~用格向量量化（Lattice VQ）替代~USQ：不存储显式码本，而是通过可学习的格基矩阵~$B$~隐式生成码本——$B$~通过~SVD（奇异值分解，Singular Value Decomposition）参数化为~$B = U \Sigma V^T$，其中~$U, V$~为可学习正交变换（通过~PyTorch~正交参数化约束），$\Sigma$~为可学习对角矩阵，保证~$B$~始终可逆且训练稳定。量化时采用~Babai~舍入技术（BRT，一种在格中快速寻找近似最近点的算法）$\hat{q}_l(\mathbf{y}) = B \lfloor B^{-1} \mathbf{y} \rfloor$，在任意格基下实现快速近似最近格点搜索。关键差异在于：USQ~的~Voronoi~单元是标准超立方体，而~SALVQ~通过学习~$B$~将其旋转和拉伸为更贴合数据分布的多面体——当所有奇异值趋近~1~时退化为旋转超立方体，沿主方向捕捉相关性以降低量化误差。SALVQ 默认仅应用于锚点的隐层特征~$\mathbf{f}$（占主要比特预算），并支持可变速率控制：训练时学习增益向量~$\mathbf{g} = [g_1, \dots, g_M]$，推理时选择~$g_i$~缩放量阶~$q_s \to g_i q_s$~即可调整码率，无需重新训练。

\subsubsection{熵编码与上下文建模}

本路线介绍利用高斯属性之间的空间相关性建立上下文模型，结合熵编码实现无损压缩的方法。

量化将连续参数离散化为有限精度的整数，但整数序列中仍存在统计冗余。熵编码（如算术编码）利用符号的概率分布进行无损压缩——出现频率高的符号用更少的比特编码。关键在于如何准确估计每个量化参数的概率分布，这需要建立高斯属性之间的上下文模型。

\paragraph{哈希网格驱动的超先验上下文建模}

HAC++（2025.01）~\cite{2501_12255}~将~Scaffold-GS~的锚点属性从``独立编码''转变为``条件编码''——不再逐锚点独立存储量化值，而是先利用结构化哈希网格提取空间上下文作为\textbf{超先验}（hyperprior），再以该先验为条件估计每个锚点属性的概率分布，从而大幅降低熵编码所需的比特数。

\paragraph{Baseline 与改了什么}

读者已从第~\ref{sec:11_5}~节了解 Scaffold-GS 的锚点表示：每个锚点用位置~$\mathbf{x}^a \in \mathbb{R}^3$~和属性集合~$A = \{\mathbf{f}^a \in \mathbb{R}^{D^a},\, \mathbf{s} \in \mathbb{R}^6,\, \mathbf{o} \in \mathbb{R}^{3K}\}$~描述相邻高斯的特征、缩放和偏移。然而，Scaffold-GS 对每个锚点的属性是独立存储的——锚点之间虽然空间上邻近且属性高度相关（相邻锚点往往有相似的颜色倾向、尺寸和方向），这些空间冗余完全未被利用。HAC++ 的核心改动是引入一个结构化哈希网格作为超先验，将独立概率~$p(A)$~替换为条件概率~$p(A \mid \mathbf{h}^a)$，其中~$\mathbf{h}^a$~是从哈希网格中提取的空间上下文特征。直觉上，这就像``先看邻居，再猜自己''——知道周围锚点的属性分布后，当前锚点的不确定性大大降低，编码所需比特随之减少。

\paragraph{哈希网格与特征提取}

\textbf{哈希网格}（hash grid）是一种多分辨率~3D~特征场，将空间划分为多层不同分辨率的网格，每层在每个格点上存储一个低维特征向量。HAC++ 构建了一个混合~3D/2D~结构化二值哈希网格~$H$，共~$L=16$~层：其中~12~层~3D~嵌入（分辨率从~$16$~到~$512$），4~层~2D~嵌入（分辨率从~$128$~到~$1024$），每层特征维度~$D^h=4$。对于位于~$\mathbf{x}_i^a$~的锚点，通过对哈希网格进行三线性插值（3D层）或双线性插值（2D层），得到该锚点的哈希特征~$\mathbf{h}_i$：

\begin{equation}
\mathbf{h}_i = \text{Interp}(\mathbf{x}_i^a,\, H)
\end{equation}

多层特征拼接后形成一个紧凑的空间描述符。哈希网格本身被二值化为~$\{-1, +1\}$（使用~STE~直通估计器），以极大压缩其存储开销。

\textbf{设计动机}：为什么哈希网格能捕捉锚点间的空间相关性？核心观察是——邻近的高斯（以及对应的锚点）往往具有相似的颜色、尺寸和方向参数，这是一种~3D~空间中的局部一致性。哈希网格作为一个覆盖整个场景的结构化特征场，天然能通过插值操作捕捉这种局部一致性：空间上接近的锚点查询到相似的哈希特征，从而获得相似的条件概率估计。论文通过实验验证了哈希特征~$\mathbf{h}$~与锚点属性~$A$~之间存在大量互信息，这为条件编码提供了信息论基础。

\paragraph{条件概率模型}

有了哈希特征~$\mathbf{h}_i$~后，HAC++ 构建条件概率模型来估计每个锚点属性的量化值分布。整个编码框架的信息流可以形式化为：

\begin{equation}
p(A,\, \mathbf{x}^a,\, H) \approx p(A \mid \mathbf{h}^a)\, p(H)
\end{equation}

其中~$p(H)$~是哈希网格本身的先验（通过统计~$+1/-1$~出现频率编码），$p(A \mid \mathbf{h}^a)$~是条件概率——即``已知空间上下文后，锚点属性的分布''。条件概率越高（不确定性越低），熵编码所需比特越少。

具体而言，哈希特征~$\mathbf{h}_i$~被送入一个~3~层~MLP~上下文模型~$\text{MLP}_c$，为每个锚点的每种属性分别预测高斯分布的均值~$\boldsymbol{\mu}_i$~和标准差~$\boldsymbol{\sigma}_i$。对于量化后的属性值~$\hat{a}_i$，其概率通过高斯分布的累积分布函数计算：

\begin{equation}
p(\hat{a}_i \mid \mathbf{h}_i) = \Phi\!\left(\frac{\hat{a}_i + \Delta_i/2 - \mu_i}{\sigma_i}\right) - \Phi\!\left(\frac{\hat{a}_i - \Delta_i/2 - \mu_i}{\sigma_i}\right)
\end{equation}

其中~$\Delta_i$~是自适应量化步长（见下文），$\Phi$~是标准正态分布的累积分布函数。这个公式的直觉含义是：量化值~$\hat{a}_i$~的概率等于``在该高斯分布下，落在~$[\hat{a}_i - \Delta_i/2,\, \hat{a}_i + \Delta_i/2]$~区间内的概率质量''。概率越集中（$\sigma_i$~越小），区间内的概率质量越大，编码~$\hat{a}_i$~所需比特越少。

此外，HAC++ 还引入了\textbf{锚点内上下文模型}（intra-anchor context）作为辅助信息。对于锚点特征~$\mathbf{f}^a$~（维度最高、内部冗余最大），将其分割为~$N^s=5$~个块，逐块用~MLP~基于已解码块预测当前块的分布参数。超先验和锚点内上下文通过\textbf{高斯混合模型}（GMM，Gaussian Mixture Model）融合，自适应地平衡两个信息源的权重。

\paragraph{自适应量化与自适应掩码}

HAC++ 的自适应量化模块（AQM）解决了``统一量化步长不适用于所有属性''的问题。不同属性有不同的数值范围：锚点特征~$\mathbf{f}^a$~适用步长~$Q_0=1$，缩放~$\mathbf{s}$~需要精细步长~$Q_0=0.001$，偏移~$\mathbf{o}$~适用~$Q_0=0.2$。进一步地，AQM 通过~MLP~从哈希特征预测一个自适应修正项：

\begin{equation}
\Delta_i = Q_0 \cdot (1 + \tanh(r_i)), \quad r_i = \text{MLP}_q(\mathbf{h}_i)
\end{equation}

这使得量化步长被限制在~$(0,\, 2Q_0)$~范围内，既保留了足够的原始信息以维持渲染保真度，又通过匹配局部属性特性来降低熵。

自适应掩码策略通过\textbf{可学习的二值掩码}剪枝无效高斯和锚点。每个高斯偏移拥有一个由可学习特征经~Sigmoid~激活后二值化的掩码~$m_i \in \{0,1\}^K$（$K=10$~为每个锚点的高斯数）。如果一个锚点的所有偏移都被掩码掉，该锚点本身也被移除。HAC++ 的关键设计是将掩码信息直接融入率失真优化中的比特消耗计算：

\begin{equation}
\mathcal{L}_{\text{entropy}} = \sum_{i=1}^{N} m_i^a \!\left[ \sum_{j} -\log_2 p(\hat{f}_{i,j}) + \sum_{k} m_{i,k} \sum_{j} -\log_2 p(\hat{f}_{i,3k+j}) \right]
\end{equation}

其中~$m_i^a$~是锚点级掩码。这一设计使掩码比例随率失真权衡参数~$\lambda$~自动调节，无需额外正则化损失，在不同码率点自动找到最优剪枝比例。

\paragraph{训练策略与压缩效果}

训练分三阶段：前~3k~迭代训练纯~Scaffold-GS~以确保锚点生成过程稳定；3k--10k~迭代对锚点属性添加均匀噪声以适应量化扰动；10k~之后完整引入~HAC++~框架联合训练上下文模型。为避免显存溢出，每次迭代随机采样~5\%~的锚点进行熵训练。

压缩效果的存储分解：码流包含五个部分——锚点属性~$A$~（经算术编码，占主导）、二值哈希网格~$H$~（经算术编码）、偏移掩码（经算术编码）、锚点位置（16~位量化后用~GPCC~压缩）和~MLP~参数（32~位直接存储）。最终在~Synthetic-NeRF、Mip-NeRF 360、BungeeNeRF、Deep Blending、Tanks and Temples~等五个数据集上，HAC++ 实现了超过~100$\times$~压缩（相对原始~3DGS）和超过~20$\times$~压缩（相对~Scaffold-GS），同时渲染质量（PSNR）甚至有所提升——主要得益于熵损失的正则化效果和锚点特征维度的增大（$D^a=50$）。

\paragraph{锚点级自回归上下文模型}

\paragraph{从超先验到自回归：信息流方向的根本改变}

回顾~HAC++ 的设计：它通过哈希网格提取一个\textbf{固定的}空间先验~$\mathbf{h}^a$，然后以此为条件独立编码每个锚点的属性——信息流是``单向''的，从哈希网格流向锚点，锚点之间不存在相互依赖。这意味着~HAC++ 虽然利用了空间位置上的统计冗余，但无法利用锚点属性本身的值之间的相关性——即使两个相邻锚点的特征向量高度相似，HAC++ 也只能通过空间位置间接推断，而不能直接看到邻居的具体属性值。

ContextGS（2024.05）~\cite{2405_20721}~的核心思想是将信息流从``单向超先验''升级为\textbf{双向自回归}：已编码锚点的具体属性值被直接用作未编码锚点的上下文。直觉上，如果~HAC++ 是``先看地图，再猜房子''，那么~ContextGS 是``先看邻居的房子，再猜自己的房子''——后者显然提供了更精确的预测信息，因为邻居的实际特征值比一张抽象的空间地图包含了更丰富的场景细节。

\paragraph{八叉树分层与锚点前向}

要实现自回归编码，首先要解决``谁先编码''的问题。ContextGS 采用\textbf{八叉树分层}（octree partitioning）策略，将~$N$~个锚点按``自底向上''的方式划分为~$K$~个不相交的层级~$\{\Lambda^0, \Lambda^1, \ldots, \Lambda^{K-1}\}$。具体而言，以最细层~$\Lambda^0$~的体素大小为基础（与~Scaffold-GS~的原始体素大小相同），逐层按比例放大体素分辨率~$\kappa_i$，在每层体素内选择一个代表锚点构成粗层。论文对所有实验统一设~$K=3$~层，并通过二分搜索自动确定各层的体素缩放参数，使相邻层锚点数量的目标比值为~$0.2$。

编码顺序从最粗层~$\Lambda^{K-1}$~到最细层~$\Lambda^0$~逐层进行。在编码第~$k$~层的锚点~$v_i^k = (\mathbf{x}_i^k, \mathbf{f}_i^k, \mathbf{s}_i^k, \mathbf{o}_i^k)$~时，所有更粗层~$\Lambda^{k+1}$~中与该锚点位于同一体素的相邻锚点~$v_j^{k+1}$~已经被解码，其特征~$\mathbf{f}_j^{k+1}$~和缩放~$\mathbf{s}_j^{k+1}$~可以直接作为上下文信息。这一设计的关键优势在于\textbf{锚点前向}（anchor forward）机制：粗层锚点在解码后不仅用作上下文，还被直接复用到最终的细层渲染中，避免了跨层重复存储。消融实验表明，移除锚点前向会导致严重的存储冗余，效果甚至不如不分层——这证明了粗层锚点的``一鱼两吃''（既提供上下文又参与渲染）是压缩效率的重要来源。

\paragraph{层级条件概率模型}

对于第~$k$~层的锚点~$v_i^k$，其特征~$\mathbf{f}_i^k$~的条件概率分布由一个层级专属~MLP~$F_\theta^k$~预测：

\begin{equation}
p_{\mathbf{f}}^k(\mathbf{f}_i^k \mid \boldsymbol{\pi}_i^k) = \left(\mathcal{N}(\boldsymbol{\mu}_i^k, \boldsymbol{\sigma}_i^k) * \mathcal{U}(-\tfrac{1}{2}\boldsymbol{\Delta}_i^k, \tfrac{1}{2}\boldsymbol{\Delta}_i^k)\right)(\hat{\mathbf{f}}_i^k), \quad (\boldsymbol{\mu}_i^k, \boldsymbol{\sigma}_i^k, \boldsymbol{\Delta}_i^k) = F_\theta^k(\boldsymbol{\pi}_i^k)
\end{equation}

其中~$\boldsymbol{\pi}_i^k$~是上下文特征。对于最粗层~$\Lambda^{K-1}$~（没有已编码锚点），上下文仅为锚点位置~$\mathbf{x}_i^{K-1}$~；对于更细层~$k < K-1$，上下文扩展为~$\boldsymbol{\pi}_i^k = [\mathbf{x}_i^k;\, \mathbf{f}_j^{k+1};\, \mathbf{s}_j^{k+1}]$，即锚点位置拼接已解码相邻锚点的特征和缩放。与~HAC++ 仅使用位置隐式推断空间相关性不同，ContextGS 直接将邻居锚点的\textbf{实际属性值}纳入条件概率模型，信息量显著更丰富。

\paragraph{超先验补充与整体压缩效果}

对于最粗层锚点（无已编码邻居可参考），仅依赖位置信息的建模能力有限。ContextGS 为此为每个锚点引入一个低维可学习超先验特征~$\mathbf{z}_i \in \mathbb{R}^{h_c}$（$h_c=4$，仅为锚点特征维度~$D^a=50$~的~$1/4$），使用全因子化密度模型（2018）~\cite{balle2018variational}~建模。最终上下文为~$\boldsymbol{\pi}_i^k = [\mathbf{x}_i^k;\, \mathbf{f}_j^{k+1};\, \mathbf{s}_j^{k+1};\, \mathbf{z}_i]$。值得注意的是，ContextGS 的超先验是\textbf{锚点级别}的低维特征（$h_c=4$），而~HAC++ 的哈希网格是一个\textbf{场景级别}的结构化特征场（16~层、每层~4~维）——ContextGS 以远低于~HAC++ 的超先验存储开销，通过自回归上下文弥补了建模能力的差距。

整体训练损失为~$\mathcal{L} = \mathcal{L}_\text{scaffold} + \lambda_e \mathcal{L}_\text{entropy} + \lambda_m \mathcal{L}_m$，其中~$\mathcal{L}_\text{entropy}$~同时包含超先验特征的编码开销和所有层级锚点属性的编码开销。消融实验显示，锚点级上下文模型可减少~21\%~存储，超先验特征额外减少~10\%，两者联合使用实现总计~26\%~的压缩增益。最终在~BungeeNeRF、Deep Blending、Mip-NeRF 360、Tanks and Temples~等数据集上，ContextGS 相对~Scaffold-GS~实现约~15$\times$~压缩，相对原始~3DGS~实现约~100$\times$~压缩，同时渲染质量与~Scaffold-GS~持平甚至略有提升。

一个关键的实际优势是解码效率。图像压缩中的自回归模型通常因逐像素串行预测而极其缓慢，但~ContextGS 的自回归是\textbf{逐层}而非逐锚点进行的——$K=3$~层仅引入~3~次循环迭代，且不同层级之间无锚点重叠，总处理锚点数不变，因此解码时间与~HAC++~等非自回归方法相当。

\paragraph{后续发展：上下文建模的进一步扩展}

ContextGS 证明了锚点级自回归上下文在~3DGS~压缩中的有效性，后续工作沿着多个方向扩展。CAT-3DGS（2025.03）~\cite{2503_00357}~引入\textbf{三平面}（triplane）结构捕捉高斯属性之间的跨维度相关性，将上下文建模从锚点间扩展到属性维度间。CAT-3DGS Pro（2025.03）~\cite{2503_12862}~用~PCA（Principal Component Analysis，主成分分析）~引导的向量-矩阵超先验替代三平面，在保持压缩率的同时将训练速度提升~3$\times$、解码速度提升~5$\times$。HEMGS（2024.11）~\cite{2411_18473}~提出混合熵模型，将可变速率预测器、超先验网络和自回归网络统一为单一框架，用同一个模型即可实现多速率压缩，免去了为每个码率点重新训练的代价。

\subsubsection{前馈压缩与通用编解码器}

本路线介绍训练通用编码器-解码器网络对任意 3DGS 场景进行单次前向传播压缩的方法。前述方法均需要逐场景优化，压缩一个场景通常需要数十分钟到数小时的训练。前馈压缩（feed-forward compression）旨在训练一个通用的编码器-解码器网络，对任意~3DGS~场景进行单次前向传播即可完成压缩。

\paragraph{前馈式高斯场景压缩}

FCGS（2024.10）~\cite{2410_08017}~代表了从"逐场景优化"到"通用前馈"的范式转变——它训练一个通用的编码器-解码器网络，对任意~3DGS~场景执行单次前向传播即可完成压缩，无需任何逐场景微调。可以把前馈压缩类比为一个"万能解压器"：就像通用的图像/视频编解码器（如~JPEG、H.264）只需训练一次就能压缩任意图片或视频一样，FCGS 只需训练一次就能压缩任意~3DGS~场景。

\paragraph{设计动机}

前述所有方法（包括~HAC++、ContextGS）都需要逐场景优化来适配压缩参数，压缩一个场景通常需要数十分钟到数小时的微调训练。这种范式在需要永久存储或服务端编码的场景下是可接受的，但对时间敏感的应用（如在线传输、移动端压缩）而言，速度成为瓶颈。一个自然的想法是：能否像图像编解码器那样，训练一个通用的压缩网络，让它在一次前向传播中完成压缩？这就是前馈压缩（feed-forward compression）的核心思路。

\paragraph{多路径熵模块（MEM）}

直接将所有高斯属性输入同一个自编码器会导致严重的质量下降。其原因在于~3DGS~中并非所有属性的敏感度相同：\textbf{几何属性}（不透明度、缩放、旋转，共~8~维）直接决定光栅化过程中的遮挡关系和深度排序——即使微小的偏差也会在渲染中被放大（例如一个高斯的不透明度解码偏大，就会错误地遮挡后方所有高斯）；而\textbf{颜色属性}（球谐系数，共~48~维）仅影响着色结果，对偏差的容忍度更高。

基于这一观察，FCGS 引入\textbf{多路径熵模块}（Multi-path Entropy Module, MEM），为每个高斯的颜色属性自适应地生成一个二值掩码~$m_i$，决定该属性走哪条路径：
\begin{itemize}
\item $m_i = 1$：属性通过自编码器~$g_a / g_s$~压缩到潜在空间再恢复，以消除冗余；
\item $m_i = 0$：属性不经自编码器，直接量化（跳过~MLP），最大程度保留原始信息。
\end{itemize}
形式化地，解码属性为：

\begin{equation}
\hat{\mathbf{x}}_i = g_s(Q(g_a(\boldsymbol{\psi}_i))) \cdot m_i + Q(\boldsymbol{\psi}_i) \cdot (1 - m_i), \quad m_i = \mathbb{I}[\mathrm{Sig}(f(\mathbf{x}_i^{\mathrm{gau}})) > \tau_m]
\end{equation}

其中~$Q$~为量化操作，$\mathrm{Sig}$~为~Sigmoid~函数，$\tau_m$~为阈值超参数。简言之，MEM 的作用是"挑选"哪些高斯的颜色属性足够鲁棒、可以安全地通过~MLP~压缩，哪些必须原样保留以保证渲染质量。而几何属性完全不经过~MLP，直接量化。掩码比率不通过额外的损失项控制，而是将掩码信息直接纳入比特消耗计算中，模型自适应学习最优掩码比率。

\paragraph{高斯间上下文模型：从稀疏到结构化}

在图像压缩中，数据天然排列在二维网格上，上下文模型的设计水到渠成。但~3DGS~中的高斯是稀疏且无序的，无法直接套用图像压缩的上下文建模思路。FCGS~的解决方案是：\textbf{利用~Morton~码序列化，将无序的高斯重新组织为结构化网格，从而建立高斯之间的长距离依赖关系}。Morton码（也称~Z-order~曲线）是一种空间填充曲线，它将三维坐标映射为一维序列，使得空间上相邻的点在一维序列中也保持相近。

具体而言，FCGS 将所有高斯的潜在表示~$\{\hat{\mathbf{y}}\}$~分为~$N^s = 4$~个批次，按顺序解码。在解码第~$n^s$~批时，利用已解码的批次~$Y_{[0, n^s-1]}$~构建三维和二维网格（多分辨率，三维网格分辨率为~$\{70, 80, 90\}$，二维网格分辨率为~$\{300, 400, 500\}$），通过三线性插值为待解码高斯提供上下文特征，再由一个小型网络预测其概率分布参数。这样，已解码高斯的空间信息被有效利用来辅助未解码高斯的熵估计。

\paragraph{高斯内上下文模型}

除了高斯之间的冗余，同一高斯内部的颜色属性也存在通道间冗余（如~RGB~三通道高度相关）。FCGS 将每个高斯的潜在表示沿通道维度分成多个块（chunk），逐块自回归解码：先解码前面的块，再基于已解码的块预测后续块的概率分布。对于~$m=1$~路径（通过自编码器），沿通道维度分为~$N^c=4$~块；对于~$m=0$~路径（直接量化），沿~RGB~轴分为~$N^c=3$~块以利用颜色分量间的冗余。

\paragraph{概率估计与训练}

上下文模型和超先验各自输出一组高斯分布参数，FCGS 用\textbf{高斯混合模型}（GMM）将三组分布参数加权组合，提供更精确的概率估计。比特消耗按信息论公式计算：

\begin{equation}
\mathrm{bit}_i = \sum_j^{D^y} (-\log_2 p(\hat{y}_{i,j})) + \sum_j^{D^z} (-\log_2 p(\hat{z}_{i,j}))
\end{equation}

即每个高斯的比特消耗等于属性潜在表示和超先验潜在表示的信息量之和。

训练数据方面，FCGS~基于~DL3DV~数据集约~7K~个多视角场景，训练生成~960P~分辨率的~3DGS~表示（耗时约~60~GPU~天），经质量筛选后保留~$6770$~个~3DGS~场景，其中~100~个用于测试。这个大规模训练集确保了压缩网络的泛化能力。

总体损失函数为：

\begin{equation}
\mathcal{L} = \mathcal{L}_{\mathrm{fidelity}}(\hat{I}, I) + \lambda \cdot \frac{1}{N^g \cdot 56} \cdot \mathcal{L}_{\mathrm{entropy}}
\end{equation}

其中~$\mathcal{L}_{\mathrm{fidelity}}$~同时包含~MSE~和~SSIM~指标，$N^g$~为高斯数量，$56$~为属性总维度（几何~8~维~$+$~颜色~48~维），$\lambda$~从~$1 \times 10^{-4}$~到~$16 \times 10^{-4}$~调节以实现可变速率。训练分三阶段：先仅用保真度损失训练~$m=1$~路径的自编码器；再联合训练上下文模型；最后端到端训练整个模型。

\paragraph{实验结果}

推理时，FCGS~在单张~GPU~上约~1~秒即可压缩~$100K$~个高斯，而逐场景优化方法通常需要数十分钟到数小时。FCGS 实现了超过~$20\times$~压缩率（相对原始~3DGS），且渲染质量可媲美大多数逐场景优化方法——尽管没有逐场景适配，这在"不公平"的比较设定下仍然表现出色。值得注意的是，FCGS 解码后的~3DGS~渲染速度与压缩前一致（如~Mip-NeRF 360~数据集上平均~FPS~从~102~微降到~91），因为~FCGS~不改变高斯的数量或结构。当与剪枝方法（Mini-Splatting、Trimming）结合时，FCGS~甚至能实现~$100\times$~压缩。此外，FCGS~对前馈~3DGS~模型（如~LGM、MVSplat）生成的场景也表现出零样本泛化能力，无需额外微调即可实现~$5\times$--$15\times$~压缩。

后续工作~FCGS~的扩展（2025.12）~\cite{2512_00877}~进一步引入长上下文建模，通过大规模上下文结构和细粒度空间-通道自回归熵模型捕捉数千个高斯之间的远程依赖，实现了~20$\times$~压缩率。

\subsubsection{后训练压缩管线}

本路线介绍不修改训练过程、直接对已训练好的3DGS模型进行端到端压缩的管线方法。

与需要修改训练过程的方法不同，后训练压缩（post-training compression）直接对已训练好的~3DGS~模型进行压缩，无需重新训练。

\paragraph{尺寸感知的后训练编解码器}

需要修改训练过程的压缩方法（如熵模型引导的训练、前馈压缩）虽然有效，但要么要求用户重新训练整个场景（耗费~30--60~分钟），要么需要在大规模~3DGS~数据集上预训练通用压缩网络。对于已经训练好~3DGS~模型的用户而言，这些方法的门槛仍然偏高。后训练压缩（post-training compression）正是为解决这一痛点而设计的——它直接对已训练好的~3DGS~模型进行压缩处理，不需要访问原始训练图像，也不需要重新训练，用户只需提供一个预训练模型和目标存储预算即可。

MesonGS++（2026.04）~\cite{2604_26799}~是目前后训练压缩路线中最完整的管线，依次包含\textbf{剪枝、属性变换、量化和熵编码}四个阶段。其核心创新在于\textbf{尺寸感知的超参数搜索}：用户只需指定目标存储预算，MesonGS++ 自动搜索最优的高斯保留率（$\tau$）和分组量化位宽（$\mathbf{b}$），无需人工调参。

\textbf{压缩管线的四个阶段}各有关键设计。在\emph{剪枝阶段}，MesonGS++ 设计了联合重要性指标~$I_g = I_d \cdot I_i$，将视角依赖得分（考虑渲染时高斯的排序与遮挡关系）与视角无关得分（基于归一化体积）相乘，比仅用不透明度或梯度作为剪枝准则更准确地衡量每个高斯对最终渲染的贡献。在\emph{属性变换阶段}，MesonGS++ 先将旋转四元数（4~个数）替换为~Euler~角（3~个数），节省一个数的存储且不会像协方差矩阵替换那样在量化后导致~50\%~的矩阵失去正定性；然后对重要属性（不透明度、缩放、旋转、零阶~SH~系数）施加\textbf{区域自适应分层变换}（Region Adaptive Hierarchical Transform, RAHT），利用八叉树结构将属性转换为低熵的~AC~系数以利于后续压缩；对占据存储~85.7\%~的高阶~SH~系数则采用\textbf{向量量化}（vector quantization）配合批量聚类策略，并采用预算优先分配策略——先估算核心组件的压缩大小，将剩余预算用于保留重要性最高的一部分高斯的原始~SH~系数。在\emph{量化阶段}，MesonGS++ 对~RAHT~变换后的~AC~系数进行分组量化（group-wise quantization），每个量化组~$c$~按如下方式量化：

\begin{equation}
c_q = \text{clamp}\!\left(\left\lfloor \frac{c}{S_c} \right\rfloor + Z_c,\; 0,\; 2^b - 1\right), \quad S_c = \frac{\max(c) - \min(c)}{2^b}, \quad Z_c = -\left\lfloor \frac{\min(c)}{S_c} \right\rfloor
\end{equation}

其中~$b$~为位宽，$S_c$~为缩放因子，$Z_c$~为零点偏移——即先将属性值缩放到~$[0, 2^b{-}1]$~的整数范围，超出范围的值被截断。不同属性组可以使用不同的位宽——例如实验发现不透明度和旋转对量化更敏感，通常被分配更高位宽，而~RGB~和缩放则使用较低位宽。最后在\emph{熵编码阶段}，八叉树坐标用~G-PCC~压缩，其余整数元素通过算术编码（torchac）和~LZ77~进行压缩打包。

直觉地说，压缩管线的超参数（保留率~$\tau$、各组的量化位宽~$\mathbf{b}$、码本大小）之间存在复杂的耦合关系——降低保留率意味着保留更多高斯但留给每个高斯的存储预算更少，需要更低位的量化来补偿；反之，提高保留率可以给每个高斯分配更多位宽，但高斯总数减少可能损失细节。MesonGS++ 将这一搜索问题形式化为\textbf{混合整数非线性规划}（Mixed Integer Nonlinear Programming, MINLP），目标是在满足存储预算约束~$S(\tau, \mathbf{b}) \leq \text{Budget}$~的条件下最小化质量损失~$Q(\tau, \mathbf{b})$。由于直接求解~MINLP~的搜索空间极大（以~$C{=}10$~个属性通道、$B{=}60$~个分组、$Q{=}16$~种位宽选项为例，搜索空间高达~$2.96 \times 10^{722}$~种组合），MesonGS++ 通过两个关键观察将问题解耦为两步：（1）在固定~$\tau$~下，压缩大小与位宽设置近似呈线性关系；（2）$\tau$~是唯一的连续变量，其余均为整数。因此，算法先对~$\tau$~进行离散采样，对每个固定的~$\tau$~值，再用~\textbf{0-1~整数线性规划}（0-1 ILP）高效求解最优位宽配置。这里~0-1 ILP~的含义是：将每个属性组的位宽选择表示为一个~one-hot~二值向量，用线性函数估计压缩大小，从而将原始的非线性问题转化为标准线性规划问题。

为加速~0-1 ILP~的求解过程，MesonGS++ 提出了三项关键技巧。第一，\textbf{线性尺寸估计器}：压缩文件的估计大小建模为~$S(\mathbf{b}) = \sum_{i,j} \omega_{ij} b_{ij} + S_\text{meta}$，其中~$\omega_{ij}$~是各量化组的大小权重，$S_\text{meta}$~是元数据和坐标的精确大小（可通过实际存储快速获取），该线性形式直接适配~ILP~求解器。第二，\textbf{量化损失替代~PSNR}：计算~PSNR~需要遍历整个训练集渲染，每次耗时超过~10~秒；而量化损失~$\mathcal{L}(i,j,b) = \| \hat{A}_{i,j}^b - A_{i,j} \|$（原始属性与量化恢复属性之间的距离）可以预先对每个组独立计算，配合~CUDA~并行量化核函数，每次仅需~0.49~秒。第三，\textbf{层级求解}：先在通道级别搜索粗粒度位宽，再根据通道预算在分组级别搜索细粒度位宽，将~ILP~的求解时间控制在~58.52~秒以内。整个超参数搜索（含校准和~ILP~求解）总耗时约~104.51~秒。

实验表明，MesonGS++ 在多个基准数据集上一致实现了优于先前后训练方法的压缩率-质量权衡。具体而言，MesonGS++ 实现了超过~34$\times$~的压缩率；在~Stump~场景上，20$\times$~压缩后的~PSNR~甚至超过了原始~3DGS（同时~Garden~场景以~15$\times$~压缩也超过了原始~3DGS）。MesonGS++ 的尺寸控制精度也很高——在~ScaffoldGS~和~4DGS~变体上，最终模型大小与指定预算的偏差仅分别为~0.8\%~和~0.5\%，证明了管线对~3DGS~变体的良好迁移性。

与~MesonGS++ 依赖启发式重要性指标（如视角依赖得分与体积的乘积~$I_g = I_d \cdot I_i$）来评估高斯重要性不同，POTR（2026.01）~\cite{2601_14821}~的核心改动是将剪枝评估从``间接推测''升级为``直接测量''——修改~3DGS~光栅化器的正向渲染过程，利用已有数据（每个像素的~$\alpha$-合成中间结果~$\mathbf{P}_{i_k}$~和透射率~$T_k$），为所有高斯\textbf{同步}计算其移除对渲染结果的精确影响。具体而言，对于像素~$p$~中排序为~$i_k$~的第~$k$~个高斯，其移除后的颜色变化（剪枝差异）为~$\Delta \mathbf{c}_k(p) = \mathbf{c}_k(p) - \mathbf{c}(p) = \mathbf{o}_{i_k} \Delta \mathbf{P}_{i_k+1} + o_{i_k} \mathbf{P}_K$，其中~$o_{i_k}$~为不透明度，$\mathbf{P}_{i_k+1}$~为移除后所有后续高斯的合成颜色。由于后续高斯的透射率被重新归一化为~$1/(1 - o_{i_k})$，这一公式精确捕捉了``移除一个红色高斯后露出背后蓝色高斯''的视觉效果——这是启发式指标无法做到的。由此推导出移除第~$k$~个高斯对~MSE（Mean Squared Error，均方误差）~的影响~$\Delta \text{MSE}_k = \frac{1}{3NP} \sum_{c,p} \Delta \text{SE}_k(c,p)$，作为剪枝的直接判据。

POTR 采用迭代剪枝策略：每次迭代按~$m(\Delta \text{MSE}_k / \text{MSE}_\text{MAX})$~排序添加候选高斯（优先移除~MSE~变化小或为负的高斯），直到达到基于重要性的预算阈值，然后重新计算所有~$\Delta \text{MSE}_k$~值以修正累积近似误差。此外，POTR 还设计了\textbf{球谐能量压缩}——利用岭回归找到一组新的球谐系数，使在训练视角子集内的颜色保持不变但在未采样方向上自由变化，配合~YCoCg~亮度-色度颜色模型（色度通道正则化强度为亮度的~3~倍），将球谐系数的熵降低~50\%~以上。整个编解码器无需微调即可实现超过~50\%~的高斯剪枝率；加入解码-微调-重编码循环（POTR-FT）后，在~Mip-NeRF 360~数据集上的~RD~性能接近或超过需要重新训练的方案，同时推理速度提升~1.5--2$\times$。

\subsubsection{基于图像/视频编码器的压缩}

本路线介绍将3DGS参数映射为2D图像或视频格式后利用成熟编解码标准进行压缩的方法。

将~3DGS~参数映射为~2D~图像或视频，然后利用成熟的图像/视频编码标准（如~HEVC、AV1）进行压缩，是一条充分利用现有编解码器生态的实用路线。

\paragraph{结构化图像编码}

GSICO（2026.01）~\cite{2601_14510}~的核心思想是将~3D~高斯基元的全部参数（位置、协方差、球谐系数、不透明度）视为一幅``展平''的~2D~图像，从而可以直接调用成熟的标准图像编解码器进行压缩。

直觉地说，3DGS~的数十万个高斯基元散布在三维空间中，位置和属性都是不规则的——它们之间没有天然的行列结构。而图像编解码器（如~JPEG、HEVC）依赖的是像素之间的\textbf{空间连贯性}（spatial coherence）：相邻像素值相近，可以用差分编码高效压缩。因此，核心挑战在于如何将不规则的~3D~点云映射为具有空间连贯性的~2D~图像。打个比方，这就像把一团杂乱无章的毛线重新缠绕到线轴上——关键是要让空间上挨得近的高斯，在~2D~图像中也排列在相邻的像素位置，这样编解码器才能``看到''平滑的数值过渡，而非随机噪声。

GSICO 提出了一种专门的映射算法来解决这个问题。其核心是利用~PLAS（Priority-based Local-Aware Sorting）排序：先根据高斯基元在~3D~空间中的邻近关系进行排序，使空间上邻近的高斯在排序后也相邻；然后将有序的参数序列按行优先的方式填充到~2D~图像中。这样一来，图像中相邻像素对应的高斯参数在~3D~空间中也彼此接近，数值自然相似——标准图像编解码器正是为这种局部平滑性而设计的。

在压缩效果方面，由于~2D~映射后的参数图像具有良好的空间相关性，标准编解码器能够高效地利用这种冗余。在~Tanks and Temples、Deep Blending~和~Mip-NeRF 360~数据集上，GSICO 实现了平均~20.2$\times$~压缩率，验证了``借用成熟编解码器生态''这一技术路线的可行性。

\paragraph{轻量级视频编码压缩}

LGSCV（2025.12）~\cite{2512_11186}~沿着~GSICO``借用成熟编解码器生态''的思路继续推进，但将后端从图像编解码器升级为视频编解码器（如~HEVC），以利用帧间预测进一步挖掘多帧~3DGS~参数之间的冗余。然而，直接套用~GSICO~的方案会遇到一个实际瓶颈：GSICO~所依赖的~PLAS~排序虽然能产生高质量的~2D~映射，但其计算代价相当昂贵。打个比方，PLAS~排序就像是用``显微镜''逐一检查每对高斯之间的关系来确定最优排列——结果虽然精确，但耗时也成正比。LGSCV~的核心贡献正是用一个更轻量的替代方案来绕过这个瓶颈。

\textbf{两阶段 Morton 扫描替代 PLAS 排序。}
与~GSICO~的~PLAS~排序不同，LGSCV~提出了两阶段~Morton~扫描来建立~3D~高斯到~2D~图像的映射。Morton~码（也称~Z-order curve）的思想很简单：将三维坐标的二进制表示逐位交错（interleave），生成一个单一整数，使得空间上接近的点在排序后索引也接近。第一阶段，对高斯基元的三维空间坐标计算~3D Morton~码并排序——这一步复杂度为~$O(N \log N)$，远低于~PLAS~的多轮局部搜索。第二阶段，给定已排序的基元序列，不再像~GSICO~那样用行优先方式填充~2D~图，而是用~2D Morton~扫描将基元映射为分块的~2D~图。直觉地说，2D Morton~扫描的填充路径呈~Z~字形，使得生成的~2D~图天然呈现\textbf{块状结构}（block structure）：每个块内的基元在空间上接近、属性值相似。这种块状结构恰好与视频编码器的基本处理单元（Coding Unit, CU）对齐——HEVC~的自适应块划分正是以方块为基础的，因此块状~2D~图能让编码器更高效地利用局部冗余。

\textbf{PCA 降维球谐系数。}
原始~3DGS~每个高斯基元有~59~个特征通道，其中~45~个是球谐交流分量（SH AC）。若直接将这~45~个通道当作~45~帧灰度图送入视频编码器，不仅编解码耗时增加，而且无法充分利用不同通道之间的上下文关联。受基于优化的~3DGS~压缩方法启发（球谐系数高度冗余，可用低秩表示重建），LGSCV~对~45~维~SH AC~特征矩阵一次性施加~PCA，仅保留前~$k$~个主成分进行压缩，其余成分直接裁剪。作者对比了三种降维策略——对全部~45~通道统一~PCA、按~R/G/B~三个颜色空间分别~PCA、按球谐阶数直接裁剪高阶项——发现统一~PCA~效果最优，因为它同时捕获了通道内和通道间的相关性。

\textbf{MiniPLAS 块内微调。}
两阶段~Morton~扫描仅基于空间坐标排序，未考虑高斯的其他属性（颜色、尺度等），因此块内可能仍有局部不连贯之处。直接在整幅~2D~图上运行~PLAS~则会破坏好不容易建立的块状结构，且速度仍然很慢。为此，LGSCV~提出了~MiniPLAS：在固定大小的块内（如~$8 \times 8$，与视频编码器的~CU~或子~CU~尺寸一致）进行像素排列微调。由于块很小，每个块的排列空间有限——以~4~像素为一个单元，每单元仅有~$4! = 24$~种排列——因此总体复杂度约为~$O(6M^2)$（$M$~为图像分辨率），远低于全图~PLAS。更重要的是，MiniPLAS~的块大小与编码器~CU~尺寸对齐，可以引导编码器快速确定最优块划分，减少率失真优化（RDO）的搜索空间，从而将编码时间降低约~50\%的同时还提升了编码效率。

\textbf{端到端压缩流程。}
在具体实现中，LGSCV~对坐标使用~20~位深度（bit depth），其他属性使用~10~位深度，总共生成~$7 + N_k$~幅三通道图像送入视频编码器。其中坐标图像采用无损模式压缩（因为~3DGS~质量对坐标精度高度敏感），其余属性采用有损压缩。值得注意的是，LGSCV~直接将图像数据视为~YUV444~格式而非~RGB——因为除颜色分量外，其他属性之间不存在~RGB~式的通道关联，而~RGB~与~YUV~之间的颜色空间转换是有损的，对坐标图会引入明显的质量下降。

在压缩效果方面，LGSCV~在~MPEG~数据集上比当时的~SOTA~方法提升约~20\%~率失真性能，同时将~2D~图生成时间从~PLAS~方案的大量计算降至约~1~秒，充分验证了``轻量级排序 + 成熟视频编码器''这一技术路线的实用性。

\begin{takeaway}
编码与压缩将高斯属性映射为2D图像后利用成熟视频编码器（VVC、HEVC）进行压缩，核心思路是"3D到2D的保序映射"。MPEG 标准化活动已建立统一的评测框架，LGSCV 的两阶段 Morton 扫描加 MiniPLAS 块内微调实现了约 20\% 率失真性能提升。关键发现是坐标精度对质量高度敏感，需无损压缩。
\end{takeaway}


%=====================================================================

\subsection{近期方法脉络：Mobile Gaussian rendering and radiance compression、渐进式传输与LOD、流式传输与动态压缩}

这一组近期工作延续 3DGS、progressive 3DGS、dynamic 3DGS，主要把改动集中在压缩、增量改进等环节。下面只保留相对既有路线的变化，便于和前文方法直接对照。

\paragraph{Monte Carlo Energy Aggregation for Mobile 3D Gaussian Splatting}
Monte Carlo Energy Aggregation for Mobile 3D Gaussian Splatting（2026.06）~\cite{2606_30017} 延续 3DGS 在“Mobile Gaussian rendering and radiance compression”上的基本设定，主要改变是：Monte Carlo Specular Energy Aggregator 与紧凑 latent 代替高阶 SH 开销。3DGS 的显式高斯表示和可微渲染能力保持为基础。
其设计动机是：针对 high-order Spherical Harmonics 带来的显著 overhead，提出面向移动端的实时高保真 GS 渲染方法。
具体而言，采样三阶 radiance residuals，将视觉显著的 specular energy 聚合进紧凑 latent space，从而减少 SH 参数和推理负担，同时保留高光细节。
适用边界是：复杂 view-dependent effects 的压缩可能仍有损失；移动端硬件差异会影响实测收益。

\paragraph{EvoGS}
EvoGS（2026.06）~\cite{2606_07179} 延续 progressive 3DGS 在“渐进式传输与LOD”上的基本设定，主要改变是：EvoGS：Evolution Tree连续分层表示，使用wavelet-inspired父子细化让子节点纠正祖先误差。3DGS 的显式高斯表示和可微渲染能力保持为基础。
其设计动机是：提出首个continuous-layering 3DGS表示，避免离散层级的误差累积、splat冗余和质量跳变。
具体而言，EvoGS把细节作为父节点到子节点的显式结构化演化信号，细层可以修正粗层误差，并形成更稀疏、可压缩的层间残差。
适用边界是：演化树构建和优化增加复杂度；动态场景和实时编辑下的层级维护仍需验证。

\paragraph{GS-NFS}
GS-NFS（2026.06）~\cite{2606_05650} 可作为“流式传输与动态压缩”中的新锚点候选。相较于 dynamic 3DGS，主要改变是：GS-NFS：GPU加速的动态3DGS压缩/解压，重构RAHT类算法以支持全帧率编码解码。3DGS 的显式高斯表示和可微渲染能力保持为基础。
其设计动机是：提出适合GPU加速的动态3DGS压缩管线，使带宽自适应流式传输从离线压缩走向实时编码解码。
具体而言，方法针对每帧高斯的位置、尺度、旋转、不透明度和颜色等属性进行可加速变换编码，重构RAHT算法以匹配GPU并行执行。
这一锚点判断来自当前分类中的机制关系：候选锚点：压缩，且相关性 0.90；后续方法可围绕 对逐帧高斯属性进行带宽自适应压缩传输 继续比较和扩展。
适用边界是：传输质量依赖码率控制；快速动态和复杂属性可能影响压缩稳定性。
\section{训练加速与高效表示}
\label{sec:11_3}

\subsection{问题背景}

本节讨论三条技术路线：（1）频率调制与粗到细训练；（2）内存有界的训练框架；（3）训练流程优化。

前两节讨论的是"训练后再压缩"的事后策略。然而，3DGS~训练过程本身的效率问题同样值得关注。原始~3DGS~的训练流程存在三个瓶颈：（1）\textbf{内存瓶颈}，ADC~不断添加高斯，峰值内存可达最终模型的~2--3~倍，在边缘设备上训练甚至无法完成；（2）\textbf{计算瓶颈}，训练中~30K~次迭代的反向传播占用了绝大部分时间，但其中大量迭代的梯度贡献微乎其微；（3）\textbf{冗余瓶颈}，训练结束时产生的大量高斯中，相当一部分是冗余的，事后剪枝不仅浪费了训练它们的计算资源，也可能损失已学习的场景信息。

训练加速与高效表示的目标是从源头控制这些瓶颈，在训练阶段就约束高斯增长、优化训练流程、减少冗余计算，使最终模型天生紧凑高效。

\subsubsection{频率调制与粗到细训练}

本节介绍利用频率域分析指导~3DGS~训练的方法，核心思想是在低频结构稳定后再逐步引入高频细节，以~Opti3DGS~为代表方法。

\paragraph{频率调制的粗到细优化}

Opti3DGS（2025.03）~\cite{2503_14475}~的核心策略是\textbf{频率调制}（frequency modulation）——通过控制训练图像的频率内容来引导~3DGS~的优化过程，实现从粗到细的渐进式训练。

\textbf{Baseline 问题。}在原始~3DGS~的训练流程中，自适应密度控制（ADC）机制根据渲染误差持续分裂和克隆高斯。由于训练图像从一开始就包含全部频率信息（从大尺度结构到精细纹理），优化器会同时尝试拟合所有尺度的信号。高频纹理（如树叶、栅栏）触发大量微小高斯的生成，低频区域（如墙面、天空）也可能因~SfM~初始点过密而产生冗余高斯。最终导致高斯数量爆炸，峰值可达所需数量的~2--3~倍。

\textbf{频率调制与课程学习。}Opti3DGS 的核心思想是：不改变~3DGS~本身，而是改变它``看到''的训练图像。具体而言，对训练图像施加低通滤波，早期只提供低频版本（只保留场景的大尺度结构），后期逐步放宽滤波范围、引入更高频的细节。这里涉及两个关键概念：\emph{频率调制}指对图像频谱的选择性截断与恢复；\emph{课程学习}（curriculum learning）指按由易到难的顺序组织训练信号——先学轮廓，再学纹理。

直觉上，这类似于画家的工作方式：画家从不直接在空白画布上逐像素描绘细节，而是先用粗线条勾勒整体构图（低频），再逐步添加色彩层次（中频），最后才处理光影纹理（高频）。这种从粗到细的工作流让每一阶段的"笔触"（对应高斯基元）都服务于明确的结构目标，避免了盲目堆砌。

\textbf{实现原理。}在具体实现上，Opti3DGS~设计了一个五级渐进管道（progressive growing pipeline）。给定原始训练图像集合~$I$，首先通过~SfM~获得初始稀疏点云作为~$G_0$：

\begin{equation}
    G_0 = \mathrm{SfM}(I)
\end{equation}

然后，对每一级~$n \in \{1, 2, 3, 4, 5\}$，使用滤波函数~$B$~对训练图像施加低通滤波，得到滤波后的图像集~$I_n$，并将上一级的高斯模型作为初始化继续优化：

\begin{equation}
    I_n = B(I, K_n, P_n), \quad G_n = \mathrm{3DGS}(G_{n-1},\, I_n)
\end{equation}

其中~$K_n$~为滤波核大小，$P_n$~为对应的填充参数。每一级图像逐渐变得更加锐利，允许更多高频信息通过。到最后一级~$n{=}5$~时，直接使用原始训练图像~$I$~进行精细优化。换言之，前四级逐步``教会''~3DGS~理解场景的宏观结构，最后一级才放手让它拟合完整细节。

低通滤波操作本身定义为标准的二维卷积：

\begin{equation}
    G[i,j] = \sum_{u=-k}^{k} \sum_{v=-k}^{k} H[u,v] \cdot F[i{-}u,\, j{-}v]
\end{equation}

其中~$H$~为卷积核，$F$~为输入图像，$G$~为滤波输出。也就是说，对图像中每个像素，取其周围~$(2k{+}1) \times (2k{+}1)$~邻域的加权平均，权重由核~$H$~决定。作者实验了均值滤波（mean filter）、高斯滤波、双边滤波和~Sobel~滤波四种方案，发现均值滤波效果最佳——它对所有邻域像素赋予相同权重，比高斯核更彻底地抑制高频信号，也比双边滤波更不容易保留边缘噪声。在频率调度策略方面，作者对比了线性、余弦、余弦重启、指数和阶跃（step）五种衰减函数，发现阶跃调度显著优于其他方案。原因在于：阶跃调度让~ADC~在固定的频率级别内充分优化当前高斯预算，而不是在持续变化的频率信号下同时处理优化和新高斯的插入。

\textbf{关键参数。}在所有实验中，Opti3DGS~使用阶跃衰减调度，初始均值滤波核大小为~$K_1{=}15$（即~$15 \times 15$~窗口），频率调制仅作用于训练的前~12\,000~次迭代（共~30\,000~次），之后使用原始图像完成剩余训练。该方法不引入任何额外的可学习参数，可即插即用地嵌入任意~3DGS~框架。

\textbf{实验结果。}在~Mip-NeRF 360、Deep Blending~和~Tanks and Temples~三个数据集（共~13~个场景）上的实验表明，Opti3DGS~应用于基线~3DGS~时实现了~62\%~的高斯数量减少、40\%~的训练~GPU~内存减少和~20\%~的优化时间减少，同时仅伴随~0.03\,dB~的平均~PSNR~下降。值得注意的是，在纹理较少的~Deep Blending~数据集上，PSNR~反而提升了~0.09\,dB，同时高斯数量减少了~26\%。此外，Opti3DGS~的渐进式管道还天然产出了多级细节表示（level-of-detail），可作为渲染、流式传输等下游应用的副产物。

Opti3DGS~的频率调制依赖于手动设计的五级阶跃调度和固定的均值滤波核（初始~$K_1{=}15$），滤波器类型和衰减策略需要人工选择。AutoOpti3DGS（2025.06）~\cite{2506_23042}~的核心改动是将滤波器本身从固定设计变为可学习参数：用离散小波变换（DWT）替代均值滤波，其中低通分析矩阵固定（保持粗尺度基础结构），高通分析矩阵~$H_0, H_1$~初始化为零（训练初期高频子带~LH/HL/HH~全部为零，等效于~Opti3DGS~的低通滤波），通过梯度下降逐步学习。高通滤波器向~Haar~小波参考收敛的过程由辅助小波损失驱动：$\mathcal{L}_\text{wavelet} = \sum_{(d)} \| H_{(d)} - H^\text{ref}_{(d)} \|_2^2$，其中~$H^\text{ref}$~为标准~Haar~高通滤波器（只有两个非零系数~$[-1/\sqrt{2}, 1/\sqrt{2}]$，结构最简且天然满足正交性）。该损失的权重~$\alpha_H$~由渲染图像与原始图像的高频子带差异自适应调节：$\alpha_H = \frac{\| X_\text{HH} - X^\text{gt}_\text{HH} \|_1}{\| X^\text{gt}_\text{HH} \|_1 + \epsilon}$，训练初期差异大时高通滤波器快速收敛（快速引入高频细节），后期差异减小时收敛放缓（避免过早引入过多高频信息触发高斯爆炸）。总损失为~$\mathcal{L} = \mathcal{L}_\text{3DGS} + \alpha_H \mathcal{L}_\text{wavelet}$，仅比~Opti3DGS~多一个超参数。核心机制是：高通滤波器从零开始逐步逼近~Haar~参考，等效于一个\textbf{自动}的粗到细课程学习——无需~Opti3DGS~的手动阶跃调度和固定滤波核大小，频率恢复速度由数据本身驱动。后续工作（2026.02）~\cite{2602_14199}~将~1~级小波分解扩展为多级，通过递归分解低频子带构建更深的课程学习序列，进一步减少了高斯数量。

\subsubsection{内存有界的训练框架}

本节介绍在训练过程中显式约束内存峰值使用量的方法，包括~Gaussians on a Diet~的峰值内存约束和~Taming 3DGS~的预算控制策略。

\paragraph{峰值内存约束的高斯训练}

Gaussians on a Diet（2026.04）~\cite{2604_20046}~提出一种基于迭代增长-剪枝循环的内存有界训练框架，将~3DGS~训练的峰值内存消耗严格控制在预设上限之下，从而实现在边缘设备上的完整训练。

\textbf{峰值内存：边缘设备训练的核心瓶颈。}~3DGS~训练的内存瓶颈并非来自最终的模型大小，而是来自训练过程中高斯数量的剧烈波动。原始~ADC~在训练早期会急剧增加高斯数量，通常在~7K--15K~迭代时达到峰值，高斯总数膨胀至最终数量的~2--3~倍。例如~bicycle~场景中，高斯数量瞬间膨胀到数百万级别。这一瞬态峰值对~GPU~内存提出了刚性需求：即使最终模型可以塞进边缘设备的内存，训练过程却因峰值超标而无法启动。现有的后训练剪枝方法（如~LightGaussian、Mini-Splatting）虽然能压缩推理时的模型体积，但训练阶段的高斯爆炸已经发生，峰值内存消耗早已不可逆。换言之，仅靠推理端压缩无法解决训练端内存不足的根本问题。

\textbf{迭代增长-剪枝循环的设计。}直觉上，如果高斯数量可以被控制在一条近似水平线上波动，而非沿抛物线冲到峰值再回落，那么训练内存就变成了一个可控的常数。Gaussians on a Diet 正是基于这一思路：每~$T_{\text{cycle}}$（例如~50）次迭代执行一轮增长-剪枝循环，交替添加新信息高斯、移除冗余高斯。具体而言，循环包含三个步骤：

\begin{enumerate}
    \item \textbf{混合梯度增长}：基于位置梯度与颜色梯度的混合判据~$\tau_i = \tau_i^{\text{pos}} + \tau_i^{\text{col}}$~识别欠拟合区域并执行克隆或分裂。其中颜色梯度~$\nabla c_k = \nabla_{c_k}\mathcal{L} \cdot \prod_{j=1}^{k-1}(1-\alpha_j)$~仅受透射率影响，不受重叠高斯颜色干扰，能够更准确地定位模糊和低对比度区域。克隆时，新高斯并非简单复制原位置，而是沿累积位置梯度方向做微小平移~$\mu_{\text{new}} = \mu_{\text{old}} + \sum_i \nabla_i / N$，避免重叠高斯收到相似梯度更新而长期滞留。
    \item \textbf{自适应高斯补偿}：这是该框架的关键创新。在训练后期（如~10K--15K~迭代），框架识别渲染误差最高的~$K$~个像素，利用~$\alpha$-混合深度将这些像素反投影至三维空间，直接在欠拟合位置生成新高斯并赋予真实像素颜色作为初始值。这种基于像素级误差的补偿机制能够恢复混合梯度增长可能遗漏的高频纹理区域。
    \item \textbf{等量剪枝}：当高斯总数超过目标峰值~$F$~时，移除与新增数量相同的最低重要性高斯。重要性通过基于射线的混合贡献度量计算：$R_i = \sum_{r \in \mathcal{R}_f} \alpha_i^r w_i^r$，其中~$w_i = \alpha_i \prod_{j=1}^{i-1}(1-\alpha_j)$~反映每个高斯对最终像素颜色的实际贡献。该指标可在现有渲染管线中高效计算，每~100~次迭代执行一次，几乎不引入额外开销。
\end{enumerate}

\textbf{设计动机与理论支撑。}该框架的核心理念源自彩票假设（Lottery Ticket Hypothesis）：在一个密集高斯模型中，存在一个稀疏子集能够达到与完整模型相当的性能。Gaussians on a Diet~将这一假设推广至~3DGS——不需要先训练一个大模型再剪枝，而是从一开始就在严格内存约束下动态搜索最优稀疏高斯子集。与~Taming 3DGS~的单向增长策略不同，迭代增长-剪枝允许框架频繁地试错：每轮循环移除``病态''高斯并在需要的位置重新生成``健康''高斯，避免了单向增长中早期容量不足导致的累积误差。实验表明，这种动态策略在视觉保真度和内存效率上均优于一次性剪枝方案。

\textbf{边缘设备实测。}在~NVIDIA Jetson AGX Xavier~上的实测表明，Gaussians on a Diet~实现了与原始~3DGS~相当的视觉质量（在~Tanks and Temples~数据集上~PSNR~提升约~0.5~dB，优于~Taming 3DGS~0.15~dB），峰值训练内存消耗降低了近~80\%，峰值高斯数量减少超过~6$\times$。此外，该方法还配合并行数据加载器动态预取训练数据，进一步节省超过~5~GB~的内存占用，使完整的~3DGS~训练流程首次在~32~GB~内存级别的边缘设备上变得可行。

\paragraph{预算控制的纯粹增密}

与 Gaussians on a Diet 的迭代增长-剪枝循环不同，Taming 3DGS（2024.06）~\cite{2406_15643}~用引导性的、纯粹建设性的增密策略替代原始~ADC，在训练全过程中只增不减——从不执行任何剪枝操作，模型大小以受控方式单调增长直至精确命中用户预设的预算上限。

\textbf{从增长-剪枝到纯粹建设性：为什么舍弃循环。}Gaussians on a Diet~的增长-剪枝循环允许频繁试错，但也引入了不稳定性：每轮剪枝移除的高斯可能恰好在下一轮增长中被重新生成，造成冗余的振荡。Taming 3DGS~基于对~3DGS~增密行为的实证分析发现，原始方法在~Mip-NeRF 360~室外场景中，每步新增高斯数量呈现二次递减趋势——即增密本身具有自然的减速特性。利用这一规律，Taming 3DGS~用一条从~SfM~初始点数出发、精确终止于预算~$B$~的抛物线调度曲线~$A(x) = -\frac{S-2N+N^2x^2}{N^2} + 2x + B$（$N$~为增密步数、$S$~为初始点数）控制每步新增高斯数量，使模型增长过程完全可预测且无峰值溢出。由于不执行剪枝，框架通过累积差值补偿自然剪枝（低不透明度高斯的逐渐消失）：每步计算当前数量与累积目标之间的差额，补增相应数量的新高斯，确保最终数量精确命中预算。

\textbf{分数引导的可控增密。}原始~ADC~仅依赖位置梯度阈值做增密决策，无法区分``真正需要更多采样''与``高斯本身被错误放置''两种情况，容易导致重复复制错误高斯。Taming 3DGS~提出一种多指标评分函数~$F$，为每个高斯计算综合得分~$S_g$~作为增密采样权重。该评分融合了八项互补信号：位置梯度~$\|\nabla_g\|_2$（权重~50，继承自原始~3DGS~的~3D~不连续性检测器）、像素覆盖数~$c_g^i$（0.1，识别投影过大的模糊高斯）、像素到中心距离~$D_g^i$（50，捕获细长``针状''高斯）、逐像素显著性~$S_v^i$（10，由~L1~损失与拉普拉斯滤波器组成，定位高频欠采样区域）、混合权重~$B_g^i$（50，优先增密对渲染贡献最大的高斯）、深度~$z_g^i$（5，兼顾前景优先与背景不遗漏）、不透明度~$o_g$（100，强抑制低不透明度``浮游''高斯的增密）、尺度~$s_g$（25，惩罚过大高斯以提升未见视角泛化性）。所有指标经中值缩放去除异常值后与光度损失相乘，再累加为加权总分。增密时，从全部高斯中按~$S_G$~作为采样权重随机抽取~$B = T - |G|$~个候选，按原始规则执行分裂（大高斯）或克隆（小高斯）。

\textbf{低频增密：给优化留出自我纠错的时间。}Taming 3DGS~将增密频率降至原始~3DGS~的~1/5（每~500~次迭代执行一次），这一看似简单的调整背后有深刻的设计动机。高位置梯度可能由两种原因引起：该区域确实采样不足，或高斯本身被错误放置。如果增密过于频繁，优化器尚未将错误放置的高斯的不透明度降低（3DGS~的自然纠错机制），新的高斯就已经被复制到同样的错误位置，导致``错误复制-错误复制''的恶性循环。低频增密给~3DGS~足够的时间和观测来通过降低不透明度消除不当高斯，再基于更可靠的评分进行增密。此外，分数计算仅需在~10~个均匀采样的训练视角上进行，而非全部视角，进一步控制了增密过程的计算开销。

\textbf{高不透明度高斯：突破单高斯表达力上限。}标准~3DGS~的高斯基元受限于固定的钟形衰减曲线，难以用少量基元精确建模不透明表面。Taming 3DGS~借鉴了层级高斯方法中使用不透明度~$>1$~的截断高斯来近似高斯簇外观的思路，在训练达到中点（15K~迭代）后将常规高斯转换为高不透明度高斯：将不透明度激活函数替换为~abs~函数，并在渲染时将混合权重从上方截断至~$1$。这一改动使单个高斯能够表达更锐利的不透明表面，在预算受限场景下对~PSNR~指标的提升尤为显著。

\textbf{管线加速。}除增密策略外，Taming 3DGS~还通过分析原始~3DGS~管线的反向传播过程，推导出数值等价但计算更高效的梯度计算和属性更新方案，进一步缩短训练时间。

\textbf{整体效果。}综合上述四项改进——可预测的抛物线调度、分数引导的随机采样增密、低频增密策略、高不透明度高斯转换——Taming 3DGS~在预设预算约束下实现了极高的帕累托效率：以~4--5$\times$~更小的模型和~4--5$\times$~更快的训练时间达到了与原始~3DGS~竞争的质量指标。与~Gaussians on a Diet~相比，Taming 3DGS~的纯粹建设性策略更为简洁（无需剪枝逻辑和误差补偿模块），但在内存预算的极致压缩下，Gaussians on a Diet~的增长-剪枝循环能够通过反复试错更精细地优化稀疏高斯子集的分布，在~Jetson AGX Xavier~等边缘设备上展现出更强的实际部署能力。

\subsubsection{训练流程优化}

本节介绍从训练流程本身出发提升效率的方法，包括 DashGaussian 的优化复杂度调度和 SkipGS 的反向传播门控策略。

\paragraph{优化复杂度调度}

\textbf{Baseline 做法：全分辨率训练的计算瓶颈。}原始~3DGS~训练始终在全分辨率下渲染和优化。渲染分辨率和基元数量共同决定每次迭代的计算量。渲染分辨率越高，需要光栅化的像素越多；基元越多，需要排序和混合的高斯越多。然而，训练早期场景的粗糙几何尚未建立，此时使用全分辨率渲染无异于用高精度仪器测量一幅草图，大量计算花费在优化尚未收敛的高频细节上。这种固定策略导致~3DGS~的训练时间主要消耗在早期低效的全分辨率迭代中。

\textbf{核心洞察：渐进频率拟合。}DashGaussian（2025.03）~\cite{2503_18402}~从频率分析的角度重新审视~3DGS~优化。根据~Nyquist~采样定理，将高分辨率图像下采样为低分辨率图像等价于丢弃高频分量、保留低频分量。因此，3DGS~优化可以重新表述为：先拟合训练视角的低频分量（低分辨率即可监督），再逐步拟合高频分量（需要全分辨率监督）。这一重新表述直接回答了一个关键问题——低分辨率训练视图能在多大程度上替代全分辨率？答案是：只要低频分量已被充分拟合，就应切换到更高分辨率以捕获残余的高频细节。

\textbf{频率引导的分辨率调度。}基于上述洞察，DashGaussian~提出自适应的分辨率调度策略。对于训练视图的某个分辨率~$(h, w)$，定义其频率显著性为
\begin{equation}
    \mathcal{X}(F) = \frac{1}{N} \sum_{n=1}^{N} \sum_{i=1}^{h} \sum_{j=1}^{w} \|F^n(i,j)\|_2,
\end{equation}

其中~$F^n = \mathrm{DFT}(I^n)$~为第~$n$~个训练视图的离散傅里叶变换（Discrete Fourier Transform, DFT）。该度量强调高频分量——高分辨率视图的~$\mathcal{X}(F)$~包含更多高频能量，低分辨率视图的~$\mathcal{X}(F_r)$~仅包含低频能量。直觉上，如果~$\mathcal{X}(F) - \mathcal{X}(F_r)$~远大于~$\mathcal{X}(F_r)$，说明高频信息尚未被充分优化，应分配更多迭代给全分辨率渲染；反之则低分辨率足以完成大部分优化。据此，全分辨率渲染应分得的迭代比例为
\begin{equation}
    f(F, F_r) = \frac{\mathcal{X}(F) - \mathcal{X}(F_r)}{\mathcal{X}(F)},
\end{equation}

切换分辨率的时间点为第~$s_r = S \cdot \mathcal{X}(F_r) / \mathcal{X}(F)$~次迭代（$S$~为总迭代数）。该策略可推广至多个中间分辨率~$r_1 < r_2 < \cdots < r_m$，其中最大下采样因子~$r_m$~满足~$\mathcal{X}(F_{r_m}) = \frac{1}{a}\mathcal{X}(F)$（$a$~默认为~4），其余因子在~$\mathcal{X}(F)$~和~$\mathcal{X}(F_{r_m})$~之间线性采样。在相邻切换点之间，分辨率通过~$r_i$~和~$r_{i-1}$~的逆平方插值平滑过渡。

\textbf{分辨率引导的基元调度。}分辨率调度只解决了一半问题。过密的高斯基元仅在推理质量上带来边际收益，却显著拖慢优化（2024）~\cite{fang2024minisplattingrepresentingscenesconstrained}~。DashGaussian~进一步提出渲染分辨率与基元数量协调增长的方案：低分辨率阶段只需少量基元描述粗糙几何，高分辨率阶段才需要大量基元捕获精细细节。具体而言，第~$i$~次迭代的基元数量由下式控制：

\begin{equation}
    P_i = P_{\text{init}} + (P_{\text{fin}} - P_{\text{init}}) \cdot \frac{1}{(r^{(i)})^{2 - i/S}},
\end{equation}

其中~$r^{(i)}$~为第~$i$~次迭代的渲染分辨率因子，指数项~$2 - i/S$~从~2~线性衰减到~1。这一设计使得训练早期（$i$~较小，指数接近~2）基元增长被强力抑制，中期随分辨率提升加速增长，形成一条凹函数曲线（面积小），从而显著降低整体优化复杂度。

\textbf{动量驱动的基元预算。}$P_{\text{fin}}$~的确定是另一个难题——现有方法通常依赖数据集先验手动设定，缺乏自适应性。DashGaussian~受动量优化器启发，将~$P_{\text{fin}}$~建模为动量，每次迭代自然增密的高斯数量~$P_{\text{add}}$~建模为力，通过指数移动平均动态更新：

\begin{equation}
    P_{\text{fin}}^{(i)} = \max\!\big(\beta_1 \cdot P_{\text{fin}}^{(i-1)},\; P_{\text{fin}}^{(i-1)} + P_{\text{add}}^{(i)}\big),
\end{equation}

其中~$\beta_1 = 0.98$~为动量系数，取最大值操作确保~$P_{\text{fin}}$~单调递增。该机制无需指定基元数量上界，即可在不同数据集上自适应确定合理的最终基元数量。

\textbf{加速效果与即插即用性。}DashGaussian~在多个~3DGS~主干上将优化加速~45.7\%，且渲染质量与基线持平甚至略有提升。具体而言，该加速来自两方面贡献：频率引导的分辨率调度减少了早期低效的全分辨率迭代，分辨率引导的基元调度避免了过度增密带来的冗余计算。该方法无需引入额外可学习参数，可作为即插即用模块嵌入任意~3DGS~训练框架。

\paragraph{后增密阶段的反向传播门控}

\textbf{锚点回顾。}DashGaussian~通过频率引导的分辨率调度减少训练前期的计算冗余，其核心假设是低分辨率训练视图足以拟合低频分量。与这种前端调度策略不同，SkipGS（2026.03）~\cite{2603_08997}~将注意力转向训练流程的另一端——增密阶段结束后的精炼阶段（post-densification phase），聚焦于反向传播的冗余问题。

\textbf{精炼阶段的反向传播冗余。}3DGS~训练通常分为两个阶段：增密阶段（高斯基元不断分裂和克隆）和精炼阶段（基元集合固定，仅微调外观与几何参数）。SkipGS~观察到，进入精炼阶段后，不同视角的收敛速率差异显著。受可见性、遮挡模式和视角相关外观等因素影响，部分视角仍能提供有效梯度信号，但许多视角的损失已稳定在平台期。然而，标准~3DGS~训练无论视角状态如何，仍然对每次采样的视角执行完整的反向传播。考虑到反向传播占据每次迭代的主要计算开销（远超前向传播），这种更新冗余导致精炼阶段大量计算被浪费。

\textbf{视角自适应的偏离分数。}为精确识别哪些视角仍需梯度更新，SkipGS~为每个训练视角维护损失序列的指数移动平均（EMA）基线：

\begin{equation}
    \hat{\ell}_v^{(t)} = \alpha\, \hat{\ell}_v^{(t-1)} + (1-\alpha)\, L_v^{(t)},
\end{equation}

其中~$\alpha \in (0,1)$~控制时间尺度，$L_v^{(t)}$~为视角~$v$~在第~$t$~次迭代被采样时的前向损失。每当视角~$v$~被采样时，SkipGS~计算尺度不变的偏离分数（deviation score）：

\begin{equation}
    s_v^{(t)} = \frac{\hat{\ell}_v^{(t)}}{L_v^{(t-1)} + \epsilon},
\end{equation}

其中~$\epsilon = 10^{-8}$~保证数值稳定性。偏离分数显著大于~1~表示该视角损失相对于近期基线出现尖峰，意味着该视角仍未充分优化（或发生了回归），值得投入反向传播计算；偏离分数低于~1~表示该视角当前表现优于近期基线，此时反向传播通常收益递减，应将计算资源优先分配给高偏离视角。无论是否执行反向传播，前向传播始终执行并更新~EMA，确保偏离分数追踪真实的视角损失趋势。

\textbf{门控判定与预算约束。}基于偏离分数，SkipGS~的门控判定为：仅当~$s_v^{(t)} > 1$~时执行反向传播（$g_t(v)=1$），否则跳过（$g_t(v)=0$）。然而，无约束的跳过可能导致优化陷入梯度饥饿。为此，SkipGS~引入两个安全保障：（1）\emph{预热窗口}——在精炼阶段的前~$W$~次迭代中无条件执行所有反向传播，用于稳定视角~EMA~并收集门控统计信息；（2）\emph{最小反向传播预算}——维持累计反向传播比率~$\beta_{\text{cum}}(t) = \frac{1}{t}\sum_{i=1}^{t} g_i$，当该比率低于目标下界~$\beta_*$~时强制执行反向传播。$\beta_*$~从预热统计中自动校准：

\begin{equation}
    \beta_* = \beta_{\text{lo}} + (1 - \beta_{\text{lo}}) \cdot \beta_W,
\end{equation}

其中~$\beta_{\text{lo}} = 0.5$~为全局下界，$\beta_W$~为预热期间门控判定会触发反向传播的比例。直觉上，若预热表明多数视角已易于优化（$\beta_W$~较小），则允许较低的最小预算；若预热表明多数视角仍需反向传播（$\beta_W$~接近~1），则强制接近全量反向传播的预算。这种自动校准机制避免了逐场景调参。

\textbf{方法正交性与加速效果。}SkipGS~仅改变何时执行反向传播，不修改渲染器、高斯表示或损失函数，因此与现有加速方法（如~DashGaussian~的分辨率调度、高斯压缩和基元数量调控）正交，可叠加使用以获得累积加速。在~Mip-NeRF 360~数据集上，SkipGS~将端到端训练时间减少~23.1\%，渲染质量与基线持平。

\begin{takeaway}
训练加速与高效表示从分辨率调度、自适应密度控制和反向传播预算三个正交维度加速 3DGS 训练。SkipGS 的门控机制仅改变何时执行反向传播，与现有加速方法正交可叠加使用，将训练时间减少 23.1\%。关键洞察是训练中并非所有视角都需要完整的反向传播，基于残差的自适应预算分配可以在不影响质量的前提下显著节省计算。
\end{takeaway}


%=====================================================================


\section{视频与流式渲染}
\label{sec:11_4}

\subsection{问题背景}

本节讨论三条技术路线：（1）渐进式场景加载；（2）分层表示与流式传输；（3）动态场景的压缩与流式传输。

前述压缩方法产生的都是一个单一的压缩文件，用户必须完整下载后才能开始渲染。对于一个~100~MB~的压缩场景，在~2~Mbps~带宽下需要约~7~分钟的黑屏等待，这在实际应用中不可接受。流式渲染（streaming rendering）的核心思想是：将压缩模型组织为渐进式结构，用户可以先下载一个低质量的"骨架"开始交互，然后在带宽允许时逐步接收细节增强数据。

视频场景的压缩与流式渲染面临额外的时域挑战。4D~高斯模型通常包含变形网络或~4D~表示，其参数量远大于静态场景，且帧间存在大量时域冗余需要消除。

\subsubsection{渐进式场景加载}

本节介绍将~3DGS~场景组织为从粗到细的渐进式结构以支持流式传输的方法，包括~PRoGS~的渐进渲染、GoDe~的可扩展压缩以及~Matryoshka~的嵌套式连续细节层次。

\paragraph{渐进式高斯渲染}

传统~3DGS~压缩方法产出一个完整的压缩文件，用户必须下载全部数据后才能开始渲染。PRoGS（2024.09）~\cite{2409_01761}~改变了这一范式：不再要求完整下载，而是为每个高斯计算一个\textbf{渲染贡献优先级}（rendering contribution priority），按重要性从高到低排序后分块传输，使用户在只接收一小部分高斯时就能看到场景的近似重建。

其核心直觉很简单：在~3DGS~的~$\alpha$-混合渲染中，不同高斯对最终像素颜色的贡献差异巨大——大尺寸、高不透明度的高斯通常比小而透明的"补丁"高斯贡献多得多。PRoGS 利用这一差异，量化每个高斯的重要性并优先传输最重要的那些。

具体而言，PRoGS 回放所有训练视角的渲染过程，逐像素记录每个高斯的\textbf{贡献分数}。在~3DGS~的~$\alpha$-混合公式中，像素颜色为
\begin{equation}
C = \sum_{i=1}^{N} T_i \, \alpha_i \, c_i, \quad T_i = \prod_{j=1}^{i-1}(1-\alpha_j),
\end{equation}

其中~$T_i$~是透射率（transmittance），$\alpha_i$~是不透明度，$c_i$~是颜色值。第~$i$~个高斯对单个像素的贡献为~$T_i \alpha_i$。将该值在所有像素和所有训练视角上求和，即得到该高斯的\textbf{全局贡献分数}：

\begin{equation}
B_{\text{id}} = \sum_{v \in V}\sum_{p \in P}(T_{\text{id},v,p}\,\alpha_{\text{id},v,p}),
\end{equation}

其中~$V$~是训练视角集合，$P$~是每个视角的像素集合。按此分数降序排列，就得到了高斯的优先级序列——最重要的~10\%高斯最先传输，用户即可获得场景的粗略预览；随后更多高斯逐步到达，渲染质量持续提升。

单纯的全局排序有一个实际问题：被更多训练视角覆盖的区域自然会获得更高的累计贡献，导致稀疏采样区域（如背景）被不公平地延迟重建。为解决此偏差，PRoGS 引入\textbf{八叉树细化}（octree refinement）：将场景空间按高斯密度划分为八叉树叶节点，在每个叶节点内独立选取贡献最高的高斯。这样确保场景各部分至少得到部分重建，而非仅优先填充高密度区域。八叉树的最大深度是一个可调参数——增大深度意味着更多地偏向密集区域的细节，减小深度则更均匀地分布重建资源。

此外，由于用户的初始位姿通常是已知的，PRoGS 还支持\textbf{视锥体内优先}（in-frustum prioritization）：在优先级序列中，额外选取当前视锥体内的高斯优先传输，加速用户首帧可见内容的到达。具体做法是将高斯投影到图像平面，保留落在视锥体内（含一定边缘余量）的高斯，提升其在传输队列中的位置。

PRoGS 的一个重要特性是\textbf{压缩方法无关性}：由于排序仅依赖贡献分数而不涉及高斯的内部参数编码，它可以与任何基于逐高斯存储的压缩方法组合使用。作者将~PRoGS~集成到~Niedermayr~等人的压缩框架中进行了验证，实验表明该组合在保持压缩效率的同时，将流式场景的初始可见时间缩短了数倍。

\paragraph{可扩展压缩与渐进细节层次}

与~PRoGS~基于渲染贡献分数的启发式排序不同，GoDe（2025.01）~\cite{2501_13558}~将排序准则建立在更坚实的理论基础上：\textbf{梯度敏感性}（gradient sensitivity）。其出发点是一个关键观察——训练好的~3DGS~模型中存在结构化的信息分布，而这一分布可以通过优化过程中自然产生的梯度信号来揭示。具体而言，在冻结模型参数后对训练视角做一次前向-反向传播，聚合每个高斯所有参数（位置、不透明度、缩放、旋转、外观系数）的梯度~$\ell_2$~范数，即得到该高斯的\textbf{重要性分数}：$s_i = \|\nabla_{\theta_i} \mathcal{L}\|_2$。这一度量直接反映每个基元对重建损失的贡献大小，且\textbf{模型无关}——不依赖特定的架构设计或渲染机制，对显式~3DGS~和神经~3DGS~变体均适用。

这一排序准则的设计动机源于对~PRoGS~一类方法的理论审视。PRoGS~的渲染贡献分数基于高斯在~$\alpha$-混合中的实际像素贡献，本质上是一个\textbf{事后统计量}——它衡量的是该高斯在渲染中实际做了多少，而非该高斯对重建质量有多重要。两者之间存在微妙的区别：一个大尺寸、高不透明度的高斯可能在渲染中贡献显著，但如果移除它后其他高斯可以轻易补偿其作用，那么它对重建质量的实际重要性可能很低。梯度敏感性则从优化的视角直接回答了移除该高斯对损失函数影响多大这一问题，提供了一阶近似估计。

GoDe~利用梯度敏感性构建\textbf{渐进层次}（progressive hierarchy），采用自上而下的迭代策略：从完整模型~$G_{L-1}$~出发，在每一步计算当前集合中所有高斯的梯度范数，将分数最低的~$k_l$~个高斯分配到当前层作为增强集~$E_l$~并从后续考虑中移除，然后重置梯度、在剩余高斯上重复此过程。这种\textbf{迭代式}构建相比一次性全局排序有显著优势——它避免了后者中常见的偏差：如果一次性按梯度大小排序所有高斯，被多个视角共享的结构性高斯会因梯度累积而始终排在高位，导致粗糙层的代表性不足。迭代式方法在每步重新评估剩余高斯的相对重要性，使层次结构更加均衡和稳定。最终，全部高斯被分解为基集~$G_0$~和一系列互不相交的增强集：$G = (G_0, E_1, \ldots, E_{L-1})$，每层的表示为~$G_l = G_0 \cup \bigcup_{i=1}^{l} E_i$。各层的规模按几何级数递增：$|G_l| = c \cdot b^l$，使得粗糙层获得更多容量，细节层的增量逐步减小——这契合信息论中的率-失真规律：前几个高斯捕获粗结构与主要外观，后续高斯的边际贡献递减。

层次构建完成后，GoDe~执行一次\textbf{量化感知联合微调}：在~$F$~次迭代中，每次均匀随机采样一个层次~$l$，仅用~$G_l$~中的高斯进行渲染和优化。随机层次采样迫使低层（被所有上层共享）的参数保持鲁棒，确保同一组基元在不同预算下都能产生合理的重建。微调期间还对主要存储参数（外观系数、不透明度、缩放、旋转）模拟量化，使用直通估计器（straight-through estimator）传播梯度，并对外观系数施加~$\ell_1$~稀疏惩罚以进一步压缩存储。

GoDe~的最终成果是一个\textbf{可扩展比特流}：基集~$G_0$~和各增强集独立压缩（使用~ZSTD~及两个共享字典，分别用于几何和外观流），解码~$G_0$~即得粗糙表示，逐步追加增强层则持续提升质量。与~PRoGS~的关键差异在于，GoDe~的渐进层次是\textbf{固定的、离线构建的}，而非运行时按需排序；这使得它可以精确控制每层的率-失真特性，并支持在\textbf{单个训练模型}上提供多个离散操作点——用户根据当前设备内存、带宽或渲染预算选择解码到哪一层即可，无需重新训练或重新编码。为缓解相邻层切换时增强高斯突然出现造成的跳变伪影，GoDe~还引入了运行时\textbf{不透明度插值}：新增层的增强高斯从零不透明度开始，逐步淡入至目标值，实现平滑的层次过渡。

\paragraph{嵌套式连续细节层次}

前述~PRoGS~和~GoDe~都遵循``先训练完整模型，再排序/分层''的两阶段范式：训练阶段产出无序高斯集合，后处理阶段按贡献或梯度敏感性排序后打包为层次结构。Matryoshka GS（MGS）（2026.03）~\cite{2603_19234}~转变了这一思路——\textbf{在训练过程中}就学习一个有序的高斯序列，使得任何前缀（前~$k$~个基元）本身就是一段连贯的场景表示，无需后处理排序或辅助数据结构。名字来源于俄罗斯套娃（Matryoshka doll）：正如每一层套娃都是完整的，MGS~的每个前缀都是一个功能完整的场景近似。

\paragraph{设计动机：连续的保真度-速度权衡}

无论是~PRoGS~的后处理排序还是~GoDe~的固定层次，都隐含了一个假设——质量档次是离散的。但在实际部署中，计算预算随硬件（从高端~GPU~到移动设备）、场景复杂度和运行时负载而\textbf{连续}变化。离散的层次切换无法平滑跟踪这种变化，且档次之间的跳变会产生可见的~pop-in/pop-out~伪影。MGS~的目标是消除这一限制：训练单个模型，在部署时通过简单地截断高斯数组的前~$k$~个元素，即可在任意预算点渲染，获得一条连续、密集的保真度-速度~Pareto~前沿。在~Mip-NeRF~360~基准上，MGS~的~$AUC_{fps}$~达到~54.46，远超下一个连续~LoD（Level of Detail，细节层次）~竞争者~CLoD-3DGS~的~28.94。

\paragraph{有序前缀表示}

为实现上述目标，MGS~首先为每个高斯基元~$g_i$~分配一个标量\textbf{重要性分数}~$s(g_i)$，并按降序排列所有高斯，得到一个置换~$\pi$，使得~$s(g_{\pi(1)}) \geq s(g_{\pi(2)}) \geq \cdots \geq s(g_{\pi(N)})$。预算为~$k$~的前缀定义为~$G_{<k} := \{g_{\pi(1)}, \ldots, g_{\pi(k)}\}$。由于前缀具有嵌套性（$G_{<k} \subseteq G_{<k'}$~当~$k < k'$），渲染~$G_{<k}$~即可得到该预算下的重建结果。论文经验性地发现，以\textbf{不透明度}降序排列作为重要性分数最为稳健——不透明度反映了基元的可见度和辐射贡献，使最早的前缀就能捕获场景的主导结构。消融实验表明，在~10\%~基元预算下，不透明度降序达到~22.2\,dB~PSNR，而次优的~SH~能量降序仅为~17.6\,dB。

\paragraph{随机预算训练}

训练单一模型使其在\emph{每个}预算级别都表现良好是~MGS~的核心挑战。直接优化所有~$N$~种预算不可行，因此作者提出\textbf{随机预算训练}（stochastic budget training）：每次训练迭代，从均匀分布~$\text{Unif}(r_{\min},\,1)$~中采样一个保留比例~$r$，计算对应的前缀大小~$k = rN$，然后\textbf{同时渲染两个版本}——前缀~$G_{<k}$~和完整集合~$G_{<N}$——并最小化联合损失：

\begin{equation}
\mathcal{L}_{\text{MGS}} := \mathcal{L}(G_{<k};\,I,\,c) \;+\; \lambda\,\mathcal{L}(G_{<N};\,I,\,c),
\end{equation}

其中~$\mathcal{L}$~是标准~3DGS~重建损失（$L_1$+D-SSIM~混合），$\lambda \geq 0$~平衡前缀项和完整集合项。前缀项迫使模型从部分子集产生强重建，完整集合项锚定满容量性能。均匀采样确保训练过程中~$[r_{\min}N,\,N]$~范围内的每个预算都被等概率访问，产生对全预算谱的无偏覆盖。每步仅需\textbf{两次}前向渲染，无论~$N$~多大。

\paragraph{动态重排序}

由于梯度更新不断修改高斯参数，重要性分数~$s(g_i)$~在训练过程中持续演化。因此，每次训练迭代后，MGS~重新计算置换~$\pi$~以保持排序约束在当前参数下成立。这确保每个前缀~$G_{<k}$~始终包含当前参数下最重要的~$k$~个基元。

\paragraph{与~PRoGS/GoDe~的关键差异}

MGS~与前述方法有三个根本区别。第一，\textbf{训练阶段排序}~vs.~\textbf{后处理排序}：PRoGS~和~GoDe~在模型训练完成后才建立层次，而~MGS~将排序融入训练循环，使模型主动学习在不同预算下都有效的基元参数。第二，\textbf{连续}~vs.~\textbf{离散}操作点：PRoGS~的贡献排序虽然理论上支持任意截断点，但并未针对部分预算进行优化训练；GoDe~产出固定的有限层次。MGS~的随机预算训练使每个整数~$k$~都是有效的操作点，在~Mip-NeRF~360~上仅评估~12~个预算比例就得到平滑的质量-速度曲线。第三，\textbf{零架构修改}：MGS~仅改变训练目标函数和排序策略，不引入任何新的网络模块、辅助数据结构或额外的存储开销。事实上，MGS~的满容量性能接近甚至超过其~3DGS-MCMC~主干网络——在~Mip-NeRF~360~和~Tanks and Temples~上仅落后~0.20\,dB（28.20~vs.~28.40，24.56~vs.~24.76），在~Deep Blending~和~BungeeNeRF~上则分别超越~0.78\,dB~和~0.09\,dB，表明随机预算目标本身可充当有益的正则化项。

\subsubsection{分层表示与流式传输}

本节介绍基于分层和分块结构的~3DGS~流式传输方法，包括~L3GS~的分层场景分发和~LapisGS~的动态不透明度优化。

\paragraph{分层高斯场景分发}

上述方法（HGS、GoDe、Matryoshka GS）分别从八叉树组织、梯度排序、随机预算等角度构建~3DGS~的层次结构，但它们聚焦于单机端的层次化表示本身，尚未解决一个更上游的问题：当~3DGS~场景存储在远程服务器上时，用户如何通过网络逐步获取场景并实时渲染？标准~3DGS~场景（如~Mip-NeRF 360~的~bicycle~场景）可达~1.52~GB、包含超过百万个高斯基元，若要求客户端完整下载后才能渲染，将导致数十秒甚至数分钟的黑屏等待。L3GS（2025.04）~\cite{2504_05517}~正是在这一背景下提出的，它构建了面向网络分发的完整分层~3DGS~系统，覆盖从场景创建、调度决策到用户实验的端到端流程。

\textbf{设计动机与问题定义。}
L3GS~的核心动机来自~VR~头显用户的需求：用户佩戴头显在~3DGS~场景中自由行走（6~自由度），期望打开场景后立即看到内容，而非等待完整下载。这带来三个耦合的挑战：（1）场景体积大，必须渐进传输；（2）用户视口不断变化，传输策略需要区分高斯基元的优先级；（3）网络带宽波动，调度算法需要自适应。L3GS~的解决思路是将~3DGS~场景组织为分层的叠加结构——先传输一个粗糙的``基础层''，再逐步叠加``增强层''，使渲染质量随传输进度平滑提升。

\textbf{分层场景的创建。}
L3GS~的分层训练包含两个步骤。

\textit{步骤一：自顶向下剪枝至目标规模。}
从默认预训练的~3DGS~模型（约~100~万基元）出发，采用~LightGaussian~的全局显著性评分（Global Significance Score, GSS）对基元排序，反复执行``剪枝—恢复''循环——剪去评分最低的~$r$~比例基元，再通过标准训练恢复视觉质量——直到总基元数精确达到目标值~$d$。例如，目标~180k~基元时，最终精确保留~$d = 180{,}000$~个基元并按~GSS~降序排列。这一``自顶向下''策略在实践中优于``自底向上''逐步增长的方法。

\textit{步骤二：逐层冻结式增量训练。}
将排序后的~$d$~个基元均匀切分为~$L$~层（如~$L=4$，每层~45k~基元），然后从低层到高层逐层训练。关键约束是：训练第~$l$~层时，冻结前~$l-1$~层的所有参数（\texttt{requires\_grad=False}），仅允许第~$l$~层新增基元的参数更新。损失函数基于前~$l$~层所有基元的联合渲染结果计算，但梯度仅回传至第~$l$~层。这种``冻结式增量训练''确保各层可以逐层叠加而不互相干扰。此外，L3GS~还利用~Gaussian Grouping~方法为每个基元附加对象~ID~作为额外特征，将场景语义分割为有意义的物体，训练过程中对象~ID~同步优化，使得同一层内不同物体自动获得不同数量的基元分配。

\textbf{调度算法：视口感知的多选择背包问题。}
给定分层场景，调度器需要在每个时间槽~$t$（长度~$T$）内决定下载哪些基元。为此，L3GS~定义了基元效用函数（splat utility），将每个基元~$i$~在物体~$j$~、层~$l$~、时刻~$t$~的效用分解为三个因子的乘积：

\begin{equation}
U_{i,j,l}[t] = \underbrace{w_{i,j,l}^{\text{close}}[t]}_{\text{接近度}} \cdot \underbrace{w_{i,j,l}^{\text{overlap}}[t]}_{\text{投影面积}} \cdot \underbrace{\alpha_{i,j,l}}_{\text{不透明度}}
\end{equation}

其中``接近度''（closeness）衡量基元中心到用户视口中心的~3D~欧氏距离，``投影面积''（overlap）是基元投影到用户~2D~视口后的面积，``不透明度''（opacity）是基元自身的固有属性。接近度和投影面积随用户视口动态变化，不透明度则为静态属性。直觉是：大尺寸、高不透明度、靠近视口中心的基元应优先下载。

在此基础上，L3GS~将调度问题形式化为一个带前序约束的优化问题。以最复杂的``分层+语义分割''场景（Case I）为例，决策变量~$x_{j,l}[t] \in \{0,1\}$~表示是否在时刻~$t$~下载物体~$j$~的第~$l$~层，优化目标为最大化总效用：

\begin{equation}
\max \sum_{j} \sum_{l} \sum_{t} \Delta U_{j,l}[t] \cdot y_{j,l}[t], \quad \text{s.t.} \quad \sum_{j,l} c_{j,l} \cdot x_{j,l}[t] \leq B[t], \;\; \forall\, t
\end{equation}

其中~$\Delta U_{j,l}[t] = U_{j,l}[t] - U_{j,l-1}[t]$~是第~$l$~层相对于第~$l-1$~层的边际效用，$c_{j,l}$~是该层的数据大小，$B[t]$~是预测的可用带宽。关键的前序约束为：第~$l$~层只能在第~$l-1$~层已被下载后才能下载（$x_{j,l}[t] \leq y_{j,l-1}[t]$）。L3GS~证明了该问题等价于带前序约束的广义指派问题（Generalized Assignment Problem），且为~NP-hard。实践中，L3GS~将每个时间槽内的子问题转化为多选择背包问题（Multiple Choice Knapsack），用动态规划在伪多项式时间内求解。

\textbf{视口与带宽预测。}
L3GS~使用~Meta Quest 3~头显采集了~6~名用户在~8~个~3DGS~场景中的~6-DoF~运动轨迹（采样率~36~Hz），并以线性回归预测未来视口（历史窗口~$hw = 0.5$~秒，预测窗口~$pw = 1$~秒）。带宽预测采用谐波均值预测器（harmonic mean predictor），基于户外~5G~用户行走轨迹模拟可变网络条件，缩放后的平均带宽为~11.8~Mbps。

\textbf{实验评估。}
在~8~个公开场景（Mip-NeRF 360、Tanks and Temples、Gaussian Grouping）上，L3GS~创建~4~层各~45k~基元的分层场景，使用~60~秒用户轨迹评估视口~SSIM。结果表明：（1）前~5~秒内，L3GS~显著优于基线——``排序下载''基线使用视口无关的全局评分，无法适应用户观看方向；``独立版本''基线在初始阶段甚至无法完成最低质量版本的下载，导致黑屏（SSIM~为~0）；（2）约~30~秒后所有~180k~基元下载完毕，性能接近预加载理想基线；（3）分层场景的视觉质量与非分层独立版本相当，证明``冻结式增量训练''未显著牺牲质量；（4）调度延迟小于~3~毫秒（2~核~2.2~GHz~CPU），满足实时性要求；（5）带语义分割的调度优于不带分割的调度，因为物体级别的决策空间更大、更灵活。

\paragraph{动态不透明度优化的自适应流式传输}

上述~L3GS~采用``基础层+增强层''的叠加结构，通过冻结式增量训练逐层构建分层场景，再由调度器按视口优先级分批下载基元。这一方案有效地将~3DGS~场景组织为可渐进传输的层次结构，但其训练过程中各层参数完全冻结——包括不透明度在内，这意味着低层基元对渲染的贡献权重无法随高层内容的引入而动态调整。当高层基元捕获了更丰富的细节后，低层中那些仅表征粗略特征的基元仍然以原始不透明度参与渲染，既造成视觉冗余（低频基元与高频基元叠加产生伪影），又导致模型体积膨胀（需要更多基元来补偿未优化的低层贡献）。LapisGS（2024.08）~\cite{2408_14823}~正是在这一背景下提出的改进方法，它继承了~L3GS~的分层渐进训练框架，但引入了三项关键设计来克服上述局限。

\textbf{与~L3GS~的核心差异：累积表示的分层结构。}
L3GS~将排序后的~$d$~个基元均匀切分为~$L$~层，每层包含独立的一组基元，层与层之间互不共享参数。LapisGS~则采用不同的构造路径：它从多尺度图像金字塔出发（将全分辨率图像集~$D_N$~逐级下采样为~$D_0, D_1, \ldots, D_N$），从最低分辨率开始逐级训练。形式化地，分层模型表示为~$\{G_0, \Delta G_1, \Delta G_2, \ldots, \Delta G_N\}$，其中~$G_0$~为基础层，$\Delta G_k$~为第~$k$~个增强层。第~$i$~级的完整表示为~$G_i = G_0 + \sum_{k=1}^{i} \Delta G_k$——即所有低层基元的累积叠加，而非~L3GS~中各层独立的基元集合。这种累积结构使得不同分辨率层之间可以共享视觉信息：低层基元捕获场景的粗略布局和低频特征，高层基元则专注于补充高频细节，避免了跨层的信息冗余。

\textbf{动态不透明度优化。}
这是~LapisGS~与~L3GS~最本质的设计差异。在~L3GS~的冻结式增量训练中，训练第~$l$~层时完全冻结前~$l-1$~层的所有参数（包括位置、协方差、球谐系数和不透明度）。LapisGS~则选择了一条中间路径：训练第~$i$~级增强层~$\Delta G_i$~时，冻结前~$i-1$~级基元的所有参数\textit{除了不透明度}。这一设计的直觉是：（1）3DGS~使用标准~$\alpha$-混合~进行渲染，每个基元的贡献由其不透明度加权，因此调整不透明度可以精细控制各层基元对最终图像的相对影响力；（2）与位置、尺度等属性不同，修改不透明度不会改变基元所承载的空间信息（几何结构和颜色特征），因此不会破坏已学习的场景布局。具体而言，当高层基元开始捕获更精细的纹理和边缘时，优化过程会自动降低低层中那些仅表征粗略特征的基元的不透明度——因为它们在高层信息到位后对渲染质量的边际贡献已显著下降。这一动态调整机制使得模型能够自适应地``修剪''低层基元的视觉影响力，而无需真正删除它们，从而在保持分层结构完整性的同时实现表示紧缩。

\textbf{有效占用图与连续层级过渡。}
LapisGS~引入了占用图（occupancy map）~$O_i$~来追踪每个基元在各层级中的不透明度状态。具体而言，在训练第~$i$~级后，$O_i$~标记所有不透明度低于阈值（$\tau = 0.005$，沿用原始~3DGS~的剪枝设置）的基元为``无效''。这些基元在流式传输和渲染时被跳过，直接减少传输数据量和渲染计算量。与~L3GS~中基于~GSS~评分的静态剪枝不同，LapisGS~的占用图是动态的——低层基元可能因不透明度优化而被标记为无效，从而在更高层级中不再传输和渲染。此外，为实现不同质量层级之间的无缝视觉过渡，LapisGS~在相邻层级之间对基元不透明度进行线性插值：对于目标分辨率~$r_t$~落在~$L_i$~和~$L_{i+1}$~之间的情形，定义连续插值因子~$t(r_t) = (r_t - r_i)/(r_{i+1} - r_i)$，通过插值不透明度实现渐进式细节融合。这与~L3GS~的离散层级切换形成对比——L3GS~的调度器以整层为粒度下载基元，层与层之间的视觉质量是阶梯式跃升的。

\textbf{实验评估与效率分析。}
LapisGS~在~4~个数据集（Synthetic Blender、Mip-NeRF 360、Tanks and Temples、Deep Blending）上评估了~19~个场景，将每个场景的图像下采样~$2\times$、$4\times$、$8\times$~构建~4~级分层表示。与四种基线对比——单尺度模型（Single）、独立多尺度模型（Multiscale，作为质量上界）、基于显著性评分的下采样方法（Downsample）——的实验结果表明：（1）在低分辨率级别上，LapisGS~以原始模型仅~23\%~的大小实现了~SSIM~提升~50.71\%、LPIPS~改善~286.53\%，证明了多层表示在大幅压缩体积的同时能够显著提升视觉保真度；（2）与多尺度模型（每级独立训练，质量上界）相比，LapisGS~在视觉质量上相当，但模型体积最多减少~$2.33\times$，验证了累积结构的跨层特征共享效率；（3）消融实验表明，若移除动态不透明度优化（即完全冻结前层参数），模型体积显著增大且视觉质量下降——因为低层基元无法被``软剪枝''，需要更多新基元来补偿未优化的低层贡献，导致冗余。综合来看，LapisGS~在~L3GS~的基础上，通过动态不透明度优化实现了分层表示的自适应紧缩，使得同一分层框架下的流式传输更加高效。

\subsubsection{动态场景的压缩与流式传输}

本节介绍面向~4D~动态场景的压缩与流式传输方法，包括~4DGS-1K~的时域冗余消除、PD-4DGS~的带宽自适应流式传输和~SwinGS~的任意长体积视频流式传输。

\paragraph{时域冗余消除的超高速动态渲染}

4DGS-1K（2025.03）~\cite{2503_16422}~是面向~4D~动态场景的一套压缩技术套件，其核心思路是：动态场景中绝大部分高斯要么寿命极短、要么在某帧根本不参与渲染——消除这两类时域冗余即可在几乎不损失画质的前提下实现存储和速度的大幅优化。

\smallskip\noindent\textbf{Baseline 的时域冗余。}
在标准~4DGS~中，场景的动态部分由大量~4D~高斯表示，每个~4D~高斯有一个时域参数~$\tau_t$~控制其寿命。作者对~Sear Steak~场景的分析揭示了两个突出问题。
其一，\textbf{短寿命高斯}泛滥：约~70\%~的高斯~$\tau_t < 0.25$，它们在时间维度上闪烁式地出现又消失，尤其集中在快速运动物体的边缘。这些高斯消耗了大量存储却仅在极短的时间窗口内有贡献。
其二，\textbf{不活跃高斯}浪费算力：在任意给定帧，约~85\%~的高斯完全不参与最终渲染，但由于~4D~高斯在光栅化前必须先分解为~3D+1D~分量，所有高斯（包括不活跃的）都会被处理，造成巨大的计算浪费。值得注意的是，活跃高斯集合在相邻~20~帧间的~IoU~超过~80\%，即帧间重叠度很高，这为后续的帧间掩码共享提供了基础。

\smallskip\noindent\textbf{全局剪枝：时空变化分数。}
为消除短寿命高斯，4DGS-1K~引入了\textbf{时空变化分数}（Spatial-Temporal Variation Score）作为新的剪枝准则，综合衡量每个~4D~高斯在空间和时间两个维度上的重要性。

空间维度上，\textbf{空间贡献分数}~$S_i^S$~借鉴~3DGS~中~$\alpha$-混合~的光线贡献累积方式，遍历所有像素计算高斯~$g_i$~对最终颜色的贡献：

\begin{equation}
S_i^S = \sum_{k=1}^{N \cdot H \cdot W} \alpha_i \prod_{j=1}^{i-1}(1-\alpha_j),
\end{equation}

其中~$\alpha_i \prod_{j=1}^{i-1}(1-\alpha_j)$~精确反映了第~$i$~个高斯经由透明度合成对像素颜色的实际贡献。

时间维度上，\textbf{时域变化分数}~$S_i^T$~利用时域不透明度函数~$p_i(t)$~的二阶导数来度量高斯寿命的稳定性：

\begin{equation}
p_i^{(2)}(t) = \bigl((t-\mu_t)^2 / \tau_t^2 - 1\bigr) \cdot \frac{1}{\tau_t} \cdot p_i(t).
\end{equation}

二阶导数幅值大意味着不透明度剧烈波动（短寿命），幅值小则表示平滑持久（长寿命）。将二阶导数经~$\sigma$~函数映射到~$(0,1)$~并在时间轴上取平均，再乘以归一化的~4D~高斯体积，得到最终的时域分数~$S_i^T = S_i^{TV} \cdot \mathrm{Norm}(V(S^{\mathrm{4D}}))$。最终，时空变化分数~$S_i = \sum_{t=0}^{T} S_i^T \cdot S_i^S$~在所有时间步上聚合空间与时域贡献。

所有~4D~高斯按~$S_i$~排序，分数低的被剪枝。在~N3V~数据集上剪枝~80\%，在~D-NeRF~上剪枝~85\%，剪枝后再微调~5{,}000~轮以补偿精度损失。单独的剪枝步骤即可实现~5$\times$~存储压缩和~5$\times$~光栅化加速，且渲染质量略有提升。

\smallskip\noindent\textbf{局部过滤：关键帧时域滤波。}
剪枝消除了短寿命高斯，但剩余高斯中仍有大量在某帧不活跃的冗余。为此，4DGS-1K~设计了\textbf{关键帧时域滤波}，利用活跃高斯的帧间高重叠度来实现掩码共享。

具体而言，以固定间隔~$\Delta_t$~选取关键帧~$\{t_i\}_{i=0}^{T}$，对每个关键帧从所有训练视角渲染并计算活跃高斯的可见性掩码~$\{m_{i,j}\}_{j=1}^{N}$（$N$~为训练视角数）。相邻关键帧之间的非关键帧直接复用最近两个关键帧~$t_l, t_r$~的掩码并集~$\bigcup_{j=1}^{N} m_{l,j} \cup \bigcup_{j=1}^{N} m_{r,j}$，仅对掩码标记为活跃的高斯执行光栅化。在~N3V~上~$\Delta_t = 20$~帧，在~D-NeRF~上选取~6~个关键帧。剪枝使高斯寿命普遍延长，从而允许使用较大的关键帧间隔而不遗漏重要高斯；之后同样微调以补偿滤波引入的精度损失。

\smallskip\noindent\textbf{整体效果。}
两个步骤协同作用：全局剪枝从根源消除短寿命冗余，局部过滤屏蔽逐帧不活跃冗余。在~N3V~数据集上，4DGS-1K~相比原始~4DGS~实现~\textbf{41$\times$~存储压缩}（含向量量化和掩码位压缩后处理）和~\textbf{9$\times$~光栅化加速}，在现代~GPU~（RTX 3090）上达到~\textbf{1000+~FPS}，在老旧~TITAN X~上仍维持~200+~FPS（原始~4DGS~为~20~FPS），PSNR~仅下降~0.04~dB。4DGS-1K~在渲染速度上达到同类~4D~压缩方法（MEGA、Compact3D）的~4--10$\times$，且仅消耗动态~3DGS~约~1\%~的存储空间即获得更高的渲染质量。

\paragraph{带宽自适应的动态场景流式传输}

4DGS-1K~聚焦存储压缩与渲染速度优化，解决的是``静态比特流够不够小、渲染够不够快''的问题；而实际部署中同样关键的是另一个维度——\textbf{渐进式传输}：能否让用户在下载了比特流的前~10\%~时即看到一个粗略但合理的画面，随下载推进逐步精细？PD-4DGS（2026.05）~\cite{2605_11427}~正是首个面向~4DGS~的渐进式压缩与按需传输框架，将单次训练产出的~4DGS~模型转化为~DASH/HLS~兼容的、可从任意前缀渲染的分层比特流。

\smallskip\noindent\textbf{分层变形分解（HDD）。}
现有~4DGS~架构（如 Compact3DGS、Efficient4DGS）仅在归纳偏置层面使用粗到细层次——训练时内部网络隐式地从粗到精逐步收敛，但最终产物仍被打包为一个整体式比特流，无法按需截断。HDD~将这一隐式层次\textbf{外化}为比特流层面的三个可独立传输的层：

\begin{itemize}
\item \textbf{Layer~0（静态脚手架）}：仅包含规范锚点集~$A = \{(x_v, f_v, s_v, o_v)\}_{v=1}^{V}$~（位置、特征、缩放、偏移）和属性预测~MLP。客户端收到此层即可渲染场景的静态快照，无需任何时域输入。由于锚点是稀疏的，Layer~0~约占~60\%~的总模型体积，但其稀疏性直接决定了后续层的质量上限。
\item \textbf{Layer~1（全局变形）}：增量传输轻量级全局变形网络~$F_{\mathrm{global}}$~（HexPlane~+~MLP），对每个锚点预测位移和特征残差，将静态场景提升为动态模式。由于变形仅作用于稀疏锚点（而非逐高斯），Layer~1~通常仅占总模型体积不到~10\%，却已能捕捉场景中主导性的大幅运动（如人物走动、相机移动）。
\item \textbf{Layer~2（局部精修）}：传输局部变形网络~$F_{\mathrm{local}}$，在粗粒度动态高斯之上预测逐高斯残差~$(\Delta\mu, \Delta\Sigma, \Delta\alpha, \Delta c)$，恢复高频细节（头发飘动、面部微表情等），达到完整保真度。
\end{itemize}

\noindent 前向过程根据已接收的前缀层数~$s \in \{0, 1, 2\}$~路由到不同计算图：$s=0$~时~$R(\Pi(A))$~仅渲染静态脚手架；$s=1$~时~$R(\Pi(A + F_{\mathrm{global}}(A, t)))$~加入全局动态；$s=2$~时~$R(\Pi((A + F_{\mathrm{global}}(A, t)) + F_{\mathrm{local}}(G, t)))$~达到完整动态渲染。三层共享相同的锚点和~MLP~基础设施，切换层次仅是增量加载权重，无需重新解码或重新训练，因此天然适配带宽波动和异构设备部署。

\smallskip\noindent\textbf{属性级率失真优化（R-DO）。}
Layer~0~占~60\%~模型体积，若基底层过大则秒级首帧的承诺无法兑现。PD-4DGS~没有采用启发式剪枝阈值（如 LightGaussian、Compressed3DGS），而是将保留/丢弃建模为\textbf{可微率失真}问题。为每个锚点~$i$~引入可学习掩码~$m_i \in [0,1]$，在训练和推理时对称地乘以高斯的不透明度和缩放：$\alpha'_i = \alpha_i m_i$，$s'_i = s_i m_i$。这种双乘法门控确保锚点随~$m_i \to 0$~同时变得透明且体积为零，消除了``软训练/硬剪枝''分布不匹配问题。

对于连续属性~$a \in \{f_{v,i}, s_{v,i}, o_{v,i}\}$，假设逐批次的~Gaussian~先验~$\mathcal{N}(\mu_a, \sigma_a^2)$，得到~Shannon~比特代价~$R(a) = -\log_2 \Phi\bigl(\frac{a + Q_a/2 - \mu_a}{\sigma_a}\bigr) + \log_2 \Phi\bigl(\frac{a - Q_a/2 - \mu_a}{\sigma_a}\bigr)$，其中~$Q_a$~为量化步长，$\Phi$~为标准正态~CDF（Cumulative Distribution Function，累积分布函数）。与~$\ell_1/\ell_2$~正则化不同，该代价项直接逼近熵编码器（导出时为~LZMA）的输出长度，将训练时的压缩信号与实际比特流大小紧密耦合。

总损失为：

\begin{equation}
\mathcal{L}_{\mathrm{total}} = \mathcal{L}_{\mathrm{render}} + \lambda_{\mathrm{layer}} \cdot \frac{1}{|V|} \sum_{i=1}^{V} m_i \bigl(R(f_{v,i}) + R(s_{v,i}) + R(o_{v,i})\bigr),
\end{equation}

其中~$\mathcal{L}_{\mathrm{render}}$~为~$\ell_1$~+~D-SSIM~损失。仅活跃锚点（$m_i > 0$）产生比特预算，优化器得到明确的保留/丢弃信号。三层使用量级递减的~$\lambda_{\mathrm{layer}} \in \{0.04, 0.01, 0.00025\}$，使~Layer~0~被激进压缩而~Layer~2~几乎保留所有锚点以追求最高保真度。

\smallskip\noindent\textbf{时域掩码一致性（TMC）。}
R-DO~稀疏化会导致掩码沿时域轴高频振荡——产生闪烁和鬼影拖尾，而~$\ell_1$~稀疏损失无法捕获此类时域伪影。TMC~通过两个互补的正则项稳定激活模式：其一，\textbf{二元熵惩罚}~$\mathcal{L}_{\mathrm{binary}} = -\frac{1}{|V|}\sum_{i=1}^{V} \bigl(m_i \log_2 m_i + (1 - m_i)\log_2(1 - m_i)\bigr)$~驱动掩码值确定性趋近~$\{0, 1\}$，消除半透明鬼影拖尾并最小化熵编码长度；其二，\textbf{空间平滑约束}~$\mathcal{L}_{\mathrm{smooth}} = \frac{1}{|S|}\sum_{(i,j)\in S} w_{ij} |m_i - m_j|$~利用空间相邻锚点共享激活状态的先验（$w_{ij} = \exp(-\|x_i - x_j\|/\sigma)$），在不直接建模时域轴的情况下隐式强制跨帧一致性。

\smallskip\noindent\textbf{容量加权渐进训练。}
架构和损失本身并不保证层间单调精化——训练时的层采样策略同样关键。若三层均匀采样~$\pi = (1/3, 1/3, 1/3)$，则高容量的~GAD/LGD~网络（参数量为基底的~2--6~倍）仅被~$1/3$~的反向传播覆盖，导致深层不优于基底：在~Dycheck~spin~场景上观察到~$\text{PSNR}_0 / \text{PSNR}_1 / \text{PSNR}_2 = 14.61/14.54/14.54$，即渐进性崩溃。

PD-4DGS~利用训练过程中掩码本身携带的场景需求信息，定义\textbf{层间掩码激活率}~$\rho = \mathbb{E}_{i \in V}[m_i^{(\ell=2)} - m_i^{(\ell=0)}] \in [0,1]$~作为场景动态程度的完全可学习统计量，并以此为插值系数在均匀分布~$\pi_{\mathrm{uniform}} = (1/3, 1/3, 1/3)$~和激进分布~$\pi_{\mathrm{aggr}} = (0.15, 0.30, 0.55)$~之间自适应混合：$\pi(\rho) = (1-\rho)\,\pi_{\mathrm{uniform}} + \rho\,\pi_{\mathrm{aggr}}$。静态场景（$\rho \to 0$）趋向均匀分布以确保基底充分训练；高动态场景（$\rho \to 1$）趋向激进分布使深层网络获得更多梯度信号。$\rho_t$~每~$K=200$~次迭代采样一次，通过~EMA（$\beta=0.05$）平滑早期噪声，前~$T_{\mathrm{warm}}=2000$~次迭代固定使用~$\pi_{\mathrm{uniform}}$~预热。所有超参数（$\beta$、$K$、$T_{\mathrm{warm}}$、$\pi_{\mathrm{uniform}}$、$\pi_{\mathrm{aggr}}$）在所有场景上完全一致——场景间差异完全由训练得到的~$\rho$~吸收。

\smallskip\noindent\textbf{整体效果。}
在~Dycheck iPhone~基准上，PD-4DGS~在匹配渲染保真度下将流式比特流减少~\textbf{60\%~以上}，在~2~Mbps~链路上将首帧延迟从~\textbf{73--930~秒降至约~1.7~秒}。与~4DGS-1K~的``全量下载后快速渲染''路线互补，PD-4DGS~走的是``渐进下载、即时渲染''路线——用户无需等待完整比特流下载即可开始观看粗略动态场景，并在网络条件允许时自动精化至完整保真度。两个方法一个优化静态压缩率、一个优化渐进传输体验，共同构成动态~3DGS~流式传输的技术版图。

\paragraph{任意长度的体积视频流式传输}

与~4DGS-1K~聚焦于单段动态场景的压缩效率、PD-4DGS~解决渐进式按需传输不同，SwinGS（2024.09）~\cite{2409_07759}~面向的是一个更根本的瓶颈：\textbf{视频长度限制}。现有~4DGS~流式传输方法（包括上述两者）均假设场景被预先训练为一段固定长度的~4D~高斯表示，一旦视频时长超出~GPU~显存或浏览器~RAM~的承载能力，训练和渲染都会崩溃。对于体量动辄数十分乃至数小时的体积视频内容而言，这一限制是实际部署的核心障碍。

\smallskip\noindent\textbf{滑动窗口 + MCMC 采样。}
SwinGS~的核心思路是将无限长视频拆分为有限长度的滑动窗口，每个窗口内独立维护一组~3D~高斯快照。具体而言，采用\textbf{时空高斯}（spatiotemporal Gaussian）表示每个窗口内的动态内容，并通过~\textbf{Markov Chain Monte Carlo（MCMC）采样}替代原始~3DGS~中的自适应密度控制：MCMC~采样以概率方式决定每帧高斯的增删与属性更新，无需存储历史状态即可在窗口内逐步累积高斯分布。窗口滑动时，旧高斯按衰减概率被丢弃，新高斯从当前帧的采样过程中生成，从而实现\textbf{无状态增量式}的场景表示更新——每个新帧只需传输该帧新增或修改的高斯，而非整个场景的全量参数。

\smallskip\noindent\textbf{WebGL 实时查看器与跨设备支持。}
在传输端，SwinGS~将每帧的增量高斯参数编码为轻量级数据包，实现\textbf{83.6\%~的传输成本降低}。在渲染端，SwinGS~实现了一个基于~WebGL~的交互式查看器，可直接在浏览器中解码增量数据包并实时渲染。该查看器无需安装任何插件或本地运行时，在现代浏览器中即可运行，包括智能手机和平板等移动设备上的~Safari~和~Chrome。这意味着用户在手机上打开一个网页链接即可观看和交互体积视频——这是将~3DGS~从实验室~GPU~推向大众消费级设备的关键一步。

\smallskip\noindent\textbf{整体效果与定位。}
SwinGS~在传输效率、渲染速度和跨设备可达性三个维度上建立了体积视频流式传输的新基准。与~4DGS-1K~通过剪枝和滤波压缩单段视频、PD-4DGS~通过分层编码加速渐进传输的定位不同，SwinGS~解决了\textbf{任意长度}体积视频的持续流式传输问题——它不追求单段场景的最优压缩比，而是通过滑动窗口机制使渲染管道的内存消耗与视频长度解耦，从而将~3DGS~流式传输从"短视频"扩展到"无限流"的范式。

\begin{takeaway}
视频与流式渲染将3DGS从离线静态场景推向实时动态视频流。SwinGS 的滑动窗口加 MCMC 采样实现了任意长度体积视频的无状态增量式传输，传输成本降低 83.6\%，WebGL 查看器使用户在手机浏览器中即可观看体积视频。这标志着 3DGS 从实验室 GPU 走向大众消费级设备的关键一步。
\end{takeaway}


%=====================================================================
\section{结构化与锚点表示}
\label{sec:11_5}

\subsection{问题背景}

本节讨论两条技术路线：（1）锚点驱动的层级表示；（2）混合与层次化表示。

原始~3DGS~中的每个高斯都是独立的：它独立存储~59~维参数，独立参与渲染，不与其他高斯共享任何信息。这种"全独立"表示虽然灵活，但导致了大量冗余。相邻高斯往往具有相似的颜色、尺寸和方向，却各自存储了一份完整的参数副本。结构化与锚点表示的核心思想是引入层级组织，用一组稀疏的锚点（anchor）代表局部区域。锚点通过小型网络预测其周围高斯的属性，从而实现参数共享和冗余消除。

\subsubsection{锚点驱动的层级表示}

本节介绍以锚点为中心组织高斯基元的层级表示方法，从奠基性的 Scaffold-GS 到二阶增强的 SOGS，展示锚点驱动范式如何从根本上减少冗余基元的产生。

\paragraph{锚点驱动的视图自适应渲染}

Scaffold-GS（2023.12）~\cite{2312_00109}~是锚点表示的奠基性工作，也是后续~HAC++、ContextGS~等熵编码方法的基础架构。其核心思想是将场景组织为由锚点定义的局部高斯簇。每个锚点存储一个紧凑的特征向量，根据当前视角和距离动态预测其管辖范围内子高斯的完整属性，而非像原始~3DGS~那样为每个高斯静态存储所有参数。

直觉地说，Scaffold-GS~将~3DGS~从"每个高斯独立行动"的组织模式转变为"层级管理"模式。在原始~3DGS~中，每个高斯独立存储~59~维参数（位置~3~维 + 四元数~4~维 + 缩放~3~维 + 不透明度~1~维 + 球谐系数~48~维），相邻高斯往往具有相似的颜色和几何，却各自保存一份完整副本，这正是冗余的根源。锚点则扮演管理者的角色：它知道"我的辖区内应该有什么样的高斯"，但具体的属性值根据当前视角动态生成，而非静态存储。这种设计带来三个优势：（1）\textbf{参数共享}，同一锚点下的~$k$~个子高斯共享锚点特征，无需重复存储相似属性；（2）\textbf{视图自适应}，高斯属性根据视角动态调整，近处取高分辨率细节、远处取低分辨率概要，自然支持多尺度渲染；（3）\textbf{冗余控制}，锚点的增删策略比逐高斯的~ADC~更稳定，从结构层面抑制过重建。

\textbf{锚点初始化与视图相关特征融合。}
Scaffold-GS~的锚点并非手动指定，而是从~SfM~稀疏点云自动生成。具体而言，将~SfM~点云体素化（体素大小为~$v_s$），取每个非空体素的中心作为锚点位置。每个锚点配备一个~32~维局部上下文特征~$\mathbf{f}_v$、一个~3~维缩放因子~$\mathbf{l}_v$~和~$k$~个可学习偏移量~$\{\mathbf{O}_i\}_{i=0}^{k-1} \in \mathbb{R}^{k \times 3}$（默认~$k=10$）。为了赋予锚点多尺度感知能力，$\mathbf{f}_v$~被扩展为一个三级特征库~$\{\mathbf{f}_v, \mathbf{f}_{v/2}, \mathbf{f}_{v/4}\}$，分别对应原始分辨率和~$2\times$、$4\times$~下采样。渲染时，根据相机~$\mathbf{x}_c$~与锚点~$\mathbf{x}_v$~之间的相对距离~$d_{vc}$~和视角方向~$\boldsymbol{\delta}_{vc}$，通过一个小型权重网络~$F_w$~预测融合权重，动态加权求和生成视图相关的锚点特征：

\begin{equation}
\tilde{\mathbf{f}}_v = w_0 \, \mathbf{f}_v + w_1 \, \mathbf{f}_{v/2} + w_2 \, \mathbf{f}_{v/4}, \quad [w_0, w_1, w_2] = F_w(d_{vc}, \boldsymbol{\delta}_{vc}),
\end{equation}

其中~$d_{vc} = \|\mathbf{x}_v - \mathbf{x}_c\|_2$，$\boldsymbol{\delta}_{vc} = (\mathbf{x}_v - \mathbf{x}_c) / \|\mathbf{x}_v - \mathbf{x}_c\|_2$。这一机制使得同一锚点在不同观察距离下自动选择合适的特征分辨率，为后续的属性解码提供视角感知的输入。

\textbf{神经高斯属性解码。}
对于视锥内每个可见锚点，Scaffold-GS~生成~$k$~个（默认~$k=10$）神经高斯。子高斯的位置由锚点位置加偏移量计算：$\boldsymbol{\mu}_i = \mathbf{x}_v + \mathbf{O}_i \cdot \mathbf{l}_v$。其余属性——不透明度~$\boldsymbol{\alpha}$、颜色~$\mathbf{c}$、旋转四元数~$\mathbf{q}$、缩放~$\mathbf{s}$——均由独立的~2~层~MLP（隐藏层~32~维，ReLU~激活）从视图相关锚点特征~$\tilde{\mathbf{f}}_v$、相对距离~$d_{vc}$~和视角方向~$\boldsymbol{\delta}_{vc}$~一次性解码：

\begin{equation}
[\alpha_0, \ldots, \alpha_{k-1}] = F_\alpha(\tilde{\mathbf{f}}_v, \, d_{vc}, \, \boldsymbol{\delta}_{vc}), \quad \mathbf{c}_i, \, \mathbf{q}_i, \, \mathbf{s}_i \text{ 由各自的 } F_c, F_q, F_s \text{ 类似解码。}
\end{equation}

这种"按需生成"策略意味着只有视锥内可见的锚点才会激活~MLP~计算，且不透明度低于阈值~$\tau_\alpha$~的神经高斯直接被过滤，大幅降低了无效计算。值得注意的是，论文实验表明无论~$k$~取何值，训练后最终激活的神经高斯数量趋同——模型自动学会选择一组充分且非冗余的高斯集合。

\textbf{锚点增长与剪枝。}
由于初始~SfM~点云在纹理匮乏区域和观测不足区域可能极度稀疏，Scaffold-GS~设计了基于梯度的增长策略来弥补锚点覆盖不足。具体而言，将神经高斯按体素大小~$v_g$~空间量化（默认~$v_g = 64$），在每个体素内累积其包含的神经高斯在~$N=100$~次迭代中的梯度均值~$\bar{g}$，若~$\bar{g} > \tau_g$~则在该体素中心部署新锚点。为支持多粒度增长，体素大小和阈值在不同量化层级递进：$v_g^{(m)} = v_g / 4^{m-1}$，$\tau_g^{(m)} = \tau_g \cdot 2^{m-1}$，其中~$m$~为量化层级索引。此外，新锚点候选还需通过随机淘汰机制进一步控制扩张速度。剪枝则更为直接：累积每个锚点的神经高斯不透明度，若低于~$0.5$~则移除该锚点。

\textbf{训练目标与收敛特性。}
总体损失函数结合了~$L_1$~颜色重建损失、SSIM~结构相似性损失和体积正则化：

\begin{equation}
\mathcal{L} = \mathcal{L}_1 + \lambda_{\text{SSIM}} \mathcal{L}_{\text{SSIM}} + \lambda_{\text{vol}} \mathcal{L}_{\text{vol}}, \quad \mathcal{L}_{\text{vol}} = \sum_{i=1}^{N_{ng}} \text{prod}(\mathbf{s}_i),
\end{equation}

其中~$\lambda_{\text{SSIM}} = 0.2$，$\lambda_{\text{vol}} = 0.001$。体积正则化~$\mathcal{L}_{\text{vol}}$~鼓励神经高斯保持小体积、减少重叠，隐式抑制冗余。训练共~30k~次迭代，收敛速度快于原始~3DGS。

在~Mip-NeRF 360~数据集上，Scaffold-GS~以更少的高斯数量和更小的存储空间达到了与原始~3DGS~相当的渲染质量（PSNR~持平），同时在~Tanks and Temples~和~Deep Blending~等包含纹理匮乏区域和镜面反射的复杂场景上超越了~3DGS。在~BungeeNeRF~等多尺度数据集上，Scaffold-GS~在~3DGS~出现模糊和针状伪影的区域展现了明显优势，验证了视图自适应属性预测对新视角的泛化能力。

由于~Scaffold-GS~建立了结构化的锚点-高斯层级关系，它成为后续大量压缩工作的天然基础。HAC++、ContextGS、CAT-3DGS~等熵编码方法均基于~Scaffold-GS~的锚点结构构建上下文模型（详见第~\ref{sec:11_2}~节）。

\paragraph{二阶锚点增强}

Scaffold-GS~的锚点特征维度~$D$~面临一个根本矛盾：增大~$D$~可以提升渲染质量，但每个锚点存储~$D$~维浮点向量直接推高模型体积；减小~$D$~则导致~MLP~解码器的输入信息不足，子高斯的颜色、几何等属性预测精度下降。SOGS（2025.03）~\cite{2503_07476}~的核心洞察是——与其堆叠更多的一阶特征维度，不如挖掘\emph{特征维度之间的协变异关系}作为额外的信息来源。

\textbf{二阶锚点：从协方差到特征增强。}
与~Scaffold-GS~的一阶锚点特征不同，SOGS 引入\textbf{二阶锚点}（second-order anchor）。其思路借鉴了信号处理中功率谱密度的思想：二阶统计量虽然不直接保留原始信号，但能紧凑地刻画分布形状。具体而言，给定全部~$N$~个锚点的特征集~$\mathbf{F}^a = [\mathbf{f}_1^a, \ldots, \mathbf{f}_N^a] \in \mathbb{R}^{N \times D}$，SOGS 将~$D$~个特征维度视为~$D$~个随机变量，$N$~个锚点视为观测样本，计算~$D \times D$~的协方差矩阵~$\boldsymbol{\Sigma}$：

\begin{equation}
\boldsymbol{\Sigma} = \frac{1}{N-1}(\mathbf{F}^a - \bar{\mathbf{F}}^T)(\mathbf{F}^a - \bar{\mathbf{F}}^T)^T, \quad \bar{\mathbf{F}} = \frac{1}{N}\mathbf{F}^a,
\end{equation}

其中~$\bar{\mathbf{F}} \in \mathbb{R}^{N \times D}$~为各维度均值。然而，两个方差较大的维度即使仅有弱相关性也可能产生大的协方差值，因此 SOGS 进一步将协方差矩阵标准化为相关矩阵~$\mathbf{R} = \mathbf{A}^{-1} \boldsymbol{\Sigma} \mathbf{A}^{-1}$，其中~$\mathbf{A} = \text{diag}(\sigma_1, \ldots, \sigma_D)$~为标准差对角矩阵，消除不同维度的尺度差异。

对~$\mathbf{R}$~进行特征分解~$\mathbf{R} = \mathbf{Q}\boldsymbol{\Lambda}\mathbf{Q}^T$，取前~$M$~个最大特征值对应的特征向量~$\mathbf{P} = [\mathbf{P}_1, \ldots, \mathbf{P}_M]$~作为主协变异方向——每个特征向量捕捉一组跨特征维度的共变模式。这些协变异方向是全局共享的（从全部锚点统计得到），但与每个锚点的局部特征~$\mathbf{f}^a$~结合后可以提取锚点特有的纹理-结构信息：

\begin{equation}
\mathbf{f}_i^t = F_i([\mathbf{P}_i, \, \mathbf{f}^a]), \quad i \in [1, M],
\end{equation}

其中~$F_i$~为两层~MLP。最终，提取的~$\{\mathbf{f}_i^t\}_{i=1}^{M}$~与原始锚点特征~$\mathbf{f}^a$~拼接后送入属性解码~MLP：

\begin{equation}
[\alpha_k, \, \mathbf{c}_k, \, \mathbf{q}_k, \, \mathbf{s}_k]_{k=0}^{K} = F(\mathbf{f}^a, \, \{\mathbf{f}_i^t\}_{i=0}^{M}, \, d_{ac}, \, \mathbf{d}_{ac}).
\end{equation}

关键在于：增强特征~$\{\mathbf{f}_i^t\}$~是从~$\mathbf{f}^a$~自身派生而来，\emph{不增加每个锚点的存储维度}——协方差和相关矩阵是在训练过程中动态计算的共享统计量，无需逐锚点保存。这使得 SOGS 可以用更小的~$D$（即更紧凑的锚点特征维度）实现比~Scaffold-GS~更高维特征才能达到的渲染质量。

\textbf{选择性梯度损失。}
为配合紧凑的锚点特征，SOGS 进一步引入选择性梯度损失（selective gradient loss）来增强纹理和几何细节的优化。标准的像素级~$L_1$~损失对纹理和结构的敏感度有限——平坦区域在损失中占主导，边缘和纹理等关键区域反而被淹没。SOGS 首先对渲染图像和真值图像分别用~Sobel~算子提取水平和垂直梯度图~$(G_x', G_y')$~和~$(G_x, G_y)$，计算方向性梯度差异~$l_x$~和~$l_y$。然后，以渲染误差图的绝对差异作为权重图~$w_x = |G_x' - G_x|$、$w_y = |G_y' - G_y|$，对梯度差异进行加权：

\begin{equation}
\mathcal{L}_s = w_x \odot l_x + w_y \odot l_y,
\end{equation}

其中~$\odot$~表示逐元素乘法。这一机制使得模型在训练过程中自适应地将优化注意力聚焦于难以渲染的纹理和结构区域，且聚焦区域随渲染误差图的演化动态调整。总体损失为：

\begin{equation}
\mathcal{L} = \lambda_1 \mathcal{L}_1 + \lambda_{\text{SSIM}} \mathcal{L}_{\text{SSIM}} + \lambda_{\text{vol}} \mathcal{L}_{\text{vol}} + \lambda_s \mathcal{L}_s.
\end{equation}

实验表明，SOGS~以更小的锚点特征维度和更少的模型尺寸，在~Mip-NeRF 360~等基准上实现了优于~Scaffold-GS~的渲染质量（PSNR~提升约~0.3--0.5~dB），验证了二阶统计量在不增加存储开销的前提下有效增强锚点特征表达能力的设计理念。

\subsubsection{混合与层次化表示}

本节介绍将锚点范式与混合模型或层次结构相结合的方法，包括~GaussianForest~的层次-混合表示、ProtoGS~的原型驱动渲染和~Sketch \& Patch++~的语义分离混合表示。

\paragraph{层次-混合表示}

Scaffold-GS~的锚点-高斯层级结构虽然实现了参数共享，但存在一个根本性的张力：大锚点特征意味着更强的表达能力，却导致模型膨胀；缩小特征维度可以压缩模型，但属性预测精度随之下降。这种"表达能力与存储开销的矛盾"在复杂场景中尤为突出——成千上万的相邻高斯彼此高度相似，却各自独立存储~59~个参数，造成大量冗余。GaussianForest（2024.06）~\cite{2406_08759}~的核心洞察是：高斯属性可以分为两类——\textbf{显式属性}（位置~$\bm{\mu}$~和不透明度~$\alpha$）对偏差高度敏感、不随视角变化且邻居间差异显著；\textbf{隐式属性}（协方差和视角相关颜色）参数量大但局部平滑、邻居间高度冗余。如果将显式属性逐高斯独立存储以保证精度，将隐式属性通过层次结构共享以消除冗余，就能同时兼顾两者优势。

GaussianForest~将场景建模为"混合高斯森林"——一组层次化组织的混合高斯树。每棵树由三层构成：根节点层~$T_R$、内部节点层~$T_I$、叶节点层~$T_L$。叶节点存储显式属性~$T_L = \{\bm{\mu}, s_s, \alpha, p\} \in \mathbb{R}^6$（仅~6~个参数，其中~$s_s$~是缩放系数，$p$~是指向父节点的指针），根节点和内部节点存储隐式特征向量。每个叶节点通过指针沿~$T_L \to T_I \to T_R$~路径向上遍历，拼接路径上的隐式特征~$\bm{f} = [\bm{f}_I, \bm{f}_R]$，再由两个轻量~MLP~解码器~$F_{\text{cov}}$~和~$F_{\text{rgb}}$~分别解码出协方差参数（旋转~$\bm{q}$~和缩放~$\bm{s}$）与视角相关颜色~$\bm{c}$。这种设计使得兄弟叶节点（共享同一父节点的叶节点）自动共享隐式属性，而每个叶节点仍保留自己独立的位置和不透明度。

直觉地说，这种设计就像一棵倒挂的族谱：叶节点是"个人"，每个人有自己独立的位置和身份；内部节点和根节点是"家族"和"宗族"，共享的几何和颜色特征沿层级传递，越往上共享范围越大。论文中的可视化生动印证了这一点——一棵高斯树可以有效地建模鼓面的完整表面，而另一棵树表示镲片的一部分。原本需要数千个独立高斯（每个~59~个参数）才能描述的鼓面，在~GaussianForest~中仅需一个根节点、数十个内部节点和轻量叶节点（每个仅~6~个参数）即可完成。

森林的自适应增长和剪枝策略控制节点总数。增长策略基于叶节点的累积梯度~$CG$——梯度高意味着局部区域复杂、当前表示不够充分。根据三个梯度阈值~$\{\tau_2, \tau_1, \tau_0\}$（$\tau_2 < \tau_1 < \tau_0$），增长分为三档：仅克隆叶节点、克隆叶节点及其父内部节点、沿完整路径克隆全部节点。这种层次化增长确保简单区域（低梯度）保持紧凑，复杂区域（如物体边界）获得更丰富的特征维度。剪枝策略则移除缩放或透明度低于阈值的叶节点，并清理无子节点的冗余内部节点。增长在训练早期逐步停止（各层有预设停止时间~$t_l = \{5k, 10k, 15k\}$），剪枝则持续到训练结束。

从空间复杂度分析，原始~3DGS~为~$O(59N)$，而~GaussianForest~的叶节点仅~$O(6N)$，非叶节点约~$O(2.2N)$~至~$O(4.3N)$，加上剪枝可将叶节点减少~1.5--3~倍，总体压缩因子约为~8--17$\times$。在~Small~设置下（$D_R=24, D_I=16$），特征维度极低；在~Large~设置下（$D_R=32, D_I=24$），特征维度更大但共享效果更强。实验中，GaussianForest~在~21~个场景（包括~Mip-NeRF 360、Tanks and Temples、Deep Blending~和~Synthetic Blender~数据集）上实现了最高~717$\times$~的模型大小缩减（相对~3DGS），渲染质量（PSNR/SSIM/LPIPS）保持可比甚至超越，同时因高斯数量大幅减少而加速了训练和渲染速度。这种层次-混合表示也因此展现出无监督~3D~分割的潜力——属于同一父节点的高斯自然聚集为语义一致的局部区域。

\paragraph{原型驱动的高效渲染}

与~GaussianForest~通过层次结构共享隐式属性不同，ProtoGS（2025.03）~\cite{2503_17486}~从聚类角度解决冗余问题。其核心动机源于一个简单观察：场景中相邻区域的同类高斯基元（尤其是纹理均匀区域）往往高度相似，却各自独立存储~59~维参数，造成大量冗余。与其像~GaussianForest~那样设计树形结构让叶节点沿路径共享特征，不如直接将相似基元"合并"为少数代表性原型。每个原型本质上就是一组相似高斯的质心向量，直接承载完整的~3DGS~参数（位置、协方差、颜色、不透明度），而非需要~MLP~解码的隐式特征。

具体而言，ProtoGS~采用三步流水线。第一步是~\textbf{SfM 锚点分片}：从~SfM~点云中随机采样~$M$~个锚点，按欧氏距离将所有高斯基元~$G \in \mathbb{R}^{N \times d}$~分配到最近的锚点所属分片。这种基于空间位置的分组确保每组分片内的基元在几何和纹理上具有局部一致性，避免将空间上不相邻但恰好落入同一网格的基元强行聚类——这是朴素网格划分的主要缺陷。第二步是~\textbf{渲染引导的~K-means~聚类}：在第~$m$~个分片内，对~$n_m$~个基元执行~K-means~聚类，推导出~$K^m$~个原型~$\bm{\mu}_{m,k}$。关键创新在于，聚类目标函数~$J^m = \sum_{i=1}^{n_m} \sum_{k=1}^{K^m} r_{ik}^m \|G_{m,i} - \bm{\mu}_{m,k}\|^2$~并非独立优化，而是与渲染损失~$L = (1-\lambda)\mathcal{L}_1 + \lambda \mathcal{L}_{\text{D-SSIM}} + \lambda_c L_c$~联合优化。这意味着聚类分配~$r_{ik}^m$~的更新受渲染误差的监督——如果某次聚类重分配导致渲染质量下降，梯度信号会修正原型参数和分配关系。相比之下，朴素的"先聚类再优化"两阶段方法因聚类阶段完全忽略渲染信号，容易产生几何和纹理受损的原型，后续优化也难以恢复。第三步是~\textbf{渐进衰减}：在训练过程中以~0.5~的衰减率逐步减半原型数量，实现自适应压缩。

这一设计的核心优势在于：原型本身就是标准~3DGS~参数向量，可以直接输入~3DGS~的光栅化渲染管线，无需任何~MLP~解码步骤。这使得~ProtoGS~在渲染效率上显著优于~Scaffold-GS~和~GaussianForest~等混合表示方法——后者需要为每个高斯运行~MLP~前向传播来解码协方差和颜色，引入了不可忽视的计算开销。实验中，ProtoGS~在~Mip-NeRF 360、Tanks and Temples、Deep Blending、DTU~和~Synthetic Blender~等~27~个场景上，以远少于基线方法的原型数量实现了可比甚至更优的渲染质量，存储开销大幅降低的同时保持了~3DGS~级别的高速渲染速度。原型可视化也印证了聚类效果：在纹理均匀区域（如沙发靠垫、墙面），ProtoGS~仅用少数紧凑原型即可覆盖原本需要数百个细碎椭球才能描述的区域。

\paragraph{语义分离的混合表示}

与~GaussianForest~通过层次结构共享隐式属性、ProtoGS~通过聚类合并相似基元不同，Sketch \& Patch++（S\&P++）（2026.01）~\cite{2601_05394}~从另一个维度审视高斯冗余问题：高斯在场景中并非同质的，它们扮演着截然不同的语义角色。有些高斯紧密排列在物体边缘、轮廓和几何不连续处，捕捉高频的边界定义特征；另一些则稀疏分布在平滑表面上，表达低频的体积填充。如果将这两类高斯混为一谈、采用统一的压缩策略，必然顾此失彼：对边界高斯过度压缩会模糊锐利边缘，对体积高斯保留过多则浪费存储。S\&P++~的核心设计正是基于这一观察，将高斯按语义角色分离为\textbf{素描高斯}（SketchGS）和\textbf{色块高斯}（PatchGS），并为每一类定制专用的压缩管线。

\textbf{分类方法：多准则密度聚类~+~自适应多项式精炼。}
S\&P++~首先需要自动识别哪些高斯属于素描、哪些属于色块。其前代工作~S\&P~依赖~Line3D++~提取~3D~线段作为几何先验，仅在人造场景的直线边缘附近识别素描高斯，无法处理曲线边界、有机结构或自然场景。S\&P++~放弃了这一外部先验，转而直接分析高斯自身的内在属性。具体而言，它扩展了~DBSCAN~聚类算法，同时考虑三个互补的距离度量——空间距离（欧氏距离，捕捉几何邻近性）、方向距离（余弦相似度，捕捉朝向一致性）和颜色距离（RGB~差异，捕捉外观相似性），只有三个约束同时满足时两个高斯才被视为邻居。这种多准则聚类从数百万高斯中发现具有多维一致性的紧凑簇——不仅限于线性特征，还能自动发现曲线和复杂结构模式。聚类之后，S\&P++~对每个簇进行\textbf{自适应多项式回归精炼}：为缩放、旋转、不透明度和颜色四类属性分别拟合多项式模型，将位置映射到属性值。回归误差低的簇表明位置与属性之间存在强相关性，这正是素描高斯的标志（边缘处高斯沿几何结构排列，属性可预测）；回归误差高的簇则表明属性变化不规则，属于色块高斯。不满足质量阈值的簇会被递归拆分——利用回归残差构建特征向量，通过~K-means~将簇分裂为子簇，直到每个子簇要么达到足够的多项式拟合质量（归入素描），要么达到最小规模阈值或最大迭代次数（归入色块）。最终，室内场景中高达~88.7\%~的高斯被分类为素描高斯，室外场景约~65.9\%~，远超前代方法基于线段检测的覆盖率。

\textbf{素描高斯编码：多项式参数化~+~级联压缩。}
素描高斯的关键优势在于其属性与位置之间的高度可预测关系。S\&P++~为每个素描簇拟合多项式回归模型，将逐高斯的属性存储转化为少量多项式系数的共享存储。原始~3DGS~中每个高斯需要~59~个参数，而多项式编码后每个素描高斯的属性仅需引用簇级系数即可重建，实现约~200--248$\times$~的素描高斯压缩比。编码后的数据再经过级联压缩管线：多项式系数从全精度浮点转为半精度（16~位），误差低于~0.01\%，再用~zlib~压缩获得约~2.14$\times$~额外压缩；素描高斯的~3D~位置（占编码后素描存储的~90\%~以上）则使用~Google Draco~几何压缩库，利用簇内空间连贯性实现约~22.3$\times$~压缩，平均位置误差低于场景直径的~0.04\%~。

\textbf{色块高斯优化：剪枝重训练~+~向量量化。}
色块高斯覆盖平滑区域，低频属性变化对精度不敏感。S\&P++~首先将已编码的素描高斯解码并固定，然后对色块高斯进行随机均匀剪枝以稀疏化分布，再在固定素描高斯的条件下重新训练色块高斯，使其补偿素描编码过程中引入的误差。重训练后，S\&P++~对色块高斯的各项属性（不透明度、缩放、四元数旋转、基础颜色和球谐系数）分别应用向量量化——用~256~项码本和~1~字节索引替代精确值，仅位置保留~16~位半精度浮点。剪枝~+~重训练~+~向量量化使色块高斯的存储减少~77\%--92\%~。

\textbf{分层渐进流式传输。}
素描-色块的语义分离天然支持分层渐进传输协议：第一层~$L_0$~仅传输紧凑的素描高斯表示（在~Playroom~场景中仅占~1.48~MB，约~8.5\%~的压缩模型大小），即可建立场景的几何骨架——物体边界、边缘和结构轮廓清晰可见；后续层~$L_1$~至~$L_4$~逐步增量传输色块高斯（按不透明度和体积覆盖率排序，优先传输视觉显著性高的部分），渐进填充体积细节直至完整质量。这种分层范式类似可伸缩视频编码中的基本层~+~增强层架构，非常适合移动端和~Web~端的~3DGS~查看器、需要快速初始加载的~VR/AR~应用、以及带宽受限的流式传输场景。

\textbf{实验表现。}在整体压缩效果方面，S\&P++~在三个数据集上实现了~116--175$\times$~的整体压缩比（Deep Blending: 703.77~MB~$\to$~4.03~MB，仅~0.57\%~的原始大小；Mip-NeRF 360: 767~MB~$\to$~5.90~MB），同时保持可比甚至超越原始~3DGS~的渲染质量（在~Mip-NeRF 360~上~PSNR~从~29.62~提升至~29.74）。相比前代~S\&P，在相同~PSNR~下~S\&P++~仅需~25\%--38\%~的存储。相比均匀剪枝基线，在相近模型大小时~S\&P++~的~PSNR~高出~1.74~dB、SSIM~高~6.7\%~、LPIPS~低~41.4\%~。值得注意的是，S\&P++~的端到端压缩时间甚至快于原始~3DGS~训练（低率配置~1559~秒 vs. 1634~秒），约为~HAC++~的~1.9$\times$~加速，这得益于其一次性聚类~+~定向优化的设计避免了层次化方法的迭代精炼开销。

\begin{takeaway}
结构化与锚点表示通过空间组织降低高斯存储和渲染开销。S\&P++ 将高斯按语义角色分离为素描（边缘定义）和色块（体积填充），为每类定制压缩管线，实现 116--175$\times$ 整体压缩比且渲染质量不降反升。其分层渐进传输协议使移动端仅需 1.48 MB 即可显示场景骨架。锚点方法（Scaffold-GS）则通过层次结构实现跨高斯的属性共享。
\end{takeaway}


%=====================================================================
\section{其他方法与硬件加速}
\label{sec:11_6}

本节介绍不属于前述主要类别但同样重要的压缩与加速方法，涵盖扩散模型辅助的极端压缩、基于对称性的压缩、神经场辅助的混合表示以及硬件加速与边缘部署四个方向。

\subsection{极端压缩与混合表示}
\label{ss:11_6a}

本节讨论三条技术路线：（1）扩散模型辅助的极端压缩；（2）基于对称性的压缩；（3）神经场辅助的混合表示。

\subsubsection{扩散模型辅助的极端压缩}

本路线介绍利用扩散模型的强大图像恢复能力，在极端压缩率下弥补传统方法质量损失的方法。

当压缩率超过~100$\times$时，传统方法的质量下降变得不可接受——纹理模糊、边缘锯齿、颜色偏移等问题集中爆发。扩散模型（diffusion model）因其强大的图像恢复能力，成为极端压缩场景下的关键补充。

\paragraph{扩散恢复的极端压缩}

传统~3DGS~压缩方法（剪枝、量化、熵编码）在~100$\times$~以下的压缩率已接近极限。当目标压缩率超过~1000$\times$~时，3D~表示的极端退化使渲染伪影不再局限于简单的模糊或块效应——纹理丢失、几何畸变和辐射度退化同时出现，且伪影模式因场景而异。这类复合伪影远比~2D~图像退化（如下采样或~JPEG~压缩）复杂，因为它们源于扭曲的~3D~表示经光栅化聚合后在~2D~平面上的投影。在此极端场景下，无论多么精细的压缩算法都无法避免严重的信息损失，因此单纯优化压缩侧已不可行。

NiFi（2026.02）~\cite{2602_04549}~提出了一个关键洞察：与其在压缩端对抗不可逆的信息损失，不如将问题转化为一个\textbf{盲图像恢复}（blind image restoration）问题——先接受极端压缩带来的严重退化，再用预训练扩散模型的强大自然图像先验来修复。这就是其"先压缩，后修复"策略的核心思想。扩散模型之所以适合这一任务，是因为~LDM~（Latent Diffusion Model）在大规模自然图像上训练后，已经内化了对纹理、结构和语义的深层理解，能够从高度退化的输入中"补全"丢失的感知细节。

NiFi 的流程分为三个阶段。

\textbf{（1）伪影合成}（Artifact Synthesis）：先用~HAC++（SOTA~压缩方法）对预训练~3DGS~进行三级渐进剪枝，基元数量按~$|G_l| = c_{\min} \cdot \exp\!\bigl((|G_{L-1}| - c_{\min}) \cdot \tfrac{L-1-l}{L-1}\bigr)$~递减，其中~$c_{\min}=4096$~为最小基元数，剪枝指标为渲染损失对各属性梯度的绝对值。剪枝后的~3DGS~再经微调、8~位量化和熵编码，得到极端低码率的压缩模型。从这些压缩模型渲染出退化视图，与原始视图配对，构成~1000~个场景的合成伪影数据集。

\textbf{（2）修复分布匹配}（Restoration Distribution Matching）：NiFi 在~Stable Diffusion~3（SD3）的~Transformer~骨干上训练两个低秩适配器（rank-64 LoRA）：修复适配器~$\theta^-$~负责从退化图像恢复干净图像，分布适配器~$\theta^+$~建模恢复后图像的分布以支持变分蒸馏。关键设计是将退化图像的潜表示~$z = E(\hat{x})$~沿扩散前向过程映射到中间时间步~$t_0 = 199/1000$，然后执行单步恢复：$\hat{z} = z_{t_0} - t_0 \, v_{\theta, \theta^-}(z_{t_0},\, t_0)$。这一重参数化使模型能在一步推理中利用扩散轨迹的随机多样性，而非直接从退化输入恢复。

\textbf{（3）感知匹配}（Perceptual Matching）：最终修复目标结合分布匹配~KL~散度、$\ell_2$~损失和~LPIPS~感知损失：$\mathcal{L}_{\theta^-} = \mathcal{L}_{\mathrm{KL}} + \lambda_2 \,\ell_2(x, \hat{x}) + \lambda_p \,\mathrm{LPIPS}(x, \hat{x})$，训练时~$\theta^-$~和~$\theta^+$~交替优化。

在实验中，NiFi 在~Mip-NeRF 360、Tanks and Temples~和~Deep Blending~三个基准上评测。以~Deep Blending~为例，原始~3DGS~模型为~555~MB，经~HAC++~极端压缩后降至~0.599~MB（$927\times$~压缩），NiFi~恢复后的感知质量（LPIPS、DISTS）与未压缩~3DGS~相当，显著优于~BM3D、SwinIR、DiffBIR~等图像恢复基线和~Difix3D~等~3DGS~恢复方法。消融实验证实中间时间步~$t_0$~对性能至关重要——缺少此重参数化时，单步恢复的质量显著下降。推理时单帧恢复仅需~100~ms（NVIDIA A40~GPU，无额外优化），具备实际部署潜力。

ExGS（2025.09）~\cite{2509_24758}~与~NiFi~共享"先极端压缩，后扩散恢复"的总体思路，但将场景定位从"盲图像恢复"转向"面向传输的3D资产压缩管线"，并引入结构化剪枝策略以更好地配合扩散模型的补全能力。

\paragraph{与 NiFi 的关键差异}

NiFi~用~HAC++~的固定剪枝方案压缩，ExGS~则设计了互补的双轴剪枝策略：（1）\textbf{细节保留剪枝}（Detail-Preserving Pruning, DPP）——沿用~LightGS~的全局显著性分数~$GS_j = \sum_{i=1}^{MHW} \mathbf{1}(G(X_j), r_i) \cdot \alpha_j \cdot T_j$~（$T_j = \prod_{i=1}^{j-1}(1-\alpha_i)$~为累积透射率），按分数排序后剔除低分基元以保留视觉显著纹理；（2）\textbf{几何保留剪枝}（Geometry-Preserving Pruning, GPP）——DPP~虽保留纹理细节，但在稀疏区域可能造成覆盖空洞。GPP~将空间划分为体素~$v$，确保每个体素内至少保留一个基元，对体素~$V$~中的基元子集最大化~$S_V = \max_{S_V |S|=k} \sum_{g_j \in S} s_j$~（$k=|V|^\beta$），平衡全局重要性与空间覆盖均匀性。此外，GPP~阶段还引入一个轻量~MLP~预测保留基元的缩放参数，使其自适应调整为对扩散补全最有利的尺度。

\paragraph{GaussPainter：掩码引导的单步扩散修复}

NiFi~在~SD3~的~Transformer~骨干上训练~LoRA~适配器，ExGS~则设计了更轻量的修复模块~GaussPainter，针对扩散模型的两大缺陷——幻觉和非实时性——做了专项优化。为抑制幻觉，GaussPainter~采用双重策略：（1）\textbf{潜空间监督}——用~VAE~编码真实图像~$y$~与退化渲染~$x$~为~$z_{hq}=E(y)$、$z_{lq}=E(x)$，最小化~$L_\text{latent}=\|z_{lq}-z_{hq}\|_2^2$，在潜空间中引导模型区分"应保留区域"与"需补全区域"；（2）\textbf{掩码引导}——将~3DGS~渲染中累积不透明度导出的可见性掩码经全局池化后投影到文本条件嵌入空间，逐元素加到~caption~特征上，自动降低可靠区域的权重并引导去噪器聚焦于缺失区域。推理时，退化渲染的潜表示在固定时间步~$t=199/1000$~执行单步去噪，配合~TAESD~轻量自编码器替换标准~VAE，实现接近实时的帧率（$512^2$~分辨率下单帧~0.23~s，NVIDIA A100）。

\paragraph{压缩性能}

ExGS~的完整压缩管线（双轴剪枝~+ 8~位量化~+ 熵编码）将原始~3DGS~模型从~354.77~MB~压缩至约~3.31~MB（$>100\times$~压缩），配合~GaussPainter~修复后恢复高质量渲染。压缩过程耗时~2.13~s，解压~0.23~s，总体效率显著优于需要多步迭代的扩散基线。

\subsubsection{基于对称性的压缩}

本路线介绍利用场景中的镜像对称性消除冗余基元，从结构层面减少需存储基元数量的压缩方法。

SymGS（2025.11）~\cite{2511_13264}~是本小节的独立方法，其核心思想是利用场景中的镜像对称性消除冗余基元，实现一条与剪枝、量化、熵编码完全不同的压缩路线——不是减少每个基元的存储位数，而是从根本上减少需要存储的基元数量。

\paragraph{Baseline 问题：对称结构的重复建模}

人造场景（建筑外观、室内布局、家具）中大量存在反射对称结构——桌子、瓶子、建筑立面等物体的一半近似是另一半的镜像。然而，3DGS~在训练过程中完全没有对称性的概念：它会独立地、重复地建模对称的两侧，为一对互为镜像的区域分配两组各自优化、互不相关的高斯基元。这种冗余在大规模场景中尤其显著——一栋对称建筑的外墙可能消耗数万乃至数十万个高斯来分别建模左半侧和右半侧，而两者的几何与外观几乎完全相同。尽管~HAC~等~SOTA~压缩方法通过自适应量化和熵编码压缩了每个锚点的属性，但它们无法消除对称基元之间的结构性冗余——因为量化只能减少每个属性的精度，熵编码只能缩短高频编码长度，两者都无法识别"这两个区域其实是彼此的镜像"。

\paragraph{改了什么：可学习的镜像平面与对称映射}

SymGS~的解决方案直观而优雅：找到场景中最重要的对称平面，丢弃对称一侧的高斯，在渲染时通过对面的高斯镜像"重建"被丢弃的那一侧。具体来说，SymGS~引入可学习的镜像平面~$M$，将保留侧高斯~$G_{\text{left}}^{M}$~关于~$M$~做镜像反射，生成对称侧高斯~$G_{\text{right}}^{M}$，从而只需存储一半的基元。镜像平面~$M$~的参数（法向量方向~$\theta, \phi$~和到原点距离~$d$）在训练中通过渲染损失的梯度回传持续优化，自动修正初始估计的离散化误差。

\paragraph{实现原理：从对称检测到镜像感知优化}

SymGS~的完整流程分为三个阶段：对称检测、镜像感知优化和迭代压缩。

\textbf{阶段一：对称检测}（Symmetry Detection）。给定初始高斯集合~$G_{\text{init}}$，SymGS~通过基于投票的聚类策略估计最显著的镜像平面。直觉上，如果场景中存在一对互为镜像的区域，那么将这对区域的任何两个视觉相似的高斯连线、取中垂面，就应该得到同一个镜像平面。SymGS~将这一直觉转化为算法：
\begin{enumerate}
\item \textbf{高斯聚类}：按颜色（HSV~空间离散化为~$N_c$~个分箱）、不透明度（$N_o$~个分箱）和缩放（$N_s$~个分箱）将高斯分入~$N_{\text{clstr}}$~个簇，确保只在视觉和几何相似的高斯之间配对，避免虚假对称检测。
\item \textbf{镜像平面参数化}：将~3D~平面~$ax + by + cz = 1$~重参数化为有界空间~$(\theta, \phi, d)$，其中~$\theta \in [0, \pi]$~和~$\phi \in [0, 2\pi]$~定义法向量的球坐标方向，$d \in [0, \text{extent}]$~定义到原点的距离（extent~为场景半径）。这一参数化确保搜索空间有限且可离散化。
\item \textbf{投票网格}：将~$(\theta, \phi, d)$~空间离散化为~3D~累加器体素网格~$A$，分辨率为~$d_\theta = \pi / \theta_{\text{res}}$、$d_\phi = 2\pi / \phi_{\text{res}}$、$d_d = \text{extent} / d_{\text{res}}$。每一对相似高斯~$(G_i, G_j)$~计算它们的中垂面参数并投一票到对应的体素；得票最高的体素即为最显著的镜像平面~$M_0$。
\end{enumerate}
投票算法的复杂度为~$O(n^2)$，SymGS~设计了~CUDA~加速版本，对约~100~万个高斯的场景，将最显著镜像平面的检测时间从数十分钟缩短至约~10~秒。

\textbf{阶段二：镜像感知优化}（Mirror-Aware Optimization）。确定~$M_0$~后，丢弃其法线方向负侧的高斯~$G_{\text{right}}^{M_0}$，仅保留~$G_{\text{left}}^{M_0}$~和不受~$M_0$~影响的高斯~$G_{\text{out}}^{M_0}$。在训练前向传播中，将~$G_{\text{left}}^{M_0}$~关于~$M_0$~做镜像反射生成对称侧（位置取镜像，缩放和不透明度直接复制），与~$G_{\text{out}}^{M_0}$~合并后送入~3DGS~可微分光栅化器渲染，计算与真实图像的~3DGS~渲染损失。反向传播时梯度同时更新~$G_{\text{left}}^{M_0}$~和~$G_{\text{out}}^{M_0}$~的属性以及~$M_0$~的参数~$(\theta, \phi, d)$，实现镜像平面与高斯属性的联合优化，自动修正初始离散化估计的不精确性。

\textbf{阶段三：迭代压缩}（Iterative Compression）。在优化完第一个镜像~$M_0$~后，SymGS~在剩余高斯~$\{G_{\text{left}}^{M_0}, G_{\text{out}}^{M_0}\}$~上重复相同的检测-优化流程，递归地发现更多镜像平面~$M_1, M_2, \ldots, M_L$，逐层压缩直至无法找到更多显著对称性。这种由粗到细的层次策略使~SymGS~能够捕获场景中不同尺度上的对称性——从建筑整体结构到家具局部细节。最终，每层只需存储保留侧高斯的均值位置~$X_{\text{left}}^{M_l}$~和镜像参数~$M_l$，以及最后一层的完整属性~$G_{\text{left}}^{M_L}$~和~$G_{\text{out}}^{M_L}$。

\paragraph{与~HAC~的集成}

SymGS~的框架可即插即用地叠加到~HAC~等锚点压缩方法上。具体集成流程为：先运行~HAC~估计锚点属性，然后对每个锚点聚合其关联高斯的平均颜色、不透明度和缩放用于聚类，后续的投票检测和镜像感知优化直接在锚点层面进行。每个镜像层的优化对象包括锚点特征、哈希网格和~HAC~的~MLP~模块。

\paragraph{实验结果}

在~Synthetic-NeRF（8~个场景）、Mip-NeRF 360（9~个场景）、Tanks and Temples（2~个场景）、Deep Blending（2~个场景）和~BungeeNeRF（6~个场景）共~27~个场景的评测中，SymGS~叠加~HAC~后平均相对~3DGS~的压缩率达到~108$\times$。相比~HAC~本身，SymGS~额外实现~1.66$\times$~压缩，其中~Deep Blending~场景收益最大（3.01$\times$~额外压缩），BungeeNeRF~次之（1.97$\times$），Mip-NeRF 360~和~Synthetic-NeRF~各约~1.4$\times$，Tanks and Temples~约~1.1$\times$——室内和建筑密集场景的对称结构最丰富，收益也最大。镜像参数的学习率设为~$10^{-3}$，距离分辨率~$d_{\text{res}}$~根据场景大小选取（Synthetic-NeRF~和~BungeeNeRF~取~0.01，其余取~0.1），角度分辨率~$\theta_{\text{res}} = \phi_{\text{res}} = 0.01$~固定。

\subsubsection{神经场辅助的混合表示}

本路线介绍将低频属性交由紧凑神经场预测、显式高斯仅保留高频细节的混合压缩方法。

HyRF（2025.09）~\cite{2509_17083}~（Hybrid Radiance Fields）是神经场辅助混合表示路线的代表性方法，其核心思想是"频率分工"：将低频变化的属性交给紧凑的网格神经场预测，仅让显式高斯保留难以被神经场捕捉的高频细节参数。

\paragraph{Baseline 问题：显式高斯的参数冗余}

原始~3DGS~中，每个高斯需要存储~59~个浮点参数，其中~48~个用于球谐系数表示视角相关颜色，~7~个用于各向异性缩放和旋转。这些参数中，大量低频变化（如缓慢变化的缩放和不透明度）完全可以通过空间连续的神经场从邻近位置推断，无需逐基元独立存储。然而，直接将所有高斯属性全部交给单一神经场预测又会导致高频细节丢失——网格表示的固定分辨率难以捕捉物体边缘处急剧变化的不透明度、薄结构的精细缩放，以及高频率的视角相关颜色变化。

\paragraph{HyRF 的混合策略}

HyRF 的解决方案是将高斯属性拆分为两个互补组件：
\begin{itemize}
\item \textbf{显式紧凑高斯}：每个高斯仅存储~8~个参数——位置~$\mathbf{p}_e \in \mathbb{R}^3$、漫反射颜色~$\mathbf{c}_e \in \mathbb{R}^3$、各向同性缩放~$s_e$~和不透明度~$\alpha_e$。这些是神经场难以准确建模的高频量。
\item \textbf{网格神经场}：基于多分辨率哈希编码（Instant-NGP），为剩余属性提供空间连续预测——各向异性缩放、旋转和视角相关颜色偏移。
\end{itemize}
相比~3DGS~的~59~维/基元，HyRF~将显式存储降至~8~维，整体模型尺寸减小约~20$\times$。

\paragraph{解耦神经场架构}

HyRF~发现用单一神经场同时预测所有属性效果不佳。几何属性（缩放、不透明度、旋转）与外观属性（视角相关颜色）之间的弱相关性使得联合学习困难。为此，HyRF~设计了双场架构：
\begin{itemize}
\item \textbf{几何场}~$\Phi_\text{geo}$：预测缩放、不透明度和旋转，输出为~$(s_n, \alpha_n, \mathbf{q}_n) = f_\text{geo}(\mathbf{f}_\text{geo})$，其中~$\mathbf{f}_\text{geo} = E(\mathbf{x}_i; \Theta_\text{geo})$~是多分辨率哈希编码的特征。
\item \textbf{辐射场}~$\Phi_\text{rad}$：预测视角相关颜色，引入相机方向的位置编码~$\mathbf{d}_c = \text{PE}(\mathbf{x}_i - \mathbf{o}_\text{cam})$，解码为~$\mathbf{c}_n = g_\text{rad}([\mathbf{f}_\text{rad}; \mathbf{d}_c])$，其中~$[\cdot;\cdot]$~表示特征拼接。
\end{itemize}
两场共享输入坐标但完全独立编码，消除了几何-外观的互相干扰。推理时，神经场预测的属性与显式高斯属性通过加法聚合：

\begin{equation}
\alpha = \sigma(\alpha_n + \alpha_e), \quad \mathbf{c} = \sigma(\mathbf{c}_n + \mathbf{c}_e), \quad \mathbf{q} = \text{Normalize}(\mathbf{q}_n), \quad s = \exp(s_n + s_e),
\end{equation}

其中~$\sigma$~为~sigmoid~激活，$\text{Normalize}$~为四元数~$L_2$~归一化。这种"神经预测+显式残差"的加法聚合方式让神经场负责低频基底，显式高斯提供高频修正。

\paragraph{混合渲染方案}

HyRF~的渲染管线包含三个阶段：（1）\textbf{可见性预剔除}——在查询神经场之前，先剔除视锥外的高斯点，实验显示这一步带来~3.9$\times$~渲染加速；（2）\textbf{前景高斯泼溅}——对可见高斯查询神经场、聚合属性后执行标准~3DGS~光栅化；（3）\textbf{神经场背景渲染}——3DGS~在远距离物体的致密化和优化上存在固有困难，常导致模糊背景。HyRF~在半径~$r=100$~的背景球体上，用辐射场预测背景颜色~$C_\text{bg} = T \cdot \mathbf{c}_s$（$T$~为前景累积透射率），再与前景混合~$C = C_\text{fg} + C_\text{bg}$。仅对透射率~$T$~低于阈值~$\tau_T=0.2$~的像素计算背景，兼顾质量和效率。

NeuralGS（2025.03）~\cite{2503_23337}~采取了比~HyRF~更激进的方向：HyRF~用神经场预测高斯属性后通过加法聚合与显式属性融合（"神经预测+显式残差"），NeuralGS~则完全用~MLP~替代高斯的属性存储，将压缩思路从"减少基元数量"彻底转向"用神经网络的参数共享替代逐基元存储"。

\paragraph{与 HyRF 的关键差异}

HyRF~保留了显式高斯属性~$(\alpha_e, \mathbf{c}_e, \mathbf{q}, s)$~作为高频修正，NeuralGS~的设计动机在于：既然~NeRF~可以用几~MB~的~MLP~参数表示复杂场景，那么高斯属性中一定存在大量可通过参数共享消除的冗余。NeuralGS~的核心压缩机制包含两个互补模块：（1）\textbf{空间条件预测}——基于~HAC~中多分辨率哈希网格~$H$~捕获的空间条件~$f_c = H(x)$，通过轻量预测网络~$P()$~预测锚点特征~$f_p = P(\text{concat}(f_c, f_r))$，其中~$f_r$~为可学习残差。由于网格存储了丰富的~3D~空间信息，预测特征~$f_p$~渲染出的图像与原始锚点特征~$f$~非常接近，因此只需编码和存储残差~$f_r$~而非完整特征~$f$~（锚点特征占模型总大小的~38.95\%），大幅节省码率。（2）\textbf{实例感知超先验}——残差~$f_r$~与空间结构的相关性较弱（这正是残差的特性），因此~HAC~中基于网格的空间熵模型对残差的概率建模不够准确。NeuralGS~设计了一个实例感知上下文编码器~$E()$~提取残差自身的紧凑超先验~$z = E(f_r)$，将~$z$~与空间条件~$f_c$~拼接后输入概率估计网络~$M()$~估计残差的概率分布参数~$(\mu, \sigma, q_{f_r})$~。超先验~$z$~通过可学习量化步长和非参数~CDF~编码进码流，解码时优先恢复以指导残差解码。

\paragraph{训练与编码}

NeuralGS~基于~HAC~实现，总体损失函数为~$\mathcal{L}_\text{total} = \mathcal{L}_\text{scaffold} + \lambda_e(\mathcal{L}_\text{entropy} + \mathcal{L}_\text{hash}) + \lambda_m \mathcal{L}_m$，其中~$\mathcal{L}_\text{scaffold}$~为~Scaffold-GS~渲染损失，$\mathcal{L}_\text{entropy}$~包含锚点属性~$(l, o, f_r)$~和超先验~$z$~的估计码率。实验显示，NeuralGS~实现了~91$\times$~的平均模型尺寸缩减而不损害视觉质量，证明了"预测+超先验"组合在锚点特征压缩上的高效性。

\subsection{硬件加速与边缘部署}
\label{ss:11_6b}

本节介绍面向手机、VR~头显等边缘设备的~3DGS~部署方案，重点关注系统级联合设计策略。

\subsubsection{硬件加速与边缘部署}

\paragraph{面向边缘设备的联合设计}

边缘设备（手机、VR~头显、无人机）对~3DGS~的部署提出了独特挑战：有限的~GPU~内存、受限的计算能力和不稳定的网络连接。第~\ref{sec:02_7}~节已详细讨论了~ASIC~加速器和~GPU~管线增强等硬件级加速方案。本节聚焦于面向边缘部署的系统级设计。

\paragraph{Baseline 挑战}

在边缘设备上直接部署~3DGS，面临一个根本性矛盾：模型压缩是必要的（设备内存有限），但压缩会显著降低渲染质量。一种自然的思路是完全依赖远程边缘服务器渲染——将用户姿态上传到配备高端~GPU~的服务器，再将渲染好的图像传回用户。然而，这种"全远程"方案在网络条件差时会导致严重的传输延迟，破坏渲染的流畅性。例如在实验中，纯远程方案在低功率（$P=10$~mW）下端到端延迟超过~80~ms，远超~60~ms~的流畅性阈值。

\paragraph{ECO-GS：边缘协作高斯泼溅}

ECO-GS（2025.10）~\cite{2510_22718}~提出了一个优雅的折中方案——让每个用户同时维护一个本地压缩小模型和一个远程完整大模型，并根据当前渲染需求自适应地在两者之间切换。直觉上，大多数视角下本地小模型的渲染质量是可以接受的，只有在那些"痛点视角"（pain points）——即本地模型与真实场景差异显著的视角——才需要借助远程大模型。核心问题在于：如何识别这些痛点视角，并在通信资源有限的约束下，智能地决定哪些用户应该切换到远程渲染？

ECO-GS~的系统模型如下：一个配备~$N$~根天线的边缘服务器同时服务~$K$~个单天线用户。对每个用户~$k$，其渲染结果为
\begin{equation}
\hat{v}_k = (1 - x_k) \cdot \phi_k(s_k) + x_k \cdot \phi_\text{edge}(s_k),
\end{equation}

其中~$x_k \in \{0,1\}$~为协作状态变量（$x_k=1$~表示使用远程渲染），$\phi_k$~和~$\phi_\text{edge}$~分别是本地小模型和远程大模型，$s_k$~是用户的~6D~观察姿态。系统需要联合优化二元的协作状态向量~$\mathbf{x}=[x_1, \ldots, x_K]^\top$~和连续的功率分配向量~$\mathbf{p}=[p_1, \ldots, p_K]^\top$。

\paragraph{IRAC：集成渲染与通信优化}

上述设计的关键在于，渲染质量的优化不能脱离通信约束独立进行——即使某个用户在痛点视角上，如果信道条件差到无法在时限内完成传输，强制切换反而降低整体体验。ECO-GS~因此提出了集成渲染与通信（Integrated Rendering And Communication, IRAC）方法，将渲染误差最小化与通信资源约束联合建模为优化问题：

\begin{align}
\min_{\mathbf{x}, \mathbf{p}} \quad & \sum_{k=1}^{K} \mathcal{L}(\hat{v}_k, v_k) \notag \\
\text{s.t.} \quad & B_k \log_2\!\left(1 + \frac{\|h_k\|_2^2 \, p_k}{\sigma_k^2}\right) \geq \frac{x_k V_k}{T - T_0}, \quad \forall k \notag \\
& \sum_{k=1}^{K} p_k \leq P, \quad \sum_{k=1}^{K} x_k \leq S,
\end{align}

其中~$\mathcal{L}$~为基于~L1~距离和结构差异度（DSSIM）的渲染误差函数，$B_k$~为带宽，$h_k$~为下行信道向量，$V_k$~为图像数据量（比特），$T=60$~ms~为端到端时延预算，$T_0=6.5$~ms~为边缘渲染时间，$P$~为总发射功率预算，$S$~为服务器最大并发用户数。

第一个约束确保远程渲染的图像能在时限内传输完毕（传输速率必须满足图像数据量与可用时间的比值）；第二个约束限制总功率不超过基站预算；第三个约束保证服务器不会过载。

\paragraph{求解方法与实验验证}

IRAC~问题是一个混合整数规划（二元变量~$\mathbf{x}$~与连续变量~$\mathbf{p}$~耦合），且目标函数中的真实图像不可直接获取。ECO-GS~提出了两种求解方法：惩罚增广最小化（Penalty Majorization Minimization, PMM）和模仿学习优化（Imitation Learning Optimization, ILO）。PMM~通过三角不等式构造目标函数的上界替代，迭代求解；ILO~则用深度神经网络学习~PMM~的决策模式，将计算时间降低~$100\times$（从~PMM~的~$>100$~ms~降至~ILO~的~1~ms~以内），同时~PSNR~仅下降~0.02~dB。

在实验中（$N=600$~天线、$K=20$~用户、$S=10$~最大并发），IRAC~相比纯本地渲染方案（UserGS）至少降低~$15\%$~的渲染误差、提升超过~1.3~dB~的~PSNR；相比"全远程"方案则严格保证~60~ms~以内的端到端延迟。本地小模型通过将每个高斯维度从~64~压缩到~14，实现~50$\times$~以上的模型缩减，但~PSNR~和~SSIM~均有明显下降，验证了边缘协作的必要性。

Transcoding（2026.05）~\cite{2605_21051}~从一个不同的角度切入边缘部署问题：ECO-GS~解决的是"已有~3DGS~模型如何在边缘与云端之间分配渲染"，Transcoding~解决的是"已有3D资产（光场点云或网格）如何高效转换为~3DGS~模型"，跳过了通常需要原始多视角图像的依赖。

\paragraph{与 ECO-GS 的关键差异}

ECO-GS~假设~3DGS~模型已通过标准管线（SfM~+ 3DGS~训练）获得，但许多实际场景中只有光场点云或网格（如~MPEG~点云压缩社区产出的~8iVSLF~数据集），原始拍摄图像并不可用。Transcoding~的核心挑战在于：没有真实多视角图像时，如何生成训练~3DGS~所需的虚拟视图和相机参数，以及如何利用已有的~3D~几何信息初始化高斯基元。方案分三步：（1）\textbf{虚拟视图生成}——用~MPEG PccAppRenderer~对输入点云/网格在~$N_c$~个均匀分布的虚拟相机位置渲染多视角图像，颜色计算采用加权插值~$\text{Colour}_j = \frac{\sum_i w_i \cdot \text{Colour}_i}{\sum_i w_i}$~，其中权重~$w_i = \max(0, (\mathbf{v}_i - \mathbf{v}_j) \cdot (\mathbf{p} - \mathbf{v}_j))^n$~（$n=10$~为锐度参数）衡量相机~$i$~与观察方向~$\mathbf{p}$~相对点~$j$~的对齐程度；（2）\textbf{自定义初始化}——不使用~SfM~从虚拟图像估计稀疏点云，而是直接将输入模型的~3D~顶点（或网格面片中心）作为高斯初始位置，将位置学习率设为~0~以保证生成的~3DGS~不偏离原始几何表面，同时保留其他属性的学习率~$>0$~。对于密集且干净的输入（如~8iVSLF），还可以关闭自适应密度控制（克隆/分裂），因为初始点密度已经足够，这同时减少了训练时间；（3）\textbf{坐标系对齐}——将渲染器的视口坐标转换为~COLMAP~相机坐标系，确保渲染管线和~3DGS~学习管线使用统一的坐标变换。实验表明，自定义初始化产生的~3DGS~模型高斯数量远少于原始密集点云中的点数，且比默认~SfM~初始化收敛更快、表面表示更干净。

FlexGS（2025.06）~\cite{2506_04174}~与~ECO-GS~都关注~3DGS~在不同能力设备上的部署，但采取了截然不同的技术路线：ECO-GS~为每个设备预训练一个固定压缩比的本地小模型，部署时无法动态调整；FlexGS~则提出"训练一次，处处部署"的弹性推理范式——训练一个模型，推理时按任意比例动态选择高斯子集，无需额外微调即可适配不同设备的算力约束。

\paragraph{与 ECO-GS 的关键差异}

ECO-GS~的本地小模型对应固定的压缩比（如~50$\times$），更换设备需要重新训练。FlexGS~的设计动机正是消除这一限制：给定目标弹性比例~$e$~（保留高斯的百分比），FlexGS~通过两个可学习模块自动完成高斯选择与属性补偿。第一个模块是\textbf{自适应高斯选择器}——受启发于~Gumbel-Softmax~松弛技术，将离散的二值选择变量~$\beta_i \in \{0,1\}$~建模为可微分类分布~$P(\beta_i = m | A_i, e)$~，其中~$A_i = (\alpha, s, q)$~为第~$i$~个高斯的属性。选择器网络~GsNet~将比例~$e$~和属性~$A_i$~分别嵌入为特征~$h_e = \text{ReLU}(e \cdot w_e^1 \cdot W_e^2)$~和~$h_i^a = \text{ReLU}(a^T \cdot W_a^1 \cdot W_a^2)$~（$a = [1 : \alpha \in A_i]$~为拼接嵌入），拼接后通过混合层~$z = [h_e, h_i^a] W$~输出~2~维~logits，再经~Gumbel-Softmax~$P(\beta_i = m) = \frac{\exp((z_m + g_m)/\tau)}{\sum_{m'}\exp((z_{m'} + g_{m'})/\tau)}$~（$g_m \sim \text{Gumbel}(0,1)$~为随机噪声，$\tau$~为温度参数）生成可微的采样掩码。训练时~$\tau$~指数衰减，分布逐渐趋近~one-hot~确定性选择。为稳定训练，FlexGS~引入全局重要性~(GI)~分数~$\text{GI}_i = \sum_{p=1}^{QHW} \mathbf{1}(G(X_i), r_p) \cdot V(s_i) \cdot \alpha_i \cdot T(i,p)$~作为先验引导（$V(s_i) = \min(|s_1 s_2 s_3|/3 / V_{90\%}, 1)$~为体积归一化），通过正则化~$L_\text{GI} = |\boldsymbol{\beta} - M_\text{GI}|$~和比例约束~$L_\text{spar} = |e - \frac{1}{N}\sum_m M|$~确保选择结果既符合~GI~先验又满足目标比例。

第二个模块是\textbf{高斯变换场}——仅选择高斯子集必然在渲染中产生空洞和细节丢失（单纯微调~SH~和不透明度不足以补偿），因此~FlexGS~引入一个空间-比例神经场~$\Psi = \{\psi_l: \mathbb{R}^4\}_{l=1}^2$~（4D~坐标~$(X, e)$~映射到多通道特征），用~6~个~2D~平面~$\psi_l(k_1, k_2)$~分解（类~HexPlane~设计）并通过双线性插值查询特征~$f = F_m(f_m)$~。多头预测器~$\phi_X, \phi_s, \phi_q$~将特征转换为位置、缩放和旋转的位移向量~$T_X = \phi_X(f)$、$T_s = \phi_s(f)$、$T_q = \phi_q(f)$~，被选高斯经变换后移至新位置~$X' = X + T_X$~以填补缺失区域。

\paragraph{训练策略}

每次训练迭代随机采样比例~$e \sim E$~，同时渲染选中高斯的图像~$I_s^e$~和全部高斯的图像~$I_f$~，对~GT~双重监督：$\mathcal{L}_\text{render} = |I_s^e - I_\text{GT}| + |I_f - I_\text{GT}|$~，最终损失~$\mathcal{L} = \mathcal{L}_\text{render} + \lambda_1 L_\text{GI} + \lambda_2 L_\text{spar}$~。这种随机比例训练确保单一模型即可在任意比例下渲染，无需为每个设备单独训练。

\begin{takeaway}
其他方法与硬件加速涵盖了弹性压缩和跨设备部署。FlexGS 通过 Gumbel-Softmax 选择器和空间-比例神经场实现单一模型适配任意设备，无需为每个目标压缩比重新训练。这一方向标志着3DGS压缩从"固定目标比"走向"弹性适配"的范式转变，对实际部署具有重要意义。
\end{takeaway}


%=====================================================================

\section{近期研究（2026 年 8--11 月）}
\paragraph{DeGS：把透明度检查和重排放在硬件混合之前}
2026.08~\cite{2608_02099}。DeGS 不让每个像素处理单元在执行混合时同步检查有效片段。它先通过跨度解析找出有效任务，再重组成密集、无冲突的执行包，使后端主要用于处理有效混合，减轻稀疏覆盖和提前终止导致的阵列空转。

作者用 Tanks and Temples、Deep Blending 与 Mip-NeRF 360 场景测试从 720p 到 8K 的输出，并比较硬件架构的利用率、吞吐和能耗。相关收益对应论文中阵列、带宽和工艺配置，不等同于现有通用 GPU 上的软件端到端加速。

\paragraph{ESVR：以椭球原语直接拟合体强度并逐原语采样}
2026.08~\cite{2608_05564}。ESVR 用可微椭球拟合三维标量强度场，依结构学习与裁剪去除冗余，再沿与各原语相交的光线采样并应用传递函数。它处理的是体数据的表示和交互渲染，可与普通 3DGS 针对多视图图像的场景重建区分。

原文以大型稀疏体数据比较表示质量、渲染能力和压缩。与 TTHRESH、ZFP 和 SZ3 的对照评估数据缩减能力，作者明确将方法定位为体表示与渲染框架；该实验也不包括常规压缩的全部其他渲染版本。

\paragraph{G²ARD-GS：冻结关键几何后蒸馏紧凑高斯场}
2026.08~\cite{2608_05704}。紧凑重建除了保留当前视图质量，还可能需要给未来轨迹的外观适配和定位提供几何框架。G²ARD-GS 逐轮降低教师模型的原语预算，保护高频表面和固定锚点，再选择对受限预算有补充信息的监督视图，在不增删原语的拓扑上恢复外观。

原文用 MatrixCity block A 的 510 个留出视图评价紧凑重建，另用 495 个离轨视图评价冻结几何后的适配，用 KingsCollege 检验相机注册。五倍至三十倍压缩的结果来自这些协议。原文指出城市主评估只有一个大型场景，方法间优化预算未对齐，且只验证一种下游几何任务。

\paragraph{用显著性引导超像素合并前馈高斯}
2026.08~\cite{2608_10712}。该方法先以显著性引导的超像素切分组织逐像素原语，编码为特征高斯，再在视角之间匹配和合并。细节层级解码把输出原语数量与场景内容的复杂度关联起来。

论文分别对 DepthSplat、AnySplat 和 DA3 比较压缩比例与 DL3DV-Bench、MipNeRF360、Tanks and Temples 的新视角画质。有无输入位姿随骨干而异，具体收益还取决于压缩档位和数据集。

\paragraph{GS²CI：从单次快照压缩测量估计三维高斯}
2026.08~\cite{2608_13502}。GS²CI 用三维视觉基础模型初始化几何和相机，再令高斯渲染符合快照压缩成像的测量过程。粗阶段得到的相机轨迹在细阶段保持不变，二维基础模型从合成视角提供伪图像来校正局部外观。

实验既有从 NeRF Synthetic、DTU 等构造的合成测量，也有 SCINeRF 的真实数据。合成实验默认用 8 的压缩比和 0.25 的掩码比例。单次测量不能唯一确定所有视角的真实外观，伪视图属于模型先验提供的额外约束，需与观测重投影分开理解。

\paragraph{RoofGS：按流水线带宽与计算瓶颈协同加速渲染}
2026.08~\cite{2608_15785}。RoofGS 用 Roofline 分析区分前端的数据搬运瓶颈与光栅阶段的计算瓶颈。方法融合前端算子并压缩属性和排序键，在光栅化中同时处理两像素，并利用透明度裁剪限定后的数值范围近似指数计算。

作者在 RTX 4090 的 4K 测试中汇总 Mip-NeRF 360、Tanks and Temples 与 Deep Blending，报告平均 10.1 倍加速、616 帧每秒和 0.028 dB 的 PSNR 降低。近似指数的误差讨论取决于裁剪范围，整体速度则对应特定 GPU、图像分辨率及基线配置。

\paragraph{QuARC-GS：以锚点和量化残差压缩动态多视点视频}
2026.08~\cite{2608_18285}。QuARC-GS 以多层锚点共享初始高斯，按时间更新熵编码的量化运动，并用视图空间阈值选择需新增的原语，尽量将增量留给新出现或解除遮挡的内容。

论文在多视图视频的比较中报告每帧存储最多降至对照的约十一分之一。作者明确指出每段视频仍需高质量且成本较高的初始重建，逐帧优化也尚达不到真正实时流式交付；逐帧存储量不能直接当作传输端到端延迟。

\paragraph{极化透明度先验：从训练早期抑制冗余高斯}
2026.08~\cite{2608_22344}。该方法结合平方误差、按模糊区域增密和极化透明度正则，推动有贡献的高斯获得较明确的不透明度，同时尽早抑制低贡献基元，降低先大量增密再统一裁剪的冗余。

论文在 Mip-NeRF 360、Deep Blending 和 Tanks and Temples 比较紧凑度、训练时间与视觉质量。作者指出模糊增密使用固定 3000 至 7000 迭代的调度，复杂纹理可能丢失高频细节，原语数较低时局部渲染缺口也可能更明显。

\paragraph{KISS-GS：将原语压实与编码增益分开分析}
2026.08~\cite{2608_26948}。KISS-GS 先用 POPSpa 对训练得到的高斯做独立压实，再用 SOG-XT 编码，其中包括排序码本和协方差平滑。若编码端拥有训练视图和 GPU，还能选做考虑编码误差的微调。

作者以多个高斯预算画出率失真曲线。在 Mip-NeRF 360 的 INRIA 质量参照点，原文给出 15.7 倍、6.6 倍和额外 2.2 倍的分阶段尺寸缩减，其中最后一项需要微调。低预算下部分远背景细节消失，码本聚类和可选微调也带来编码时间。

\paragraph{非均匀量化：先估计原语重要性再合并落入同格的高斯}
2026.08~\cite{2608_28272}。该方法用原语重要性引导递归加权 Lloyd-Max 量化高斯中心，当多个原语量化到同一位置时执行重要性加权的几何合并，再编码剩余属性。量化和合并可接到点云编码器上，降低离散化和重复原语造成的编码浪费。

论文比较 G-PCC 和 V-PCC 等设置，报告测试中优于均匀体素化，并在若干配置得到更好的率失真。作者同时指出，MPEG Objects 的部分 G-PCC 配置中 Adaptive Voxelization 的率失真更好，故不能概括为所有设置都领先。

\paragraph{CC-4DGS：分别压缩动态变形和规范高斯属性}
2026.09~\cite{2609_02184}。CC-4DGS 用计算式变形场替换多分辨率哈希表，并压缩规范坐标中的高斯属性。该组合同时针对动态网络存储与点云存储，便于在保留动态渲染的前提下减少模型容量。

作者在 N3DV 和 Technicolor 与多种动态方法比较，并对 Swift4D 和本方法施加相同点云压缩以对齐存储比较。Swift4D 仍具有较高渲染速率，因此存储收益和实时速度属于需要分别衡量的量。

\paragraph{AnyGS2Mesh：从既有高斯和原分辨率视图直接生成网格}
2026.09~\cite{2609_03304}。AnyGS2Mesh 把高斯原语变成与视图对应的几何特征，以流式分块编码聚合可变数量的视图；混合深度细化将预测的局部信息与高斯渲染的尺度结合。最后经深度融合输出显式网格，无需对每个输入重新优化一套网格模型。

作者在 DTU 与 Tanks and Temples 验证，并用不同高斯重建作为输入检查可插接性。前馈免去的是逐场景的网格优化，高斯重建已存在，表面质量仍受输入表征和可用视图的证据限制。

\paragraph{分频率阶段训练多个高斯场以降低活跃原语数量}
2026.09~\cite{2609_03334}。该方法先将图像拆成低频底层和拉普拉斯残差，再从粗到细训练各频带的高斯场。训练时把不再活跃的层归档，以减少同时参与优化的原语；推断时分别渲染各层并合成结果。

论文在三类数据集的 1K 设置比较训练时间、原语量和画质。训练活跃内存的减少不能直接解释为推断只需一个场，部署仍需统计各层合成的成本。

\paragraph{从高斯到芯片：用实际负载而非原语数估计效率}
2026.09~\cite{2609_06157}。该研究将高斯存储、投影覆盖、高斯与图块关联、像素候选和有效混合分开统计，再映射到投影、排序、混合和训练更新。它用同输入 GPU 测量与文献可比性检查，展示某一阶段的计算下降可能被传输等开销抵消。

论文强调压缩和剪枝可能改变输出质量，视图一致性和更新策略也会影响重建。即便相同场景和设置也需核对计时覆盖和实际质量；更少原语、更短内核时间或更快迭代都不能单独代表端到端成果。

\paragraph{LinearMask-GS：先稳定重要性排序再按目标比例裁剪}
2026.09~\cite{2609_10095}。LinearMask-GS 把短期学习的高斯掩码激活改为线性增量，延缓过早二值化；跨视图最大池化的重要性用于在指定预算按排序裁剪，再对剩余原语细化。它将掩码学习与最后的硬选择分成两个阶段。

论文在 Mip-NeRF 360、Tanks and Temples 和 Deep Blending 检查重建并验证多种高斯骨干的迁移。解决的是掩码过早饱和与排序不稳定，真实率失真和训练耗时仍需以各场景的保留率、细化预算及相应渲染指标理解。

\paragraph{2D GauSS-MI：按面元方向估计主动采集的信息量}
2026.09~\cite{2609_21516}。该方法对增量二维高斯场中的每个面元估计视角相关的重建可靠性，以显式法向组织方向证据，再用互信息评价候选相机位置对几何与视觉重建的价值，驱动下一视点选择。

实验在 Habitat 中用带 RGB-D 的四旋翼遍历 Replica 八个房间与办公室场景。表 III 同时报告图像及几何指标，其中该方法 PSNR 为 36.87 dB，Chamfer 指标上 ActiveGS 较优。结果是仿真场景、既定传感器与导航方案下的比较。

\paragraph{CRP-GS：用锚点对应与共享先验联合压缩高斯}
2026.09~\cite{2609_23005}。CRP-GS 按特征一致性联系锚点，以对应关系建立条件编码层级，并从共享特征中聚合低冗余先验。率失真训练同时约束锚点表示与编码概率，而不是仅利用锚点距离识别可共用的信息。

作者在 Mip-NeRF 360、Tanks and Temples、Deep Blending 和 BungeeNeRF 比较存储与渲染。实现要等到第 30000 次迭代后锚点趋于稳定再开启层级训练，并继续训练 10000 次，因而结果包含特定训练计划而非任意阶段的后置压缩。

\paragraph{GARO：用梯度活动与曲率联合判断动态高斯冗余}
2026.09~\cite{2609_23509}。动态场景的增密可能积累对后续优化贡献较低的高斯，但低梯度本身不足以判断原语是否可删。GARO 将冗余判定加入自适应密度控制，先按优化活动选出候选，再用局部几何估计判断是否处在可由附近原语表达的低曲率区域。

原文在合成及真实动态场景评价速度、显存和渲染质量，包含与 D-3DGS 的 NeRF-DS 定性对照。结果支持所测数据上的速度与质量折中。几何平坦不保证外观简单，含细薄结构或高频纹理的场景仍应分别报告剪枝效果。

\paragraph{在不同动态高斯估计器的输出上独立进行压缩}
2026.09~\cite{2609_25633}。该工作把动态压缩作为对逐帧高斯序列的后处理，与高斯估计器本身分开。这样可把同一压缩器用于跟踪、半跟踪及逐帧独立的输入，分别测试不同时间对应质量下的率失真。

作者在 N3DV 六个三百帧序列比较多种输入估计器。论文的 BD-rate 以未压缩输入高斯渲染为失真参照，隔离压缩损失。因此这些 BD-rate 不能与直接对真实拍摄图像评分的端到端文献表混比。

\paragraph{观测 Gram 矩阵：用训练相机对球谐压缩定权}
2026.09~\cite{2609_28997}。该方法按训练视线方向与透明度混合贡献累积观测 Gram 矩阵，为球谐截断、码率分配和量化定权。权重反映不同系数组合在已观测方向上的颜色误差，而非简单均匀处理各阶系数。

实验用九个 Mip-NeRF 360、两个 Tanks and Temples 和两个 Deep Blending 场景的原始预训练高斯，训练相机累积统计，留出相机评价。作者指出一阶模型忽略颜色截断和共享像素中原语间的交叉项，远离捕捉视角时相对截断的优势可能消失。

\paragraph{双阶段 Morton 排序：为块式视频编码重排高斯属性}
2026.09~\cite{2609_29041}。该方法先把量化高斯位置按三维 Morton 索引排序，再用二维映射组织规则属性图。相近位置对应的属性也更易落在二维相邻块，便于传统块式编码器利用空间相关性。

论文在 GSCodec Studio 上比较 PLAS、普通 Morton 与提出的排序，采用 MPEG 测试条件中的静态单帧和动态多帧，列出率失真及计算结果。结论依赖选用的平台、码流与场景运动，不能仅据一种排列推断所有动态资产的压缩收益。

\paragraph{ADATEX4D：在冻结变形骨干上为纹理按方向分配容量}
2026.09~\cite{2609_29963}。ADATEX4D 不再给各动态高斯分配同样大小的纹理，而是依据可见归一梯度、时间峰值和变形尺度对纹理两轴分别扩容，同时冻结已有四维高斯骨干。

在 12 个场景汇总中，原文相比 Textured 4DGS 报告纹理尺寸由 109.3 MiB 降至 51.7 MiB、完整模型存储下降 24.4\%，但平均渲染由 100.8 降至 84.1 帧每秒、训练时间上升。作者指出当前稠密纹理实例化限制运行时收益，不能把省存储视为提速。

\paragraph{实用高斯压缩：在不依赖空间聚合的条件下整数化解码}
2026.09~\cite{2609_30245}。这项工作由坐标派生几何上下文，再对紧凑锚点潜变量与属性作因果编码。线性变换与激活构成的上下文网络配合率失真学习、量化感知训练和整数推理，简化解码并降低平台浮点差异。

论文在多个公开场景比较码流大小、渲染质量和总编码、解码时间，区分不同训练预算与量化感知配置。跨平台一致性是实验所用整数解码路径的结果，不能据此对未测硬件或任意码流作兼容性保证。

\paragraph{用采集规模估计容量，再按训练质量调整高斯数量}
2026.09~\cite{2609_31248}。固定一百万高斯的默认预算对不同大小的采集不一定合适。这项研究在 LiteGS 骨干上将容量上限和实际使用量拆开：初始容量按总采集像素及视图重叠决定，训练再根据拟合质量增减原语。前者反映观察了多少表面，后者反映这些表面需要多少局部细节。

原文在 13 个标准场景和 8 个较大场景使用同一配置作比较。在标准基准，bicycle 比 LiteGS 多用 18\% 高斯，stump 则少用 14\%，两者 PSNR 均有提升；truck 缩到 LiteGS 的 0.68 倍数量时 PSNR 差距约 0.1 dB，但 LPIPS 较差。紧凑度与质量的取舍依场景和指标而变，结果针对离线场景训练。

\paragraph{无线边缘联邦高斯：用重要度和时延选择端侧训练子集}
2026.09~\cite{2609_32177}。该方法在每台设备按原语重要度与开销选择可训练的高斯，以适应异构计算、存储与通信资源。不同设备剪枝后原语结构不再相同，服务器按上传的掩码恢复参数与梯度的共同索引，之后才能做结构一致的联邦聚合。

论文把真实的 Mill19 Building、Rubble 和 UrbanScene3D Polytech 图像不均匀分给模拟设备，再比较无线边缘设置下的方案。真实图像与模拟链路需要区分；该实验并未对运营网络部署、链路故障或联邦更新的隐私性作保证。

\paragraph{FeCoSplat：用初次高斯重建反馈指导第二阶段编码}
2026.09~\cite{2609_33330}。FeCoSplat 先编码多视图的几何相关特征，解码的第一阶段高斯状态用来指导发送端补发未充分表示的信息。两步共享变码率特征编码器和隐式高斯状态，接收端用轻量交互恢复可渲染原语。

论文在 RE10K 与 ACID 的率失真比较中分别考察码率和渲染质量。反馈描述的是方法内部对第一次重建的分析，不能直接解释为测过网络往返时延；发送端的第二阶段也仍需消耗计算。

\paragraph{SAGA：以粗细四维锚点压缩体积视频}
2026.09~\cite{2609_33969}。SAGA 让稀疏粗锚点提供细锚点的共享特征，再从细锚点解码高斯运动和外观。遗漏的局部细节通过锚点坐标中的哈希残差补偿。熵编码按层组织已解码的时空上下文，并用固定容量记忆减少重复信息。

论文在 Neu3D、Panoptic Sports 和 MPEG MIV 的条件下比较码率失真、解码与新视图质量。实时解码并不意味着实时编码，方法仍要逐场景优化。快速运动、拓扑变化和稀少的相机覆盖是作者指出的边界。

\paragraph{全景高斯编码：让层级球面原语的保留由率失真决定}
2026.09~\cite{2609_34367}。方法将原语位置锚定在层级 HEALPix 网格，无需为每个存在的原语另传坐标。量化透明度所导出的存在掩码与熵模型共同决定哪些原语的失真下降值得其码率，可解码整个球面或视口。

论文主要在 100 张 2048×1024 的等距柱状全景测试，并用 100 张 SUN360 图像验证跨数据集表现。评价用球面加权 PSNR 和六个视口指标，避免平面指标过度加权极区。这些是静态全景编解码实验，不能直接代表动态头显传输时延。

\paragraph{NRF-GS：用共享神经残差提高外观表达效率}
2026.09~\cite{2609_37115}。标准 3DGS 给每个原语单独存储解析球谐系数。NRF-GS 保留原语的空间表示与基本颜色，为高斯配置潜码，用共享神经残差场预测视角相关外观。方向编码再按低频与高频分开，从而把反射细节集中交由共享模型表达。

论文将参数表达能力的等内存模拟与完整重建实验分开，在 Mip-NeRF 360、DL3DV 与 Tanks and Temples 报告高斯数、训练时间及图像指标。原文所称原语最多减少 50\% 是所测设置的上限；共有网络成本也应统计。补充表标明某个 Mip-NeRF 360 对照因基线显存不足排除了 Bicycle 和 Garden。

\paragraph{TSGL：由解码端可见属性重建高斯编码图}
2026.09~\cite{2609_38635}。TSGL 先由真实属性建立教师图，再用解码端已知的输入训练小学生网络重建近似邻接，以图傅里叶基压缩属性而无需传教师图。场景级的通道变换另用于移除高阶球谐系数的跨通道相关。

论文用 Mip-NeRF 360、Tanks and Temples 与 Deep Blending 的 13 个场景测试，分别报告数据集平均尺寸和渲染指标。作者在 Mip-NeRF 360 报告 29.32 MB，而 MesonGS 为 34.41 MB。结论依据所比较的后训练方法和测试渲染设置，图估计所需的学习并未被取消。

\paragraph{ByteSplat：减少分布式高斯训练的显存和跨卡流量}
2026.10~\cite{2610_02851}。ByteSplat 融合光栅化前后向处理以减少中间值访问，用硬件感知剪枝适配共享内存，并略去跨卡消息中的零梯度，在原语归属 GPU 汇总其更新。优化目标同时覆盖片外带宽和互连。

论文在各有 8 块 GPU 的 A6000 与 A100 服务器上，针对六个数据集测量含启动时间的端到端训练，报告最高 6.1 倍加速。最高加速属于特定工作负载，不能视为每个场景或不同互连的固定收益。

\paragraph{PocketSplat：先为世界坐标潜单元分配预算再解码}
2026.10~\cite{2610_03192}。PocketSplat 将经标定的多视图证据汇合到世界坐标的局部潜单元，按目标预算精确分配整数配额，仅解码预算内的高斯。这样可以直接输出预定数量的原语，而不是生成密集高斯后临时截取。

论文在 DL3DV 用四个条件视图和 50 个评价帧评估、在 Tanks and Temples 与 MegaDepth 检查域外迁移，另对 Mip-NeRF 360 做设备侧测试。方法要求相机内参与位姿，端侧结果对应文中使用的设备和可部署对照。

\paragraph{Mobile-4DGS：把动态高斯部署约束写进表示与调度}
2026.10~\cite{2610_05289}。Mobile-4DGS 将外观修正预先烘焙进低阶球谐参数，由多视角误差指导高斯增密与裁剪。动态部分采用逐原语的显式时间运算，静态部分是零运动的特例。移动播放时只在深度排序仍满足几何条件时复用静态绘制列表。

论文在静态及动态场景测图像质量、体积、训练和实机吞吐，目标是资源受限下的取舍。缓存排序只适用于通过相机运动界限检查的帧，轻量渲染结果也不能直接推成不限资源设定下所有指标领先。

\paragraph{SteadySplats：以空间复用和方差训练降低随机渲染噪声}
2026.10~\cite{2610_05576}。SteadySplats 将低样本随机高斯渲染的噪声分为估计方差与样本数问题。一方面复用邻域和历时样本，另一方面训练时约束同一视线上贡献颜色的方差，使少量抽样也能接近按序透明度合成的结果。

论文以 7 个 Mip-NeRF 360 场景的 21 条相机轨迹、3152 帧验证，并将按序混合渲染作为参照。作者指出方差训练会稀释光度损失、略降充分收敛时的部分峰值质量指标；样本复用也并非零开销。

\paragraph{GSCV：用跨帧相容的属性图接入标准视频编码器}
2026.10~\cite{2610_07795}。GSCV 以参考布局约束相邻帧的属性排列，使同一三维结构在二维高位深属性图中尽量能被帧间预测跟踪。属性图再由视频编码器压缩，以利用已有标准而保持较高的精度上限。

论文比较 MPEG 测试中的已追踪和半追踪高斯序列，以 GSCodec Studio 与 G-PCC 为对照，报告码率和重建质量。作者指出重排后的属性图不同于自然视频，通用编码器仍会限制压缩率；多个视频属性流也会改变系统整合开销。

\paragraph{混合分层光线追踪：按区域组织高斯和三角形求交}
2026.10~\cite{2610_13498}。该工作按空间局部的密度选择加速结构，复用相邻光线的遍历状态，并在混合求交路径复用高斯和三角形的部分计算。目标是减轻局部不均衡和对邻近光线重复访问。

作者用 RTX A5000 上的 RaySplats 作为参照，并以匹配视角校验 C++ 渲染器输出。端到端速度把实测 CUDA 与模型估算的修改版 RT Core 结合，因而是架构评估，不能当成现有芯片上完整实现的实测加速。

\paragraph{SplatQ：在固定预算下联合选择相互影响的高斯}
2026.10~\cite{2610_13503}。按每个原语独立得分剪枝可能同时移除互相补偿的一组高斯，也容易让图像局部失去过多容量。SplatQ 按空间与外观划分后分配区域预算，并用颜色响应的二阶近似构造固定预算的二次选择问题，将成对影响与单点贡献共同用于删减。

原文在五个基准的 27 个场景对照重建与渲染性能，并单独衡量区域间质量不均。在其中一组消融，区域预算分配和带二次项的贪心解分别改变宏平均 PSNR。对照结果依赖预训练场景、剪枝预算和微调配置；离线求解开销也应与剪枝后的渲染帧率分别比较。

\paragraph{TETRIS-4DGS：以重要性和运动组织后训练量化}
2026.10~\cite{2610_14663}。TETRIS-4DGS 从已有动态高斯估计逐原语重要性，据此自适应量化，并按重要性与运动相关性把属性排列为二维映射，用常规图像和视频编码器分别处理相应数据。

作者从 MeetRoom 三个场景生成重要性训练标签，在 Neural 3D Video 每场景的前 300 帧评价实际压缩率、渲染失真和编解码用时。这里的后训练压缩指输入的四维高斯不用逐场景重新训练，重要性估计器仍需从已重建的资料学习。

\paragraph{PLEDGE：在声明相机域内为渐进码流提供误差界}
2026.10~\cite{2610_17605}。PLEDGE 将可接受的屏幕位置、栅格足迹、颜色和透明度偏差转换为逐高斯量化约束，并组织分层位平面，使有效截断点上的量化上界可计算。

作者将误差界验证和率失真比较分开，同时检查渐进解码。论文明确把界定义相对已剪枝模型、限制在声明的相机域，并以光栅足迹代替真实透视轮廓；对合成颜色的界还要求几何和不透明度固定。

\paragraph{RDCompGS：在高斯属性拟合时考虑二维编码率失真}
2026.10~\cite{2610_19538}。RDCompGS 在高斯表示优化中利用变换和量化系数的分布估计码率梯度，并调节率约束，使结果适合后续二维编码器，同时保留渲染质量目标。这样把原语优化与其最终能否高效编码关联起来。

论文在 Tanks and Temples、Deep Blending 和 Mip-NeRF 360 上对照 H.265/HEVC 与 JPEG XL 压缩，并比较 PSNR、SSIM 与 LPIPS。收益对应这些实际编码设置和可微码率估计，不能据此保证其他格式同幅受益。

\paragraph{前馈多视图特征编码：在接收端才构造高斯场景}
2026.10~\cite{2610_20590}。该工作先对联合外观和几何的多视图特征编码，再在接收端从解码特征前馈地产生可渲染高斯。率失真联合训练让码流控制特征信息量，并与先重建再压缩的链路区分。

作者在 RE10K 训练并在 ACID 和 DL3DV 跨域测试，前两个评估各使用 8 个输入视图，DL3DV 使用 12 个，三者均有 8 个目标视图。每个质量权重单独训练；与先生成再压缩方案的降码率不能泛化到其他输入数量或其他基线。

\paragraph{ResStreamGS：按残差敏感度分流动态视频高斯}
2026.10~\cite{2610_21121}。ResStreamGS 将视频高斯分成在一组图像内可复用的稳定部分，以及对帧间变化敏感、要逐帧编码的部分。轻量的探测适配和多属性残差评分指导选取分流比例。

论文在 Neu3D 的五段和 MeetRoom 的一段全场景视频评价，每段 300 帧、每组 20 帧，包含少量局部动态与大面积稳定背景。码流与重建比较应按作者选定的全场景序列和对照条件解释，大比例全局运动会改变可复用分量的前提。

\paragraph{LENS：在嵌入式 FPGA 协同实现定位与固定点高斯渲染}
2026.10~\cite{2610_25068}。LENS 分别实现坐标回归和有序的定点透明度混合，前者针对姿态精度，后者通过表示、递推和数据搬运协同设计减少边缘设备上的冗余带宽。

实验以 KR260 和 Jetson Orin NX 为板级比较，选用 Cambridge、Wayspots 定位和 RealEstate10K 新视图设置，训练另列为边界实验。文中的性能功耗比指相应姿态回归与渲染任务，作者不把 FPGA 定位为通用训练加速器。

\paragraph{CCT-GS：用可接入标准编码器的瓶颈图压缩高阶球谐}
2026.10~\cite{2610_26920}。CCT-GS 在多个场景训练高阶球谐图的变换，以可微编码器代理把高维属性整理为低维瓶颈图。目标场景可进一步调节前处理器，解码端在冻结主体的情况下用 LoRA 参数恢复场景相关细节。

作者按 MPEG GSC 的前向场景单帧设置、用固定留出视图比较视频式基线。训练与场景适配仍是单独步骤；兼容标准视频编码指瓶颈图的编码格式，完整高斯重建仍需论文的学习式处理。

\paragraph{KoalaGS：让同一套高斯在预算递减时保持结构}
2026.10~\cite{2610_28615}。KoalaGS 在训练中用可预先计算的图像复杂度引导原语重定位和保留排序。推理选择固定排序的前缀，从同一个模型取得不同渲染预算，而无需为每个数量重新建模。

作者以固定预算质量、完整预算区间的指标曲线以及过渡稳定性三方面比较五个基线。该方法的排序是视角无关的，并不等同于相机依赖的可见性层次结构；随机重定位仍会在重复训练之间带来一定差异。

\paragraph{ClunoGS：在同一优化中压缩高斯的语义与外观}
2026.10~\cite{2610_30013}。ClunoGS 用一个约束目标组织语义感知原语稀疏化和向量量化，以原对偶过程交替更新场景参数、压缩变量及对偶量。对语义与 RGB 同时施加监督，缓解各阶段独立压缩造成的语义丢失。

论文在 LERF 与 ScanNet-10 评价，并在同一 RGB 基础场景及相同高斯保留比例下比较先剪再量化的组合，也测试先学习语义与先压缩两种顺序。结果表明相应压缩率下的定位和分割改善，不能推出该优化对所有语言表示或码率更优。

\paragraph{BoundedGS：分别限制光栅工作区和候选梯度存储}
2026.10~\cite{2610_34233}。BoundedGS 用有负载上限的光栅单元缩短中间状态寿命，对可能可见的候选基元再做详细投影，并让反向梯度按候选紧凑存储。其设计针对高分辨率训练峰值工作区，而不是把参数和优化状态本身全部替换。

作者在固定约 600 万高斯的分辨率扫描中报告 1K 至 16K 的渲染工作区保持约 1.073 至 1.077 GiB，在五个原生 5K 场景单独报告完整训练的工作区下降。工作区按附录用分配峰值扣除模型参数与优化器等项目计算，不是整机内存或显存总预算。

\paragraph{ParetoGS：以稀疏适配锚点连接多个渲染预算}
2026.10~\cite{2610_34235}。ParetoGS 以训练完成的完整模型为起点，交替收缩原语集合并在固定成员内调整参数，形成嵌套的预算锚点。持久身份让锚点可共享成员，残差图集编码附加状态，解码时用插值覆盖锚点之间的请求预算。

论文在 Mip-NeRF 360、Tanks and Temples 与 Deep Blending 共 13 个场景对照质量和预算，使用相同留出视角。适配状态是相对完整模型保存的额外信息，需与仅存一个全模型前缀的方法按共同基线和总存储口径比较。

\paragraph{RacerGS：利用因果组上下文压缩高斯锚点残差}
2026.10~\cite{2610_34263}。RacerGS 将哈希编码投影为紧凑空间上下文，以预测的属性均值和量化步长压缩残差。FiLM 调制的循环网络再利用已解码组帮助估计后续组的熵分布。

作者在多数据集报告相对原始 3DGS 与 Scaffold-GS 的平均体积压缩，并以相应渲染指标比较。论文指出自回归上下文必须按顺序解码属性组，新场景仍需逐场景优化，压缩比不能理解成无损或实时流式解码能力。

\paragraph{MatrixSplat：将高维特征合成映射到矩阵并行单元}
2026.11~\cite{2611_15153}。MatrixSplat 以顺序无关的加权累加改写特征渲染，并将打包原语送入线程束的矩阵乘加。分批按需加载通道以减轻共享内存压力，目标是使增加语义等输出特征时不被串行混合与重复读入限制。

论文把 Mip-NeRF 360 的渲染评价与 LERF、3D-OVS 的语义结果分别比较，还报告独立内核和端到端计时。独立内核实验的加速不等于训练或完整渲染吞吐；顺序无关近似也不能当作逐像素精确的有序透明度合成。

\paragraph{GDM-GS：把沿轨迹的高斯选择重写为空间数据查询}
2026.11~\cite{2611_20546}。GDM-GS 为结构化场景的锚点建立空间索引，在物化高斯之前查询并筛除无关锚点。相邻视点可共享已解码的数据，CPU 和 GPU 按调度共同处理锚点查询，再对各视点投影混合。

论文在 BungeeNeRF 的八个场景、1600×900 分辨率对照 Scaffold-GS。轨迹窗口和相邻视图相关性是收益来源，故空间查询、共享解码带来的加速不直接代表随机相机访问或未预知位姿的交互负载。

\paragraph{选择性冻结：缓存低梯度高斯的光栅贡献以减少训练开销}
2026.11~\cite{2611_22059}。该方法逐步冻结低梯度原语，复用其对颜色和透明度合成的贡献，并重组反向光栅以兼容缓存。专用流水阵列与冻结原语缓存面向在训练中重复访问的模式。

论文分别评估冻结后的重建质量与专用架构的速度和能耗，报告最高冻结 90\% 高斯。专用吞吐与能效数字来自架构和缓存评估，不能作为购买的通用 GPU 上软件提速或流片结果引用。

\paragraph{CueSplat：为动态高斯调度稀疏精确变形}
2026.11~\cite{2611_22095}。CueSplat 固定已训练的动态高斯骨干，用历史状态预测多数原语的新参数，仅让路由器挑出的原语执行精确变形。场景预算规定精确更新规模，间歇恢复全量变形校正累积偏差。

论文用 Technicolor、N3DV 和 NeRF-DS 共 21 段序列和四种骨干比较版间吞吐及质量，基线与加速版本共用预训练检查点和渲染器。预算在场景层面校准，跨场景平均 PSNR 表现不能保证逐帧质量，短 Technicolor 序列的刷新设置也与其他数据不同。

\paragraph{Cross3DGS：在 Vulkan 中按设备可行性选择统一训练调度}
2026.11~\cite{2611_22240}。Cross3DGS 保持投影、混合、梯度和优化器的逻辑一致，仅调整内存布局与并行方式。程序根据驱动报出的硬件限制、任务规模和已知驱动问题筛除不可行方案，再按固定偏好顺序选一种调度，不在目标设备上实时搜索最快实现。

论文报告六家厂商的八个配置能完成训练，并用数值容差测试单步更新。浮点累加次序不同不会保持逐位一致，单步测试也不包含阶段间的 MCMC 原语细化，已观察到的设备与驱动规则不能直接推广到所有同厂配置。


\section*{本章小结}
\addcontentsline{toc}{section}{本章小结}

\begin{futurebox}{本章总结与展望}{ch11_summary}
本章围绕~3DGS~模型压缩与加速这一实用化关键问题，系统梳理了~161~篇相关工作。

\textbf{高斯剪枝与稀疏化}（第~\ref{sec:11_1}~节）从减少基元数量的角度出发，发展出基于重要性分数、可学习掩码、几何一致性和最优传输四条技术路线。LightGaussian~确立了全局显著性+渐进剪枝的基本范式，PUP 3D-GS~用二阶灵敏度分析提供了原则性理论基础，MaskGaussian~的概率性掩码实现了动态保留决策，Gaussian Herding~的最优传输视角将简化建模为全局分布匹配。

\textbf{编码与压缩}（第~\ref{sec:11_2}~节）从压缩基元参数的角度出发，发展出量化与向量量化、熵编码与上下文建模、前馈压缩和后训练管线四类方法。Compact 3DGS~建立了"掩码+神经场+码本"的压缩范式，HAC++~和~ContextGS~将熵编码中的超先验和自回归上下文模型引入~3DGS，FCGS~开创了前馈压缩的通用编解码路线，MesonGS++~实现了尺寸感知的后训练压缩管线。

\textbf{训练加速与高效表示}（第~\ref{sec:11_3}~节）从源头控制高斯增长和训练效率。Opti3DGS~的频率调制和~Gaussians on a Diet~的内存有界训练分别解决了冗余增长和峰值内存问题。

\textbf{视频与流式渲染}（第~\ref{sec:11_4}~节）将压缩与流式传输结合。Matryoshka GS~的嵌套式连续细节层次和~PD-4DGS~的分层变形分解使场景可以在下载过程中逐步可渲染。

\textbf{结构化与锚点表示}（第~\ref{sec:11_5}~节）从根本上改变了~3DGS~的组织方式。Scaffold-GS~的锚点驱动层级表示不仅减少了冗余，还为后续的熵编码提供了结构化基础，成为压缩方法的通用~backbone。

\textbf{其他方法}（第~\ref{sec:11_6}~节）包括扩散模型辅助的极端压缩（NiFi）、对称性利用（SymGS）和神经场混合表示（HyRF）等新兴方向。

当前~3DGS~压缩领域呈现出清晰的演进趋势：（1）从"训练后压缩"到"训练中控制"——在训练阶段就约束模型规模；（2）从"固定码率"到"可伸缩码率"——单一模型支持多速率点；（3）从"静态压缩"到"流式渐进"——边下载边渲染；（4）从"逐场景优化"到"通用前馈"——一次训练通用编码器。这些趋势共同推动~3DGS~从实验室走向手机、VR~头显等边缘设备的实际部署。
\end{futurebox}
